% Les mécanismes de l'immunité et ses dysfonctionnements — SVTEEHB Tle TI — Cours signé OnBuch+
% Programme officiel MINESEC. Compiler avec : tectonic N06-mecanismes-immunite-dysfonctionnements.tex
\newcommand{\DOCMATIERE}{SVTEEHB}
\newcommand{\DOCNIVEAU}{Tle TI}
\newcommand{\DOCMODULE}{Module II : L'Éducation à la Santé}
\newcommand{\DOCLECON}{N06}
\newcommand{\DOCTITRE}{Les mécanismes de l'immunité et ses dysfonctionnements}
% OnBuch+ — préambule « cours signés » ENRICHI (compilé avec tectonic / XeLaTeX)
% Système visuel « Pop » d'OnBuch+ : fond crème (jamais blanc), bordures encre
% épaisses, ombres « dures » décalées, titres Archivo Black en capitales.
% Version MATHÉMATIQUES : graphes (pgfplots/tikz), boîtes pédagogiques
% (définition, propriété, méthode, attention, le savais-tu, exercice résolu,
% exercice, TP, bilan), page de garde et signature OnBuch+ ancrée en bas.
%
% Macros attendues dans main.tex AVANT \input{preamble} :
%   \DOCMATIERE  \DOCNIVEAU  \DOCMODULE  \DOCLECON (numéro)  \DOCTITRE
\documentclass[11pt]{article}

\usepackage{fontspec}
\usepackage[french]{babel}
\usepackage[a4paper,margin=2cm,top=3.4cm,bottom=2.8cm,headheight=1.4cm,footskip=1.3cm]{geometry}
\usepackage{amsmath,amssymb,amsfonts}
\usepackage{mathtools} % \xmapsto, \coloneqq... souvent utilisés par les modèles
\usepackage{cancel} % simplifications barrées (bilans de forces)
% \abs{z} (module, valeur absolue) et \norm{u} (norme), utilisés par les modèles
\providecommand{\abs}{}\let\abs\relax\DeclarePairedDelimiter{\abs}{\lvert}{\rvert}
\providecommand{\norm}{}\let\norm\relax\DeclarePairedDelimiter{\norm}{\lVert}{\rVert}
\usepackage[table]{xcolor} % [table] : fournit aussi \rowcolors (alternance de lignes)
\usepackage{pagecolor}
\usepackage{tikz}
\usetikzlibrary{arrows.meta,positioning,calc,shapes.geometric,shapes.misc,shapes.arrows,decorations.pathmorphing,decorations.pathreplacing,decorations.markings,patterns,shadows,mindmap,trees,fit,backgrounds,matrix,intersections,angles,quotes}
\usepackage{pgfplots}
\usepackage[version=4]{mhchem} % équations nucléaires \ce{^{235}_{92}U}
\usepackage[european,siunitx]{circuitikz} % schémas électriques
\pgfplotsset{compat=1.18}
\usepgfplotslibrary{fillbetween,groupplots} % aires entre courbes (calcul intégral)
\usepackage{enumitem}
\usepackage{fancyhdr}
% [most] charge toutes les bibliothèques tcolorbox (listings, minted, raster,
% documentation...) et gonfle inutilement la mémoire TeX sur un document long
% (~300 boîtes) : on ne charge que ce qui est réellement utilisé.
\usepackage[skins,breakable]{tcolorbox}
\usepackage{graphicx}
\usepackage{colortbl}
\usepackage{booktabs}
\usepackage{tabularx}
\usepackage{multirow}
\usepackage{diagbox} % case d'en-tête diagonale des tableaux à double entrée
% Colonnes extensibles centrée / gauche / droite (C, L, R), souvent utilisées par les modèles
\newcolumntype{C}{>{\centering\arraybackslash}X}
\newcolumntype{L}{>{\raggedright\arraybackslash}X}
\newcolumntype{R}{>{\raggedleft\arraybackslash}X}
\usepackage{array}
\usepackage{multicol}
\usepackage{siunitx}
% parse-numbers=false : accepte aussi \SI{2,5 \times 10^{-3}}{...}
\sisetup{locale=FR,output-decimal-marker={,},group-minimum-digits=4,parse-numbers=false}
\DeclareSIUnit\ohm{\ensuremath{\symup{\Omega}}} % U+2126 absent de la police
\DeclareSIUnit\division{div} % réglages d'oscilloscope (ms/div, V/div)
\DeclareSIUnit\tr{tr} % tours (vitesse de rotation, ex. \SI{1500}{\tr\per\minute})
\DeclareSIUnit\molar{M} % concentration molaire (ex. \SI{3}{\molar}, \SI{1}{\micro\molar})
\DeclareSIUnit\an{an} % année (le modèle confond parfois avec \year, un compteur TeX) — ex. \SI{5}{\kilo\meter\per\an}
\DeclareSIUnit\FCFA{FCFA} % franc CFA (ex. \SI{500}{\FCFA\per\kilo\gram})
\DeclareSIUnit\degreeBrix{\textdegree Brix} % teneur en sucre d'un sirop/jus (réfractométrie)
\DeclareSIUnit\month{mois} % le modèle utilise parfois \month, un compteur TeX, comme unité de durée
\DeclareSIUnit\microliter{\micro\liter} % le modèle écrit parfois \microliter en un seul token
\DeclareSIUnit\million{millions} % dénombrement (ex. \SI{15}{\million\per\mL} spermatozoïdes)
\usepackage{qrcode}
% \degree : commande fréquemment utilisée par les modèles pour un angle ou une
% température en dehors de \SI{}{} (non fournie par les paquets chargés).
\providecommand{\degree}{^{\circ}}
% \qed (fin de démonstration), utilisé par les modèles sans amsthm
\providecommand{\qed}{\hfill$\square$}
% \barre : double barre des valeurs interdites dans un tableau de variations (tkz-tab)
\providecommand{\barre}{\ensuremath{\|}}
% environnement proof (amsthm non chargé) utilisé par les modèles
\ifdefined\proof\else\newenvironment{proof}[1][Démonstration]{\par\noindent\textit{#1.}\ }{\qed\par}\fi
\usepackage{lastpage}
\usepackage{hyperref}
\hypersetup{hidelinks}

% ============================================================
% Palette « Pop » OnBuch+ — mêmes hex que la classe P (plus_ui.dart)
% ============================================================
\definecolor{poporange}{HTML}{F59321}
\definecolor{popdark}{HTML}{E07A0C}
\definecolor{popcream}{HTML}{FFF6E8}
\definecolor{popink}{HTML}{1C1714}
\definecolor{poppurple}{HTML}{7A5AE0}
\definecolor{popgreen}{HTML}{2E9E6A}
\definecolor{poppink}{HTML}{E0457B}
\definecolor{popblue}{HTML}{2F7DE1}
\definecolor{popgold}{HTML}{C98A12}
\definecolor{popmuted}{HTML}{8A7F76}
\definecolor{popline}{HTML}{EADFD2}
\definecolor{popgray}{HTML}{8A7F76}
\colorlet{popgrey}{popgray}
% Alias tolérants (le modèle nomme parfois les couleurs différemment)
\colorlet{popwhite}{white}
\colorlet{popblack}{popink}
\colorlet{popred}{poppink}
\colorlet{popyellow}{popgold}
% Teintes claires (fonds de tableaux/encadrés) — le modèle nomme parfois
% les couleurs avec un suffixe L (ex. popblueL) sans les avoir définies.
\colorlet{poporangeL}{poporange!20}
\colorlet{popdarkL}{popdark!20}
\colorlet{poppurpleL}{poppurple!20}
\colorlet{popgreenL}{popgreen!20}
\colorlet{poppinkL}{poppink!20}
\colorlet{popblueL}{popblue!20}
\colorlet{popgoldL}{popgold!20}
\colorlet{popredL}{popred!20}
\colorlet{popgrayL}{popgray!20}
% Teintes claires pour fonds de tableaux / zones
\colorlet{popblueL}{popblue!10!white}
\colorlet{poporangeL}{poporange!14!white}
\colorlet{popgreenL}{popgreen!12!white}
\colorlet{poppurpleL}{poppurple!12!white}
\colorlet{poppinkL}{poppink!12!white}
\colorlet{popgoldL}{popgold!14!white}

\pagecolor{popcream}

% ============================================================
% Typographie
% ============================================================
\defaultfontfeatures{Ligatures=TeX}
\setmainfont{PlusJakartaSans-Medium.ttf}[Path=../../../fonts/,BoldFont=PlusJakartaSans-Bold.ttf]
% Maths en sans-serif (Fira Math) avec chiffres et lettres droites de Plus
% Jakarta, pour que les formules se fondent dans le texte.
\usepackage[math-style=ISO,bold-style=ISO]{unicode-math}
\setmathfont{FiraMath-Regular.otf}[Path=../../../fonts/]
\setmathfont{PlusJakartaSans-Medium.ttf}[Path=../../../fonts/,range={up/num,up/{Latin,latin},\mathup}]
\usepackage{icomma} % virgule décimale sans espace en mode math
% Caractères mathématiques Unicode absents des polices de texte (Plus Jakarta,
% Archivo Black) : rendus par leur équivalent mathématique, en texte comme en
% formule — sinon ils s'affichent en carré vide (ex. « Arithmétique dans ℤ »).
\usepackage{newunicodechar}
\newunicodechar{ℝ}{\ensuremath{\mathbb{R}}}
\newunicodechar{ℤ}{\ensuremath{\mathbb{Z}}}
\newunicodechar{ℂ}{\ensuremath{\mathbb{C}}}
\newunicodechar{ℕ}{\ensuremath{\mathbb{N}}}
\newunicodechar{ℚ}{\ensuremath{\mathbb{Q}}}
\newunicodechar{𝔻}{\ensuremath{\mathbb{D}}}
\newunicodechar{α}{\ensuremath{\alpha}}
\newunicodechar{β}{\ensuremath{\beta}}
\newunicodechar{γ}{\ensuremath{\gamma}}
\newunicodechar{δ}{\ensuremath{\delta}}
\newunicodechar{ε}{\ensuremath{\varepsilon}}
\newunicodechar{θ}{\ensuremath{\theta}}
\newunicodechar{λ}{\ensuremath{\lambda}}
\newunicodechar{μ}{\ensuremath{\mu}}
\newunicodechar{σ}{\ensuremath{\sigma}}
\newunicodechar{τ}{\ensuremath{\tau}}
\newunicodechar{φ}{\ensuremath{\varphi}}
\newunicodechar{ψ}{\ensuremath{\psi}}
\newunicodechar{ω}{\ensuremath{\omega}}
\newunicodechar{Σ}{\ensuremath{\Sigma}}
% majuscules produites par \MakeUppercase dans les titres (α-aminés -> Α)
\newunicodechar{Α}{\ensuremath{\alpha}}
\newunicodechar{Β}{\ensuremath{\beta}}
\newunicodechar{Γ}{\ensuremath{\Gamma}}
\newunicodechar{Δ}{\ensuremath{\Delta}}
\newunicodechar{Ε}{\ensuremath{\varepsilon}}
\newunicodechar{Θ}{\ensuremath{\Theta}}
\newunicodechar{Λ}{\ensuremath{\Lambda}}
\newunicodechar{Μ}{\ensuremath{\mu}}
\newunicodechar{Τ}{\ensuremath{\tau}}
\newunicodechar{Φ}{\ensuremath{\Phi}}
\newunicodechar{Ψ}{\ensuremath{\Psi}}
\newunicodechar{Ω}{\ensuremath{\Omega}}
\newunicodechar{↦}{\ensuremath{\mapsto}}
\newunicodechar{∈}{\ensuremath{\in}}
\newunicodechar{∉}{\ensuremath{\notin}}
\newunicodechar{∧}{\ensuremath{\wedge}}
\newunicodechar{⇔}{\ensuremath{\Leftrightarrow}}
\newunicodechar{⟺}{\ensuremath{\Longleftrightarrow}}
\newunicodechar{⇒}{\ensuremath{\Rightarrow}}
\newunicodechar{⟹}{\ensuremath{\Longrightarrow}}
\newunicodechar{∘}{\ensuremath{\circ}}
\newunicodechar{∪}{\ensuremath{\cup}}
\newunicodechar{∩}{\ensuremath{\cap}}
\newunicodechar{≡}{\ensuremath{\equiv}}
\newunicodechar{⊂}{\ensuremath{\subset}}
\newunicodechar{⊥}{\ensuremath{\perp}}
\newunicodechar{∃}{\ensuremath{\exists}}
\newunicodechar{∀}{\ensuremath{\forall}}
\newunicodechar{∛}{\ensuremath{\sqrt[3]{\,}}}
\newunicodechar{∜}{\ensuremath{\sqrt[4]{\,}}}
\newunicodechar{⃗}{\ensuremath{{}^{\to}}}
\newunicodechar{ⁿ}{\ifmmode^{n}\else\textsuperscript{\ensuremath{n}}\fi}
\newunicodechar{ˣ}{\ifmmode^{x}\else\textsuperscript{\ensuremath{x}}\fi}
\newunicodechar{ᵇ}{\ifmmode^{b}\else\textsuperscript{\ensuremath{b}}\fi}
\newunicodechar{ʳ}{\ifmmode^{r}\else\textsuperscript{\ensuremath{r}}\fi}
\newunicodechar{ᵘ}{\ifmmode^{u}\else\textsuperscript{\ensuremath{u}}\fi}
\newunicodechar{ʸ}{\ifmmode^{y}\else\textsuperscript{\ensuremath{y}}\fi}
\newunicodechar{ᵃ}{\ifmmode^{a}\else\textsuperscript{\ensuremath{a}}\fi}
\newunicodechar{ᵉ}{\ifmmode^{e}\else\textsuperscript{\ensuremath{e}}\fi}
\newunicodechar{ᵏ}{\ifmmode^{k}\else\textsuperscript{\ensuremath{k}}\fi}
\newunicodechar{ᵖ}{\ifmmode^{p}\else\textsuperscript{\ensuremath{p}}\fi}
\newunicodechar{ᴺ}{\ifmmode^{N}\else\textsuperscript{\ensuremath{N}}\fi}
\newunicodechar{ᶜ}{\ifmmode^{c}\else\textsuperscript{\ensuremath{c}}\fi}
\newunicodechar{ʰ}{\ifmmode^{h}\else\textsuperscript{\ensuremath{h}}\fi}
\newunicodechar{ᵠ}{\ifmmode^{q}\else\textsuperscript{\ensuremath{q}}\fi}
\newunicodechar{ₙ}{\ifmmode_{n}\else\textsubscript{\ensuremath{n}}\fi}
\newunicodechar{ₐ}{\ifmmode_{a}\else\textsubscript{\ensuremath{a}}\fi}
\newunicodechar{ᵢ}{\ifmmode_{i}\else\textsubscript{\ensuremath{i}}\fi}
\newunicodechar{ᵣ}{\ifmmode_{r}\else\textsubscript{\ensuremath{r}}\fi}
\newunicodechar{ₖ}{\ifmmode_{k}\else\textsubscript{\ensuremath{k}}\fi}
\newunicodechar{ᵦ}{\ifmmode_{\beta}\else\textsubscript{\ensuremath{\beta}}\fi}
\newunicodechar{ₓ}{\ifmmode_{x}\else\textsubscript{\ensuremath{x}}\fi}
\newfontfamily\popdisplay{ArchivoBlack-Regular.ttf}[Path=../../../fonts/]
\newfontfamily\poplogo{SpaceGrotesk-Bold.ttf}[Path=../../../fonts/]
\newfontfamily\popsemi{PlusJakartaSans-SemiBold.ttf}[Path=../../../fonts/]
\newfontfamily\popxbold{PlusJakartaSans-ExtraBold.ttf}[Path=../../../fonts/]
\newfontfamily\popxboldls{PlusJakartaSans-ExtraBold.ttf}[Path=../../../fonts/,LetterSpace=3.0]

\setlist[enumerate]{leftmargin=1.7em,itemsep=5pt,topsep=4pt}
\setlist[itemize]{leftmargin=1.7em,itemsep=4pt,topsep=4pt,label={\color{poporange}\textbullet}}
\setlength{\parindent}{0pt}
\setlength{\parskip}{5pt}
\renewcommand{\arraystretch}{1.25}

% pgfplots : style Pop par défaut
\pgfplotsset{
  every axis/.append style={axis line style={popink,line width=1pt},tick style={popink},
    grid style={popline,line width=0.5pt},label style={font=\small},tick label style={font=\footnotesize}},
  popcurve/.style={poporange,line width=1.8pt,no markers,smooth},
}
% Même style disponible dans un \draw TikZ simple (hors pgfplots)
\tikzset{popcurve/.style={poporange,line width=1.8pt}}
\tikzset{popgrid/.style={popline,very thin}}
\colorlet{popcurve}{poporange} % le nom du style est parfois utilisé comme couleur

% ============================================================
% Pill / squiggle
% ============================================================
\newcommand{\poppill}[3]{%
  \tikz[baseline=-0.55ex]\node[draw=popink,line width=1.2pt,rounded corners=2.6pt,%
    fill=#2,inner xsep=6.5pt,inner ysep=3pt]{%
    \popxboldls\fontsize{7}{7}\selectfont\color{#3}\MakeUppercase{#1}};%
}
\newcommand{\popsquiggle}[1]{%
  \tikz[baseline=-0.4ex]{
    \draw[#1,line width=2.6pt,line cap=round]
      plot [smooth] coordinates {(0,0.09) (0.34,0.01) (0.68,0.21) (1.02,0.01) (1.36,0.10)};
  }%
}
% Mot-clé surligné (marqueur orange sous le texte)
\newcommand{\cle}[1]{{\popxbold\color{popdark}#1}}

% ============================================================
% En-tête / pied
% ============================================================
\pagestyle{fancy}
\fancyhf{}
\renewcommand{\headrulewidth}{0pt}
\renewcommand{\footrulewidth}{0pt}
\lhead{\raisebox{-0.15ex}{\poplogo\fontsize{13}{13}\selectfont\color{popink}OnBuch\color{poporange}+}}
\rhead{\poppill{\DOCMATIERE\ \textbullet\ \DOCNIVEAU\ \textbullet\ Leçon \DOCLECON}{white}{popink}}
\lfoot{\popxbold\fontsize{7.6}{7.6}\selectfont\color{popmuted}\MakeUppercase{Cours signé OnBuch+}\ \textbullet\ {\popsemi\DOCTITRE}}
\rfoot{\tikz[baseline=-0.6ex]\node[rounded corners=8.5pt,fill=popink,minimum height=17pt,minimum width=17pt,inner xsep=5pt,inner sep=0]{\popxbold\fontsize{8}{8}\selectfont\color{popcream}\thepage\,/\,\pageref*{LastPage}};}

% ============================================================
% Titres
% ============================================================
\newcommand{\chaptitre}[2]{%
  \begin{tcolorbox}[enhanced,colback=poporange,colframe=popink,boxrule=2.6pt,arc=11pt,
      drop shadow=popink,left=15pt,right=15pt,top=12pt,bottom=11pt,before skip=3pt,after skip=18pt]
    \poppill{#2}{popink}{popcream}\par\vspace{9pt}
    {\raggedright\hyphenpenalty=10000\exhyphenpenalty=10000\popdisplay\fontsize{22}{24}\selectfont\color{popcream}\MakeUppercase{#1}\par}
  \end{tcolorbox}
}
% Section numérotée : pastille orange avec le numéro + titre Archivo
\newcounter{popsec}
\newcommand{\coursec}[1]{%
  \stepcounter{popsec}\setcounter{popsub}{0}%
  \par\vspace{16pt}\noindent
  \begin{minipage}{\linewidth}
  \tikz[baseline=-0.7ex]\node[draw=popink,line width=1.6pt,fill=poporange,rounded corners=4pt,minimum size=22pt,inner sep=0]{\popdisplay\fontsize{12}{12}\selectfont\color{popcream}\thepopsec};%
  \hspace{8pt}{\popdisplay\fontsize{15}{17}\selectfont\color{popink}\MakeUppercase{#1}}\par
  \vspace{4pt}\noindent{\color{poporange}\rule{36pt}{3.6pt}}
  \end{minipage}\par\nopagebreak\vspace{8pt}
}
\newcounter{popsub}
\newcommand{\courssub}[1]{%
  \stepcounter{popsub}%
  \par\vspace{9pt}\noindent{\popxbold\fontsize{11.5}{13}\selectfont\color{popink}{\color{poporange}\thepopsec.\thepopsub}\hspace{6pt}#1}\par\nopagebreak\vspace{4pt}
}

% ============================================================
% Boîtes pédagogiques — même recette : fond blanc, bordure épaisse colorée,
% ombre dure encre, badge plein. Titre optionnel [#1].
% ============================================================
\tcbset{popbase/.style={breakable,enhanced,colback=white,boxrule=2.3pt,arc=9pt,
  shadow={3pt}{-3pt}{0pt}{popink},left=14pt,right=14pt,top=11pt,bottom=11pt,before skip=11pt,after skip=15pt,
  fontupper=\popsemi\fontsize{10.6}{14.5}\selectfont\color{popink},
  fontlower=\popsemi\fontsize{10.6}{14.5}\selectfont\color{popink}}}
% Coche dessinée (le glyphe ✓ n'existe pas dans les polices du document)
\newcommand{\popcheck}[1]{\tikz[baseline=-0.5ex]\draw[#1,line width=1.6pt,line cap=round,line join=round](-0.13,0.0)--(-0.04,-0.09)--(0.13,0.1);}
\providecommand{\coche}{\popcheck{popgreen}}
\providecommand{\resultat}[1]{\par\smallskip\noindent\fcolorbox{popgreen}{popgreenL}{\parbox{\dimexpr\linewidth-2\fboxsep-2\fboxrule\relax}{\textbf{Résultat.} #1}}\par\smallskip}
\newcommand{\popboxhead}[3]{\poppill{#1}{#2}{white}\ifx\relax#3\relax\else\hspace{6pt}{\popxbold\fontsize{10.6}{12}\selectfont\color{popink}#3}\fi\par\vspace{7pt}}

\newtcolorbox{aretenir}[1][]{popbase,colframe=popgreen,before upper={\popboxhead{À retenir}{popgreen}{#1}}}
\newtcolorbox{exemplebox}[1][]{popbase,colframe=poporange,before upper={\popboxhead{Exemple}{poporange}{#1}}}
\newtcolorbox{definition}[1][]{popbase,colframe=popblue,before upper={\popboxhead{Définition}{popblue}{#1}}}
\newtcolorbox{propriete}[1][]{popbase,colframe=poppurple,before upper={\popboxhead{Propriété}{poppurple}{#1}}}
\newtcolorbox{methode}[1][]{popbase,colframe=popgold,before upper={\popboxhead{Méthode}{popgold}{#1}}}
\newtcolorbox{attention}[1][]{popbase,colframe=poppink,before upper={\popboxhead{Attention}{poppink}{#1}}}
\newtcolorbox{experience}[1][]{popbase,colframe=popblue,colback=popblueL,before upper={\popboxhead{Expérience}{popblue}{#1}}}
% « Le savais-tu ? » : carte violette pleine (contexte, culture, Cameroun)
\newtcolorbox{savaistu}[1][]{popbase,colback=poppurple,colframe=popink,
  fontupper=\popsemi\fontsize{10.4}{14.2}\selectfont\color{white},
  before upper={\poppill{Le savais-tu\,?}{white}{poppurple}\ifx\relax#1\relax\else\hspace{6pt}{\popxbold\color{white}#1}\fi\par\vspace{7pt}}}
% Exercice résolu : énoncé en haut, \tcblower puis la solution sur fond crème
\newcounter{popexo}\newcounter{popexr}
\newtcolorbox{exoresolu}[1][]{popbase,colframe=poporange,colbacklower=popcream,
  segmentation style={popink,line width=1.2pt,dashed},
  before upper={\stepcounter{popexr}\popboxhead{Exercice résolu \thepopexr}{poporange}{#1}},
  before lower={\poppill{Solution}{popink}{popcream}\par\vspace{6pt}}}
% Exercice (non résolu en place) — la correction est dans la partie Corrigés
\newtcolorbox{exercice}[1][]{popbase,colframe=popink,
  before upper={\stepcounter{popexo}\popboxhead{Exercice \thepopexo}{popink}{#1}}}
% Corrigé
\newtcolorbox{corrige}[1][]{popbase,colframe=popgreen,colback=popgreenL,
  before upper={\popboxhead{Corrigé}{popgreen}{#1}}}
% Objectifs / prérequis / bilan
\newtcolorbox{objectifs}{popbase,colframe=popink,colback=poporangeL,
  before upper={\popboxhead{Objectifs de la leçon}{popink}{}}}
\newtcolorbox{prerequis}{popbase,colframe=popink,colback=popblueL,
  before upper={\popboxhead{Prérequis}{popblue}{}}}
\newtcolorbox{bilan}{popbase,colframe=popink,colback=white,boxrule=2.8pt,
  before upper={\popboxhead{Fiche bilan}{poporange}{L'essentiel à connaître pour l'examen}}}
% Figure encadrée avec légende
\newtcolorbox{popfigure}[1][]{enhanced,breakable=false,colback=white,colframe=popline,boxrule=1.4pt,arc=8pt,
  left=10pt,right=10pt,top=10pt,bottom=8pt,before skip=10pt,after skip=12pt,halign=center,
  lower separated=false,halign lower=center,
  fontlower=\popsemi\fontsize{9}{11.5}\selectfont\color{popmuted},
  #1}
\newcommand{\legende}[1]{\par\vspace{6pt}{\popsemi\fontsize{9}{11.5}\selectfont\color{popmuted}#1}}

% ============================================================
% Signature OnBuch+
% ============================================================
\newcommand{\poplogoinline}[1]{{\poplogo\fontsize{#1}{#1}\selectfont\color{popink}OnBuch\color{poporange}+}}
\newcommand{\signatureonbuch}{%
  \begin{tcolorbox}[enhanced,colback=popink,colframe=popink,boxrule=2.2pt,arc=10pt,
      drop shadow=poporange,left=14pt,right=14pt,top=10pt,bottom=10pt,before skip=0pt,after skip=0pt]
    \tikz[baseline=-0.7ex]\node[circle,fill=popgreen,draw=popcream,line width=1.3pt,minimum size=22pt,inner sep=0]{\popcheck{white}};%
    \hspace{9pt}\begin{minipage}[c]{0.62\linewidth}
      {\popxbold\fontsize{12}{13}\selectfont\color{popcream}Signé\ \textbullet\ {\poplogo OnBuch\color{poporange}+}}\par\vspace{2pt}
      {\raggedright\popsemi\fontsize{8}{10.5}\selectfont\color{popline}\MakeUppercase{Équipe pédagogique OnBuch+}\\\MakeUppercase{Conforme au programme officiel MINESEC}\par}
    \end{minipage}\hfill
    {\poplogo\fontsize{22}{22}\selectfont\color{popcream}OnBuch\color{poporange}+}
  \end{tcolorbox}%
}

% ============================================================
% Page de garde
% #1 titre · #2 matière · #3 niveau · #4 module · #5 n° leçon · #6 durée ·
% #7 description / objectifs (texte ou itemize)
% ============================================================
\newcommand{\pagegarde}[7]{%
  \thispagestyle{empty}
  \newgeometry{a4paper,margin=1.7cm,top=1.6cm,bottom=1.4cm}%
  \begin{tikzpicture}[remember picture,overlay]
    \fill[poporange!18] ([xshift=-3.2cm,yshift=-3.6cm]current page.north east) circle (4.6cm);
    \fill[poppurple!14] ([xshift=2.4cm,yshift=7.6cm]current page.south west) circle (3.8cm);
  \end{tikzpicture}%
  {\poplogo\fontsize{30}{30}\selectfont\color{popink}OnBuch\color{poporange}+}\hfill
  \raisebox{0.6ex}{\poppill{Cours signé \textbullet\ Premium}{popink}{popcream}}\par
  \vspace{1pt}\popsquiggle{poppurple}\par
  \vspace{8pt}
  \begin{tcolorbox}[enhanced,colback=poporange,colframe=popink,boxrule=2.8pt,arc=13pt,
      drop shadow=popink,left=18pt,right=18pt,top=12pt,bottom=13pt,after skip=10pt]
    \poppill{#2 \textbullet\ #3}{popink}{popcream}\hspace{5pt}\poppill{#4}{white}{popink}\par\vspace{8pt}
    {\popdisplay\fontsize{14}{14}\selectfont\color{popink}LEÇON #5}\par\vspace{4pt}
    {\raggedright\hyphenpenalty=10000\exhyphenpenalty=10000\popdisplay\fontsize{25}{27}\selectfont\color{popcream}\MakeUppercase{#1}\par}
    \vspace{8pt}
    {\popxbold\fontsize{9.5}{9.5}\selectfont\color{popink}\MakeUppercase{Durée conseillée : #6}}
  \end{tcolorbox}
  \begin{tcolorbox}[enhanced,colback=white,colframe=popink,boxrule=2.2pt,arc=10pt,
      drop shadow=popink,left=16pt,right=16pt,top=10pt,bottom=10pt,after skip=8pt]
    \poppill{Dans cette leçon}{poppurple}{white}\par\vspace{5pt}
    \popsemi\fontsize{9.3}{12.6}\selectfont\color{popink}#7
  \end{tcolorbox}
  \noindent\begin{minipage}[t]{0.487\linewidth}
    \begin{tcolorbox}[enhanced,colback=popgreen,colframe=popink,boxrule=2.2pt,arc=10pt,
        drop shadow=popink,left=11pt,right=11pt,top=10pt,bottom=10pt,height=3.1cm,valign=center]
      \tcbox[colback=white,colframe=popink,boxrule=1.3pt,arc=5pt,boxsep=0pt,nobeforeafter,
        left=3pt,right=3pt,top=3pt,bottom=3pt,valign=center]{\qrcode[height=1.9cm]{https://play.google.com/store/apps/details?id=com.luvvix.onbuch&hl=fr}}%
      \hspace{8pt}\begin{minipage}[c]{0.5\linewidth}
        {\popdisplay\fontsize{11}{12.5}\selectfont\color{white}\MakeUppercase{Télécharge OnBuch}}\par\vspace{4pt}
        {\raggedright\popsemi\fontsize{8.4}{11}\selectfont\color{white}Gratuit \textbullet\ Google Play\\Tuteur IA, annales, concours, résultats\par}
      \end{minipage}
    \end{tcolorbox}
  \end{minipage}\hfill
  \begin{minipage}[t]{0.487\linewidth}
    \begin{tcolorbox}[enhanced,colback=poppurple,colframe=popink,boxrule=2.2pt,arc=10pt,
        drop shadow=popink,left=11pt,right=11pt,top=10pt,bottom=10pt,height=3.1cm,valign=center]
      \tikz[baseline=-0.7ex]\node[circle,fill=white,draw=popink,line width=1.3pt,minimum size=1.6cm,inner sep=0]{\popdisplay\fontsize{24}{24}\selectfont\color{poppurple}+};%
      \hspace{8pt}\begin{minipage}[c]{0.6\linewidth}
        {\popdisplay\fontsize{11}{12.5}\selectfont\color{white}\MakeUppercase{Passe à OnBuch+}}\par\vspace{4pt}
        {\raggedright\popsemi\fontsize{8.4}{11}\selectfont\color{white}Tous les cours signés, Tuteur IA illimité, fiches PDF — onglet OnBuch+ de l'app\par}
      \end{minipage}
    \end{tcolorbox}
  \end{minipage}
  \vspace{8pt}
  \popcontenu
  \vfill
  \signatureonbuch
  \restoregeometry
  \newpage
}

% Rangée « Ce cours contient » (page de garde)
\newcommand{\poptile}[3]{% #1 couleur #2 gros texte #3 légende
  \begin{tcolorbox}[enhanced,colback=#1,colframe=popink,boxrule=2pt,arc=9pt,shadow={2.5pt}{-2.5pt}{0pt}{popink},
      left=6pt,right=6pt,top=9pt,bottom=9pt,halign=center,valign=center,height=1.75cm,width=\linewidth,nobeforeafter]
    {\popdisplay\fontsize{12.5}{13}\selectfont\color{white}\MakeUppercase{#2}}\par\vspace{3pt}
    {\centering\popsemi\fontsize{7.8}{9.5}\selectfont\color{white}#3\par}
  \end{tcolorbox}}
\newcommand{\popcontenu}{%
  \par\noindent{\popxboldls\fontsize{7.5}{7.5}\selectfont\color{popmuted}CE COURS SIGNÉ CONTIENT}\par\vspace{7pt}
  \noindent\begin{minipage}[t]{0.235\linewidth}\poptile{popblue}{Cours}{complet, illustré, pas à pas}\end{minipage}\hfill
  \begin{minipage}[t]{0.235\linewidth}\poptile{poporange}{Méthodes}{et exercices résolus}\end{minipage}\hfill
  \begin{minipage}[t]{0.235\linewidth}\poptile{poppink}{Exercices}{type examen + corrigés}\end{minipage}\hfill
  \begin{minipage}[t]{0.235\linewidth}\poptile{popgreen}{Fiche bilan}{l'essentiel à retenir}\end{minipage}\par}

% Sommaire
\newtcolorbox{sommaire}{enhanced,colback=white,colframe=popink,boxrule=2.2pt,arc=10pt,
  drop shadow=popink,left=18pt,right=18pt,top=13pt,bottom=13pt,before skip=6pt,after skip=20pt,
  before upper={\poppill{Sommaire}{poppurple}{white}\par\vspace{9pt}\popsemi\fontsize{10.6}{14.5}\selectfont\color{popink}}}

% Signature de fin de cours — poussée tout en bas de la dernière page
\newcommand{\findecours}{%
  \par\vfill\nopagebreak
  \begin{center}\popsquiggle{poporange}\end{center}\vspace{4pt}
  \signatureonbuch
}
\AtBeginDocument{\def\llbracket{\mathopen{[\mkern-3.5mu[}}\def\rrbracket{\mathclose{]\mkern-3.5mu]}}} % intervalles d'entiers (glyphe ⟦ absent de la police)
% Glyphes absents de Fira Math : redéfinis à partir de glyphes disponibles
\AtBeginDocument{%
  \def\setminus{\mathbin{\backslash}}\let\smallsetminus\setminus
  \def\vdots{\mathord{\vbox{\baselineskip3.2pt\lineskiplimit0pt\kern4pt\hbox{.}\hbox{.}\hbox{.}}}}%
  \def\ddots{\mathinner{\mkern1mu\raise6.5pt\hbox{.}\mkern2mu\raise3.6pt\hbox{.}\mkern2mu\raise0.7pt\hbox{.}\mkern1mu}}%
  \def\triangle{\mathord{\scalebox{0.9}{$\symup{\Delta}$}}}%
  \def\nabla{\mathord{\raisebox{\depth}{\scalebox{1}[-1]{$\symup{\Delta}$}}}}%
  \def\checkmark{\popcheck{popdark}}%
  \def\star{\mathbin{\ast}}\def\bigstar{\mathord{\ast}}%
  \def\diamond{\mathbin{\scalebox{0.6}{$\Diamond$}}}%
  \def\bigcup{\mathop{\mathchoice{\raisebox{-0.3em}{\scalebox{1.7}{$\cup$}}}{\scalebox{1.25}{$\cup$}}{\cup}{\cup}}\displaylimits}%
  \def\bigcap{\mathop{\mathchoice{\raisebox{-0.3em}{\scalebox{1.7}{$\cap$}}}{\scalebox{1.25}{$\cap$}}{\cap}{\cap}}\displaylimits}%
  \def\lesseqqgtr{\mathrel{\lessgtr}}\def\bigoplus{\mathop{\oplus}\displaylimits}%
}
\providecommand{\ssi}{si et seulement si} % abréviation fréquente
\AtBeginDocument{\providecommand{\wideparen}{\overparen}} % arc de cercle

\begin{document}

\pagegarde{Les mécanismes de l'immunité et ses dysfonctionnements}{SVTEEHB}{Tle TI}{Module II : L'Éducation à la Santé}{N06}{5 h}{Cette leçon explore l'origine des cellules immunitaires, les mécanismes des réponses immunitaires non spécifiques et spécifiques (cellulaire et humorale) ainsi que la mémoire immunitaire, fondement de la vaccination. Elle analyse ensuite les dysfonctionnements du système immunitaire (maladies auto-immunes, VIH/Sida) et développe des compétences de prévention et de solidarité face au Sida.
\vspace{4pt}
\begin{multicols}{2}\small
\begin{itemize}
\item À la fin de cette leçon, je sais déterminer les lieux de formation et de maturation des cellules immunitaires.
\item À la fin de cette leçon, je sais identifier les structures de reconnaissance du soi et du non-soi (CMH, récepteurs des lymphocytes).
\item À la fin de cette leçon, je sais rappeler les mécanismes de la réponse immunitaire non spécifique (barrières, phagocytose, inflammation).
\item À la fin de cette leçon, je sais interpréter des résultats d'expériences mettant en évidence la réponse à médiation cellulaire (greffes, destruction des cellules cibles).
\item À la fin de cette leçon, je sais interpréter des résultats d'expériences mettant en évidence la réponse à médiation humorale (production d'anticorps, sérothérapie).
\item À la fin de cette leçon, je sais exploiter des courbes de cinétique des anticorps pour expliquer la mémoire immunitaire et le principe de la vaccination.
\end{itemize}\end{multicols}}

\begin{sommaire}
\begin{multicols}{2}
\begin{enumerate}
\item Origine et maturation des cellules immunitaires
\item La reconnaissance du soi et du non-soi
\item La réponse immunitaire non spécifique (rappels)
\item La réponse immunitaire spécifique : médiation cellulaire et humorale
\item La mémoire immunitaire et la vaccination
\item Les dysfonctionnements : maladies auto-immunes
\item Le VIH/Sida : mécanismes, prévention et solidarité
\item Activité d'intégration
\item Méthodes et exercices résolus
\item Exercices
\item Corrigés des exercices
\item Fiche bilan
\end{enumerate}
\end{multicols}
\end{sommaire}

\begin{objectifs}
\begin{itemize}
\item À la fin de cette leçon, je sais déterminer les lieux de formation et de maturation des cellules immunitaires.
\item À la fin de cette leçon, je sais identifier les structures de reconnaissance du soi et du non-soi (CMH, récepteurs des lymphocytes).
\item À la fin de cette leçon, je sais rappeler les mécanismes de la réponse immunitaire non spécifique (barrières, phagocytose, inflammation).
\item À la fin de cette leçon, je sais interpréter des résultats d'expériences mettant en évidence la réponse à médiation cellulaire (greffes, destruction des cellules cibles).
\item À la fin de cette leçon, je sais interpréter des résultats d'expériences mettant en évidence la réponse à médiation humorale (production d'anticorps, sérothérapie).
\item À la fin de cette leçon, je sais exploiter des courbes de cinétique des anticorps pour expliquer la mémoire immunitaire et le principe de la vaccination.
\item À la fin de cette leçon, je sais expliquer les mécanismes des maladies auto-immunes à partir d'exemples (vitiligo, diabète juvénile, myasthénie).
\item À la fin de cette leçon, je sais expliquer le mode d'action du VIH sur les lymphocytes T4 et les conséquences de l'immunodéficience.
\item À la fin de cette leçon, je sais produire des outils de sensibilisation à la prévention du VIH/Sida et adopter des comportements solidaires envers les personnes vivant avec le VIH.
\end{itemize}
\end{objectifs}

\begin{prerequis}
\begin{itemize}
\item Les cellules de l'organisme : structure générale, notion de membrane plasmique et de marqueurs membranaires (Seconde/Première).
\item Les groupes sanguins et la notion d'antigène/anticorps (Première).
\item La réaction immunitaire non spécifique : phagocytose et réaction inflammatoire (Première).
\item Les protéines : structure et spécificité d'action (enzymes, anticorps) (Première).
\item La notion de virus et de cycle de multiplication virale (Première).
\end{itemize}
\end{prerequis}

\newpage

\coursec{Origine et maturation des cellules immunitaires}

Au service des urgences de l'hôpital de district de Bafoussam, deux patients arrivent presque simultanément. Le premier, un adolescent de 14 ans, est en détresse métabolique aiguë (cétoacidose diabétique) : son pancréas ne produit plus d'insuline, ses îlots de Langerhans ayant été détruits par ses propres lymphocytes. Le second, un jeune père de famille séropositif depuis trois ans, non observant de son traitement antirétroviral, succombe à une pneumonie à \emph{Pneumocystis jirovecii}, un champignon microscopique que tout système immunitaire sain maintient en silence. Comment le même arsenal de défense peut-il, chez l'un, se retourner contre l'organisme (maladie auto-immune), et chez l'autre, s'effondrer face à des agents opportunistes (immunodéficience) ? Pourquoi la vaccination antivariolique (jadis obligatoire) ou antiamarile (encore recommandée au Cameroun) confère-t-elle une protection à vie, alors qu'une simple plaie infectée guérit sans laisser de mémoire spécifique ? Pour répondre, il faut remonter à la source : d'où viennent nos soldats de l'immunité, où s'entraînent-ils, et comment sont-ils déployés ?

\courssub{Les cellules immunitaires : des leucocytes aux fonctions spécialisées}

\begin{definition}[Leucocyte et lignées cellulaires]
Un \cle{leucocyte} (ou globule blanc) est une cellule nucléée du sang et de la lymphe, acteur de la défense de l'organisme. On distingue trois grandes lignées aux morphologies et fonctions distinctes :
\begin{itemize}
    \item Les \cle{polynucléaires} (neutrophiles, éosinophiles, basophiles) : cytoplasme granuleux, noyau multilobé. Premiers intervenants de l'inflammation (phagocytose pour les neutrophiles, dégranulation pour les basophiles/éosinophiles).
    \item Les \cle{monocytes} (grands noyaux réniformes) qui, dans les tissus, se différencient en \cle{macrophages} : phagocytose massive, présentation d'antigènes, sécrétion de cytokines.
    \item Les \cle{lymphocytes} (noyau rond occupant quasi tout le volume, cytoplasme mince) : cellules de l'immunité \emph{spécifique}. Deux sous-types majeurs : les \cle{lymphocytes B} (LB) et les \cle{lymphocytes T} (LT).
\end{itemize}
\end{definition}

\begin{popfigure}
\begin{tikzpicture}[
    scale=0.85, every node/.style={transform shape},
    stem/.style={draw=popdark, thick, fill=popcream, rounded corners=4pt, minimum width=3cm, minimum height=1.1cm, font=\small\bfseries, align=center},
    prog/.style={draw=popblue, thick, fill=popblueL, rounded corners=4pt, minimum width=2.8cm, minimum height=0.9cm, font=\small, align=center},
    mature/.style={draw=popgreen, thick, fill=popgreenL, rounded corners=4pt, minimum width=2.6cm, minimum height=0.85cm, font=\small, align=center},
    arrow/.style={-Stealth, thick, popdark},
    dasharrow/.style={-Stealth, thick, poporange, dashed},
]
% Niveaux y
\def\yStem{6}
\def\yProg{3.5}
\def\yMature{0.5}

% Cellule souche
\node[stem, align=center] (HSC) at (0,\yStem){Cellule souche hématopoïétique\\$($CD34$^+$$)$\\Moelle osseuse rouge};

% Progéniteurs
\node[prog, align=center] (CMP) at (-5.5,\yProg){Progéniteur myéloïde commun\\ (CMP)};
\node[prog, align=center] (CLP) at (5.5,\yProg){Progéniteur lymphoïde commun\\ (CLP)};

% Lignée myéloïde
\node[mature, align=center] (GMP) at (-8,\yMature){Progéniteur granulocytaire-\\monocytaire (GMP)};
\node[mature] (Mono) at (-11,-2) {Monocyte $\rightarrow$ Macrophage};
\node[mature] (PolyN) at (-8,-2) {Polynucléaire neutrophile};
\node[mature] (PolyE) at (-5,-2) {Polynucléaire éosinophile};
\node[mature] (PolyB) at (-2,-2) {Polynucléaire basophile};

% Lignée lymphoïde
\node[mature] (preB) at (3,\yMature) {Pré-lymphocyte B};
\node[mature, align=center] (LB) at (3,-2){Lymphocyte B mature\\ (Moelle osseuse)};
\node[mature] (preT) at (8,\yMature) {Pré-lymphocyte T};
\node[mature] (Thymus) at (8,-2) {Maturation dans le \textbf{Thymus}};
\node[mature] (LT4) at (11,-2) {LT CD4$^+$ (Auxiliaire)};
\node[mature] (LT8) at (8,-4) {LT CD8$^+$ (Cytotoxique)};

% Flèches
\draw[arrow] (HSC.south west) -- (CMP.north);
\draw[arrow] (HSC.south east) -- (CLP.north);

\draw[arrow] (CMP.south west) -- (GMP.north);
\draw[arrow] (GMP.south west) -- (Mono.north);
\draw[arrow] (GMP.south) -- (PolyN.north);
\draw[arrow] (GMP.south east) -- (PolyE.north);
\draw[arrow] (GMP.east) -- (PolyB.west);

\draw[arrow] (CLP.south west) -- (preB.north);
\draw[arrow] (preB.south) -- (LB.north);

\draw[arrow] (CLP.south east) -- (preT.north);
\draw[dasharrow] (preT.south) -- (Thymus.north); % Migration vers thymus
\draw[arrow] (Thymus.south east) -- (LT4.north west);
\draw[arrow] (Thymus.south west) -- (LT8.north);

% Annotations
\node[align=center, font=\small\itshape, popmuted] at (-11,-3.2) {Immunité non spécifique};
\node[align=center, font=\small\itshape, popmuted] at (3,-3.2) {Immunité humorale};
\node[align=center, font=\small\itshape, popmuted] at (9.5,-4.5) {Immunité cellulaire};

% Légende boîtes
\node[draw=popdark, fill=popcream, rounded corners=2pt, inner sep=2pt, font=\tiny, anchor=north west] at (-11.5,6.5) {Cellule souche};
\node[draw=popblue, fill=popblueL, rounded corners=2pt, inner sep=2pt, font=\tiny, anchor=north west] at (-11.5,6.0) {Progéniteur engagé};
\node[draw=popgreen, fill=popgreenL, rounded corners=2pt, inner sep=2pt, font=\tiny, anchor=north west] at (-11.5,5.5) {Cellule effectrice mature};
\node[draw=poporange, dashed, rounded corners=2pt, inner sep=2pt, font=\tiny, anchor=north west] at (-11.5,5.0) {Migration / Maturation};
\end{tikzpicture}
\legende{Schéma de l'hématopoïèse humaine : différenciation des cellules souches de la moelle osseuse rouge vers les effecteurs de l'immunité. Les flèches pleines indiquent la différenciation ; la flèche pointillée la migration des pré-LT vers le thymus pour maturation.}
\end{popfigure}

\courssub{L'origine commune : la moelle osseuse rouge, lieu de formation de tous les leucocytes}

\begin{propriete}[Origine hématopoïétique]
Toutes les cellules immunitaires (polynucléaires, monocytes/macrophages, lymphocytes B et T) dérivent d'une \cle{cellule souche hématopoïétique} (CSH, phénotype CD34$^+$) située dans la \cle{moelle osseuse rouge} (os plats : sternum, iliaques, côtes, vertèbres ; épiphyses des os longs). Cette origine commune est un savoir admis, solidement établi par l'expérience d'irradiation-greffe.
\end{propriete}

\begin{experience}[Mise en évidence du rôle de la moelle osseuse : irradiation-greffe chez la souris]
\textbf{Problématique :} La moelle osseuse est-elle le lieu de formation des cellules immunitaires ?
\textbf{Matériel :} Souris isogéniques (même patrimoine génétique), source de rayonnements ionisants (dose létale : \SI{6}{\gray} à \SI{9}{\gray}), moelle osseuse de donneur sain.
\textbf{Protocole :}
\begin{enumerate}
    \item \textbf{Groupe Témoin (T) :} Souris non irradiée, non greffée.
    \item \textbf{Groupe Irradié (I) :} Souris irradiée (destruction de la moelle), non greffée.
    \item \textbf{Groupe Greffé (G) :} Souris irradiée $+$ injection intraveineuse de cellules de moelle osseuse de donneur sain (\num{1e6} à \num{1e7} cellules).
\end{enumerate}
\textbf{Observations (à J+14) :}
\begin{center}
\begin{tabularx}{\linewidth}{|l|X|X|X|}\hline
\textbf{Paramètre} & \textbf{Témoin (T)} & \textbf{Irradié (I)} & \textbf{Greffé (G)} \\ \hline
Survie & 100\% & 0\% (mort J+7 à J+12) & 90\% \\
Numération leucocytaire & Normale (\SI{5e3}{\per\micro\liter}--\SI{1e4}{\per\micro\liter}) & Aplasie profonde ($<$ \SI{500}{\per\micro\liter}) & Restituée à la normale \\
Réponse immunitaire (test antigène) & Positive & Absente & Positive \\ \hline
\end{tabularx}
\end{center}
\textbf{Interprétation :} L'irradiation détruit la moelle $\rightarrow$ aplasie médullaire $\rightarrow$ disparition des leucocytes $\rightarrow$ mort par infection (absence de défense). La greffe de moelle restaure la production de leucocytes et la compétence immunitaire.
\textbf{Conclusion :} La moelle osseuse est le site de \textbf{formation} de toutes les cellules immunitaires. C'est le principe de la greffe de moelle osseuse (allogreffe) utilisée en thérapeutique.
\end{experience}

\begin{methode}[Interpréter une expérience d'irradiation-greffe de moelle]
Pour conclure sur le lieu de formation des cellules immunitaires à partir d'un tel protocole :
\begin{enumerate}
    \item \textbf{Identifier la variable indépendante :} l'irradiation (destruction moelle) et la greffe (apport cellules moelle).
    \item \textbf{Comparer les groupes :} Témoin vs Irradié (effet de la destruction) ; Irradié vs Greffé (effet de la restauration).
    \item \textbf{Analyser les paramètres mesurés :} survie, numération leucocytaire (formule sanguine), capacité de réponse immunitaire (ex. production d'anticorps après injection d'antigène).
    \item \textbf{Conclure :} Si la destruction de la moelle abolit l'immunité et que sa restauration par greffe la rétablit, la moelle est le lieu de formation des cellules effectrices de l'immunité.
\end{enumerate}
\end{methode}

\begin{savaistu}[La greffe de moelle osseuse au Cameroun]
La greffe de moelle osseuse (ou de cellules souches hématopoïétiques) est une thérapie curative pour certaines leucémies aiguës, lymphomes, aplasies médullaires sévères ou drépanocytose majeure. Au Cameroun, bien que techniquement lourde (nécessite unité d'hématologie, conditionnement myéloablatif, isolement stérile, donneur HLA-compatible), elle commence à être réalisée dans des centres de référence (ex. Hôpital Général de Yaoundé, Centre Mère-Enfant de la Fondation Chantal Biya). Le donneur est souvent un frère ou une sœur HLA-identique (probabilité 25\% par fratrie). Ce savoir de Terminale C/TI est donc directement à la base d'une thérapeutique de pointe accessible dans notre pays.
\end{savaistu}

\courssub{La maturation différentielle : moelle pour les LB, thymus pour les LT}

Si la \textbf{formation} (prolifération, engagement lignager) a lieu dans la moelle pour tous, la \textbf{maturation} (acquisition des récepteurs fonctionnels, sélection du répertoire non-auto-réactif) sépare les deux lignées lymphoïdes.

\begin{definition}[Organes lymphoïdes primaires et secondaires]
\begin{itemize}
    \item \textbf{Organes lymphoïdes primaires (centraux) :} Sites de formation \emph{et/ou} maturation des lymphocytes à partir de précurseurs. Chez l'Homme : \textbf{moelle osseuse rouge} (formation de tous $+$ maturation des LB) et \textbf{thymus} (maturation des LT).
    \item \textbf{Organes lymphoïdes secondaires (périphériques) :} Sites de rencontre entre lymphocytes matures naïfs et antigènes, et de déclenchement de la réponse adaptative. Ganglions lymphatiques, rate, amygdales, plaques de Peyer (iléon), appendice.
\end{itemize}
\end{definition}

\begin{experience}[Rôle du thymus dans la maturation des LT : thymectomie néonatale chez la souris]
\textbf{Problématique :} Quel est le rôle du thymus dans l'immunité ?
\textbf{Protocole (Miller, 1961) :} Thymectomie chirurgicale chez la souris \textbf{nouveau-née} (J+1 à J+3) vs souris sham-opérée (témoin). Analyse à l'âge adulte.
\textbf{Résultats :}
\begin{itemize}
    \item Atrophie des zones lymphoïdes périphériques (ganglions, rate : zones T dépourvues de lymphocytes).
    \item Déficit sévère de la \textbf{réponse à médiation cellulaire} (rejet de greffe allogénique, hypersensibilité retardée type tuberculine).
    \item Déficit en \textbf{production d'anticorps} vis-à-vis des antigènes \emph{thymodépendants} (protéines), mais conservation de la réponse aux antigènes \emph{thymoindépendants} (polysaccharides bactériens).
    \item Sensibilité accrue aux infections virales, fongiques, intracellulaires.
\end{itemize}
\textbf{Interprétation :} L'ablation précoce du thymus empêche la maturation des pré-LT en LT fonctionnels (CD4$^+$ et CD8$^+$). Les LT auxiliaires (CD4$^+$) étant nécessaires pour aider les LB à produire des anticorps de classe IgG (commutation de classe), leur absence explique le déficit humorale partiel.
\textbf{Conclusion :} Le thymus est l'organe \textbf{indispensable à la maturation des lymphocytes T}. Sans thymus, pas de LT matures $\rightarrow$ pas d'immunité cellulaire $\rightarrow$ aide auxiliaire absente pour les LB $\rightarrow$ immunité humorale altérée.
\end{experience}

\begin{attention}[Piège classique : Formation $\neq$ Maturation]
\begin{itemize}
    \item \textbf{Formation (hématopoïèse) :} Moelle osseuse rouge \textbf{POUR TOUS} (polynucléaires, monocytes, LB, LT).
    \item \textbf{Maturation (acquisition compétences fonctionnelles) :}
    \begin{itemize}
        \item LB : \textbf{Moelle osseuse} (chez l'Homme et la souris ; \emph{bourse de Fabricius} chez l'oiseau, d'où le "B").
        \item LT : \textbf{Thymus} (d'où le "T").
    \end{itemize}
    \item Ne jamais écrire : « Les LT se forment dans le thymus. » Ils s'y \emph{maturent} seulement. Les pré-LT migrent de la moelle vers le thymus via le sang.
\end{itemize}
\end{attention}

\begin{popfigure}
\begin{tikzpicture}[scale=0.9, every node/.style={transform shape}]
% Silhouette simplifiée
\draw[popdark, line width=1.5pt] (0,0) ellipse (2.5 and 7); % Corps
\draw[popdark, line width=1.5pt] (0,7) circle (1.2); % Tête
% Moelle osseuse (fémurs, sternum, bassin)
\foreach \x/\y in {-1.5,-0.5, 1.5,-0.5, -1.5,-3.5, 1.5,-3.5, 0,5.5} {
    \fill[popred!60] (\x,\y) circle (0.25cm);
    \node[popred, font=\tiny\bfseries, anchor=north] at (\x,\y-0.35) {Moelle};
}
% Thymus
\fill[poppurple!70] (-0.8,6.2) ellipse (0.5 and 0.7);
\node[poppurple, font=\small\bfseries, anchor=south] at (-0.8,5.3) {Thymus};
% Rate
\fill[poporange!80] (1.8,3.5) ellipse (0.6 and 0.4);
\node[poporange, font=\small\bfseries, anchor=west] at (2.4,3.5) {Rate};
% Ganglions (cervicaux, axillaires, inguinaux)
\foreach \x/\y in {-2.2,6.0, 2.2,6.0, -2.5,4.0, 2.5,4.0, -2.0,1.0, 2.0,1.0, -1.8,-2.5, 1.8,-2.5} {
    \fill[popblue!70] (\x,\y) circle (0.18cm);
}
\node[popblue, font=\small\bfseries, anchor=east] at (-2.8,6.0) {Ganglions};
% Flèches migration
\draw[popgreen, thick, -Stealth, bend left=20] (-0.8,5.5) to node[midway, above, font=\tiny, popgreen] {Pré-LT} (0.5,3.5);
\draw[popgreen, thick, -Stealth, bend right=20] (0,5.5) to node[midway, above, font=\tiny, popgreen] {LB, LT matures} (-2.0,4.0);
\draw[popgreen, thick, -Stealth] (0,5.5) -- (0,3.5) node[midway, right, font=\tiny, popgreen] {Vers rate};
\draw[popgreen, thick, -Stealth] (0,-0.5) -- (0,-2.0) node[midway, right, font=\tiny, popgreen] {Vers gangl. inguin.};

% Légende intégrée
\begin{scope}[xshift=5.5cm, yshift=2cm]
\node[draw=popred!60, fill=popred!10, rounded corners=2pt, inner sep=3pt, font=\small, anchor=west] at (0,1.5) {Organes primaires : Moelle (formation tous, maturation LB), Thymus (maturation LT)};
\node[draw=popblue!70, fill=popblue!10, rounded corners=2pt, inner sep=3pt, font=\small, anchor=west] at (0,0.8) {Organes secondaires : Ganglions, Rate, Amygdales, Plaques de Peyer};
\node[draw=popgreen!70, fill=popgreen!10, rounded corners=2pt, inner sep=3pt, font=\small, anchor=west] at (0,0.1) {Flèches vertes : Migration des lymphocytes matures (naïfs) vers organes secondaires};
\end{scope}
\end{tikzpicture}
\legende{Localisation des organes lymphoïdes chez l'Homme. En rouge : moelle osseuse (os plats et épiphyses) et thymus (médiastin antérieur, involue après la puberté). En bleu : organes lymphoïdes secondaires (ganglions, rate). Les flèches vertes schématisent la migration des lymphocytes matures (naïfs) via la circulation sanguine et lymphatique vers les organes secondaires où ils rencontreront l'antigène.}
\end{popfigure}

\courssub{Bilan récapitulatif et exercice résolu}

\begin{center}
\begin{tabularx}{\linewidth}{|>{\bfseries}l|X|X|X|}\hline
\textbf{Type cellulaire} & \textbf{Lieu de formation} & \textbf{Lieu de maturation} & \textbf{Rôle principal (effecteur)} \\ \hline
Polynucléaire neutrophile & Moelle osseuse rouge & Moelle osseuse rouge & Phagocytose bactérienne (1\textsuperscript{ère} ligne, non spécifique) \\ \hline
Monocyte $\rightarrow$ Macrophage & Moelle osseuse rouge & Moelle (monocyte) $\rightarrow$ Tissus (macrophage) & Phagocytose massive, présentation d'antigène (APC), cytokines \\ \hline
Lymphocyte B (LB) & Moelle osseuse rouge & \textbf{Moelle osseuse rouge} (Homme) & Différenciation en plasmocytes $\rightarrow$ sécrétion d'anticorps (immunité \textbf{humorale}) \\ \hline
Lymphocyte T CD4$^+$ (LT auxiliaire) & Moelle osseuse rouge & \textbf{Thymus} & Orchestration : aide LB (IgG), activation macrophages, activation LT8 (cytokines : IL-2, IFN-$\gamma$) \\ \hline
Lymphocyte T CD8$^+$ (LT cytotoxique) & Moelle osseuse rouge & \textbf{Thymus} & Lyse directe cellules cibles infectées ou tumorales (perforine, granzymes, FasL) \\ \hline
\end{tabularx}
\end{center}

\begin{exoresolu}[Analyse d'un cas clinique : Syndrome de DiGeorge]
\textbf{Énoncé :} Un nourrisson de 3 mois présente une hypocalcémie tétanique, une cardiopathie congénitale (tronc artériel commun) et des infections à répétition (candidose buccale, pneumonies virales). Le bilan immunologique montre : taux d'IgG, IgA, IgM normaux ou légèrement bas ; nombre de LB normaux ; \textbf{nombre de LT CD3$^+$ très bas} ($<$ 500/\textmu L) ; absence de réponse à la tuberculine (IDR) ; absence d'image thymique à la radiographie.
\begin{enumerate}
    \item Identifier l'organe atteint à l'origine de ce tableau.
    \item Expliquer le mécanisme physiopathologique reliant l'atteinte de cet organe aux anomalies immunologiques observées.
    \item Préciser pourquoi le taux d'immunoglobulines n'est pas effondré.
\end{enumerate}
\tcblower
\textbf{Solution détaillée :}
\begin{enumerate}
    \item L'association cardiopathie $+$ hypocalcémie (hypoparathyroïdie) $+$ absence de thymus à la radiographie oriente vers le \textbf{syndrome de DiGeorge} (délétion 22q11.2). L'organe principalement atteint sur le plan immunologique est le \textbf{thymus} (aplasie/hypoplasie thymique).
    \item Le thymus est le site de maturation des LT. Son absence empêche la différenciation des pré-LT (venus de la moelle) en LT matures CD4$^+$ et CD8$^+$.
    \begin{itemize}
        \item $\rightarrow$ Pénurie de LT circulants (lymphopénie T sévère).
        \item $\rightarrow$ Pas d'immunité cellulaire : pas de rejet de greffe, pas de DHR (tuberculine négative), pas de lyse cellulaire par LT8, pas d'activation des macrophages $\rightarrow$ infections intracellulaires (virus, champignons \emph{Candida}).
        \item $\rightarrow$ Pas d'aide auxiliaire (LT CD4$^+$) pour les LB.
    \end{itemize}
    \item Les LB se forment et mûrissent dans la \textbf{moelle osseuse}, qui est intacte. Ils peuvent donc répondre aux antigènes \emph{thymoindépendants} (polysaccharides bactériens capsulaires) et produire des IgM, et dans une moindre mesure IgG/IgA sans aide T. D'où un taux d'Ig globalement conservé, mais une incapacité à produire des anticorps de haute affinité, de classe IgG, et à générer une mémoire vaccinale efficace contre les antigènes protéiques (virus, toxines).
\end{enumerate}
\textbf{Conclusion :} Ce cas illustre parfaitement la dissociation formation/maturation : moelle intacte $\rightarrow$ LB OK ; thymus absent $\rightarrow$ LT absents $\rightarrow$ immunité cellulaire effondrée $+$ immunité humorale partielle (dépendante du thymus).
\end{exoresolu}

\begin{exercice}[Entraînement : QCM Bac]
Parmi les affirmations suivantes, laquelle est \textbf{exacte} concernant l'origine et la maturation des lymphocytes chez l'Homme adulte ?
\begin{enumerate}
    \item Les lymphocytes B et T se forment et mûrissent dans le thymus.
    \item Les lymphocytes B se forment dans la moelle osseuse et mûrissent dans le thymus ; les lymphocytes T se forment et mûrissent dans la moelle osseuse.
    \item Les lymphocytes B et T se forment dans la moelle osseuse ; les lymphocytes B mûrissent dans la moelle osseuse et les lymphocytes T mûrissent dans le thymus.
    \item Les lymphocytes B se forment dans la rate et mûrissent dans les ganglions ; les lymphocytes T se forment dans le thymus et mûrissent dans la moelle osseuse.
\end{enumerate}
\textit{Réponse attendue : \textbf{3}. Justification : Formation commune en moelle (CSH). Maturation différentielle : Moelle pour LB (Homme), Thymus pour LT. La rate et les ganglions sont des organes secondaires (site de rencontre Ag/Lymphocyte naïf).}
\end{exercice}

\begin{corrige}[Exercice n°1]
La réponse 3 est la seule correcte. Rappel : La moelle osseuse rouge est l'organe hématopoïétique principal chez l'adulte (formation de toutes les lignées). La maturation des LB y a lieu (chez les mammifères). Les pré-LT migrent par le sang vers le thymus pour y subir leur maturation (sélection positive/négative). Les organes secondaires (rate, ganglions) n'interviennent qu'après maturation, pour l'activation clonale.
\end{corrige}

\coursec{La reconnaissance du soi et du non-soi}

Dans la section précédente, nous avons vu que les lymphocytes naissent dans la moelle osseuse, puis acquièrent leur maturité dans le thymus (lymphocytes T) ou dans la moelle osseuse elle-même (lymphocytes B), avant de coloniser les organes lymphoïdes secondaires (rate, ganglions, amygdales, plaques de Peyer). Mais une question fondamentale reste en suspens : une fois matures et en circulation dans l'organisme, comment ces cellules savent-elles \emph{quoi} attaquer et \emph{quoi} respecter ? Pourquoi ton système immunitaire détruit-il le virus de la grippe ou le \emph{Plasmodium falciparum} agent du paludisme, mais épargne-t-il tes propres cellules ? Et pourquoi, à l'inverse, rejette-t-il violemment un rein greffé provenant d'un donneur incompatible ? Toute la réponse tient en une capacité remarquable : la \cle{reconnaissance du soi et du non-soi}. C'est l'objet de cette section.

\courssub{Les marqueurs du soi : le Complexe Majeur d'Histocompatibilité (CMH)}

\textbf{Une carte d'identité moléculaire à la surface de chaque cellule}

Observe un instant ta main : elle est constituée de milliards de cellules. Chacune d'elles porte à sa surface, insérées dans la membrane plasmique, des molécules particulières qui jouent le rôle de « carte d'identité » de l'individu. Ces molécules appartiennent à un ensemble appelé \cle{Complexe Majeur d'Histocompatibilité}, abrégé \cle{CMH}.

\begin{definition}[Le Complexe Majeur d'Histocompatibilité (CMH ou HLA)]
Le \cle{CMH} est un ensemble de glycoprotéines membranaires présentes à la surface des cellules d'un individu et qui caractérisent cet individu. Chez l'Homme, on le nomme système \cle{HLA} (\emph{Human Leucocyte Antigen}), car il a été découvert sur les leucocytes. Les gènes codant le CMH sont situés sur le \textbf{chromosome 6}. On distingue deux classes principales :
\begin{itemize}
\item le \cle{CMH de classe I}, présent à la surface de \textbf{toutes les cellules nucléées} de l'organisme (donc absents des érythrocytes) ;
\item le \cle{CMH de classe II}, présent uniquement à la surface des \textbf{cellules présentatrices d'antigènes professionnelles} (CPA) : macrophages, cellules dendritiques, lymphocytes B.
\end{itemize}
\end{definition}

\begin{propriete}[Unicité du CMH — propriété admise]
Le CMH est \textbf{unique pour chaque individu} : deux personnes prises au hasard n'ont quasiment jamais le même ensemble de molécules HLA, en raison de l'extrême polymorphisme des gènes du CMH (plusieurs dizaines, voire centaines d'allèles possibles pour chaque gène : HLA-A, HLA-B, HLA-C pour la classe I ; HLA-DP, HLA-DQ, HLA-DR pour la classe II). La seule exception concerne les \textbf{jumeaux monozygotes} (vrais jumeaux), issus du même zygote, qui possèdent un patrimoine génétique identique et donc le même CMH. Cette unicité du CMH conditionne le rejet des greffes : tout tissu étranger portant un CMH différent est reconnu comme non-soi et attaqué.
\end{propriete}

\courssub{Mise en évidence expérimentale : les greffes de peau chez la souris}

Comment a-t-on démontré que le système immunitaire distingue le soi du non-soi ? Par des expériences classiques de greffes de peau réalisées chez la souris par Peter Medawar (prix Nobel 1960), que tu dois savoir interpréter pour le Baccalauréat.

\begin{experience}[Greffes de peau chez la souris : protocole et résultats (d'après Medawar)]
\textbf{Matériel :} deux lignées de souris génétiquement différentes (lignée A et lignée B), du matériel de chirurgie stérile.\\
\textbf{Protocole et observations :}
\begin{itemize}
\item \textbf{Expérience 1 (autogreffe) :} on prélève un fragment de peau sur le dos d'une souris A et on le greffe sur le flanc de la \emph{même} souris A. \textbf{Résultat :} la greffe est acceptée, la peau cicatrise normalement, vascularisation et repousse des poils sont observées.
\item \textbf{Expérience 2 (allogreffe) :} on prélève un fragment de peau d'une souris A et on le greffe sur une souris B (individu différent, de lignée différente). \textbf{Résultat :} la greffe est rejetée en \SIrange{10}{14}{jours} : la peau greffée noircit, nécrose et tombe (réaction de rejet aigu).
\item \textbf{Expérience 3 (greffe entre jumeaux) :} on greffe la peau d'une souris sur sa sœur jumelle monozygote. \textbf{Résultat :} la greffe est acceptée.
\item \textbf{Expérience 4 (seconde allogreffe) :} on réalise une seconde greffe de peau de la souris A sur la souris B qui a déjà rejeté une première greffe de A. \textbf{Résultat :} le rejet est plus rapide (environ \SI{6}{jours}), signe d'une mémoire immunitaire.
\end{itemize}
\textbf{Interprétation :} la souris accepte son propre tissu (soi) et celui de son jumeau génétiquement identique (même CMH), mais rejette le tissu d'un individu génétiquement différent (non-soi, CMH différent). Le rejet accéléré lors de la seconde greffe prouve l'implication de la mémoire immunitaire.
\textbf{Conclusion :} le système immunitaire distingue le \cle{soi} du \cle{non-soi} ; cette discrimination repose sur les marqueurs membranaires (CMH) propres à chaque individu.
\end{experience}

\begin{popfigure}
\begin{center}
\begin{tabularx}{\linewidth}{|l|X|X|X|}
\hline
\rowcolor{popblueL} \textbf{Type de greffe} & \textbf{Donneur} & \textbf{Receveur} & \textbf{Issue} \\
\hline
Autogreffe & Souris A & Souris A (elle-même) & \textcolor{popgreen}{\textbf{Acceptation}} : cicatrisation normale, vascularisation \\
\hline
Allogreffe (1ʳᵉ) & Souris A & Souris B (lignée différente) & \textcolor{poporange}{\textbf{Rejet}} en 10--14 j : nécrose, dépilation \\
\hline
Greffe entre jumeaux & Souris A$_1$ & Souris A$_2$ (jumelle monozygote) & \textcolor{popgreen}{\textbf{Acceptation}} \\
\hline
Allogreffe (2ᵉ) & Souris A & Souris B (déjà greffée) & \textcolor{poporange}{\textbf{Rejet accéléré}} (\textasciitilde 6 j) : mémoire immunitaire \\
\hline
\end{tabularx}
\end{center}
\legende{Tableau-schéma des résultats des greffes de peau chez la souris (expériences de Medawar). L'autogreffe et la greffe entre jumeaux monozygotes (même CMH) sont acceptées ; l'allogreffe (CMH différent) est rejetée. Le rejet est plus rapide lors d'une seconde greffe du même donneur, prouvant l'existence d'une mémoire immunitaire.}
\end{popfigure}

\begin{methode}[Conclure à partir de résultats de greffes — démarche type Bac]
Face à un tableau ou un document décrivant des greffes, adopte la démarche suivante :
\begin{enumerate}
\item \textbf{Comparer} les situations : identifier le donneur et le receveur dans chaque cas (même individu ? même lignée ? individus différents ? jumeaux ?).
\item \textbf{Dégager la variable} qui change d'une expérience à l'autre : ici, l'identité génétique (et donc le CMH) entre donneur et receveur, ou le rang de la greffe (1ʳᵉ ou 2ᵉ).
\item \textbf{Formuler une interprétation} : le rejet n'a lieu que lorsque le CMH du donneur diffère de celui du receveur ; la rapidité du rejet lors de la seconde greffe témoigne d'une mémoire.
\item \textbf{Conclure} en répondant au problème posé : le système immunitaire reconnaît le soi (toléré) et rejette le non-soi (éliminé) ; le CMH est la molécule-support de cette reconnaissance ; le rejet est une réponse immunitaire spécifique à mémoire.
\end{enumerate}
\end{methode}

\begin{savaistu}[Peter Medawar et la naissance de l'immunologie des greffes]
Le biologiste britannique \textbf{Peter Medawar} (1915--1987), travaillant pendant la Seconde Guerre mondiale sur les brûlures des pilotes de la RAF, a montré que le rejet de greffe est un phénomène \emph{immunologique} et non chirurgical. Ses expériences sur la souris (1940--1950) ont établi les bases de la compatibilité tissulaire (HLA) et de la tolérance immunologique. Il reçoit le prix Nobel de physiologie ou médecine en 1960 avec Macfarlane Burnet. Au Cameroun, ses découvertes permettent aujourd'hui les greffes rénales au Centre Hospitalier Universitaire de Yaoundé ou à l'Hôpital Laquintinie de Douala, grâce au typage HLA et aux immunosuppresseurs.
\end{savaistu}

\courssub{Le soi et le non-soi : définitions}

\begin{definition}[Soi et non-soi]
\begin{itemize}
\item Le \cle{soi} est l'ensemble des molécules (protéines, glycoprotéines, glucides...) normalement présentes dans l'organisme d'un individu et qui sont \textbf{tolérées} par son système immunitaire. Le CMH de l'individu est le marqueur fondamental du soi.
\item Le \cle{non-soi} est l'ensemble des molécules étrangères à l'individu, capables de déclencher une réponse immunitaire. Ces molécules étrangères sont appelées \cle{antigènes}.
\end{itemize}
\end{definition}

\begin{definition}[Antigène et épitope]
Un \cle{antigène} est une molécule (le plus souvent une protéine ou un glucide de haut poids moléculaire, \textgreater \SI{10000}{Da}) étrangère à l'organisme, capable d'être reconnue par le système immunitaire et de déclencher une réponse immunitaire. Un antigène n'est pas reconnu dans sa totalité : la petite région précise de l'antigène (quelques acides aminés ou sucres) effectivement reconnue par un récepteur de lymphocyte ou par un anticorps s'appelle un \cle{épitope} (ou déterminant antigénique). Un même antigène peut porter \textbf{plusieurs épitopes différents}, reconnus par des clones lymphocytaires différents.
\end{definition}

\begin{attention}[Piège classique : antigène $\neq$ épitope]
Un anticorps (ou un récepteur lymphocytaire) ne reconnaît \textbf{pas un antigène entier}, mais un \textbf{épitope précis} de cet antigène, comme une clé ne reconnaît qu'une seule serrure. Conséquence : une seule bactérie, portant de nombreux antigènes de surface à épitopes variés, peut activer simultanément plusieurs clones de lymphocytes B et T différents. Ne rédige jamais « l'anticorps reconnaît l'antigène » sans préciser qu'il se fixe sur un épitope déterminé. De même, un antigène peut être \emph{immunogène} (déclenche une réponse) sans être l'agent pathogène entier.
\end{attention}

\courssub{Les structures de reconnaissance : BCR, TCR et CMH}

Pour reconnaître le non-soi, les lymphocytes possèdent à leur surface des \cle{récepteurs spécifiques}, dont chaque lymphocyte porte un seul type (un seul « modèle »), généré aléatoirement lors de sa maturation par recombinaison V(D)J :

\begin{itemize}
\item les \textbf{lymphocytes B} portent des récepteurs \cle{BCR} (\emph{B Cell Receptor}), qui sont en réalité des anticorps membranaires (IgM et IgD) capables de reconnaître directement un épitope \textbf{conformationnel} (dépendant de la structure 3D) d'un antigène \textbf{libre} (soluble dans le sang, la lymphe, ou à la surface d'un pathogène) ;
\item les \textbf{lymphocytes T} portent des récepteurs \cle{TCR} (\emph{T Cell Receptor}), qui ne reconnaissent un épitope que s'il est \textbf{présenté} par une molécule de CMH à la surface d'une cellule : le CMH de classe I présente des peptides (épitopes \textbf{linéaires}, 8--10 acides aminés) aux lymphocytes T CD8$^+$ (cytotoxiques) ; le CMH de classe II présente des peptides (13--18 acides aminés) aux lymphocytes T CD4$^+$ (auxiliaires).
\end{itemize}

Ainsi, le CMH n'est pas seulement une carte d'identité : c'est aussi un \textbf{présentoir} qui expose les épitopes peptidiques aux lymphocytes T. La reconnaissance du non-soi résulte donc d'un double dialogue : le récepteur du lymphocyte (BCR ou TCR) d'un côté, l'épitope (libre ou associé au CMH) de l'autre.

\begin{popfigure}
\begin{center}
\begin{tikzpicture}[scale=1.0, every node/.style={font=\small}]
% Membrane cellulaire (double couche phospholipidique simplifiée)
\fill[popgoldL] (-7,0) rectangle (7,1.2);
\foreach \x in {-6.5,-5.7,...,6.5}{
  \draw[popgold,fill=popgold] (\x,1.05) circle (0.08);
  \draw[popgold,fill=popgold] (\x,0.15) circle (0.08);
  \draw[popgold,thin] (\x,0.95) -- (\x,0.25);
}
\node[popdark] at (0,-0.5) {\textbf{Membrane plasmique} (cytoplasme en dessous, milieu extracellulaire au-dessus)};

% CMH classe I (deux chaînes lourdes + b2-microglobuline + peptide)
\draw[popblue,very thick] (-5.2,0.1) -- (-5.2,1.8);
\draw[popblue,very thick] (-4.8,0.1) -- (-4.8,1.8);
\fill[popblue!50] (-5.0,2.0) circle (0.15); % groove
\fill[poporange] (-5.0,2.0) circle (0.06); % peptide
\node[popblue,align=center] at (-5.0,2.8) {\textbf{CMH classe I}\\ (toutes cellules nucléées)};
\node[poporange] at (-5.0,2.0) {\tiny pep};

% CMH classe II (deux chaînes alpha/beta + peptide)
\draw[popgreen,very thick] (-2.2,0.1) -- (-2.2,1.8);
\draw[popgreen,very thick] (-1.8,0.1) -- (-1.8,1.8);
\fill[popgreen!50] (-2.0,2.0) circle (0.15);
\fill[poppink] (-2.0,2.0) circle (0.06);
\node[popgreen!70!black,align=center] at (-2.0,2.8) {\textbf{CMH classe II}\\ (CPA : M$\phi$, DC, LB)};
\node[poppink] at (-2.0,2.0) {\tiny pep};

% BCR (IgM membranaire : 2 chaînes lourdes + 2 légères)
\draw[poporange,very thick] (1.5,0.1) -- (1.5,1.6);
\draw[poporange,very thick] (1.9,0.1) -- (1.9,1.6);
\draw[poporange,very thick] (1.3,1.6) -- (1.3,2.0);
\draw[poporange,very thick] (2.1,1.6) -- (2.1,2.0);
\draw[poporange,thick] (1.3,2.0) -- (1.7,2.0) -- (1.7,2.2) -- (1.3,2.2) -- cycle; % binding site
\node[poporange!80!black,align=center] at (1.7,2.9) {\textbf{BCR (IgM/IgD)}\\ (Lymphocyte B)};
\node[poporange] at (1.7,2.1) {\tiny Ag libre};

% TCR (heterodimere alpha/beta + CD3)
\draw[poppurple,very thick] (4.2,0.1) -- (4.2,1.7);
\draw[poppurple,very thick] (4.6,0.1) -- (4.6,1.7);
\draw[poppurple,very thick] (4.2,1.7) -- (4.0,2.1);
\draw[poppurple,very thick] (4.6,1.7) -- (4.8,2.1);
\node[poppurple,align=center] at (4.4,2.9) {\textbf{TCR $\alpha\beta$ + CD3}\\ (Lymphocyte T)};

% Antigène libre au-dessus
\draw[popink,fill=poppinkL] (-3.5,3.5) circle (0.5);
\fill[poppink] (-3.3,3.7) circle (0.07);
\fill[poppink] (-3.7,3.3) circle (0.07);
\fill[poppink] (-3.6,3.8) circle (0.07);
\node[popink, align=center] at (-3.5,4.3){\textbf{Antigène libre}\\ (multiples épitopes)};

% Flèches de reconnaissance
\draw[->,poporange,thick] (1.7,2.2) -- (-3.5,3.5) node[midway,above,sloped,font=\tiny] {Reconnaissance directe};
\draw[->,poppurple,thick] (4.4,2.1) -- (-2.0,2.15) node[midway,above,font=\tiny] {Reconnaissance CMH-restreinte};
\draw[->,popblue,thick] (-5.0,2.15) -- (-5.0,2.5);
\draw[->,popgreen,thick] (-2.0,2.15) -- (-2.0,2.5);
\end{tikzpicture}
\end{center}
\legende{Schéma des structures de reconnaissance à la surface des cellules. Le CMH de classe I (bleu) est présent sur toutes les cellules nucléées et présente des peptides intracellulaires aux LT CD8. Le CMH de classe II (vert) est présent sur les CPA et présente des peptides extracellulaires/internalisés aux LT CD4. Le BCR (orange, lymphocyte B) reconnaît directement un épitope conformationnel d'un antigène libre. Le TCR (violet, lymphocyte T) ne reconnaît un épitope linéaire que présenté par une molécule de CMH (restriction au CMH).}
\end{popfigure}

\courssub{La tolérance du soi : comment l'organisme évite de s'attaquer lui-même}

Si les lymphocytes génèrent leurs récepteurs au hasard lors de leur maturation (recombinaison V(D)J), une fraction non négligeable de ces récepteurs reconnaît forcément... nos propres molécules (auto-antigènes) ! Que deviennent ces lymphocytes potentiellement dangereux ?

\begin{propriete}[Tolérance centrale du soi — savoir admis]
Lors de leur maturation (dans le thymus pour les LT, dans la moelle osseuse pour les LB), les lymphocytes sont soumis à une \cle{sélection négative} (ou éducation centrale) : tout lymphocyte dont le récepteur reconnaît avec une forte affinité une molécule du soi (auto-antigène présenté par le CMH) est \textbf{éliminé} par mort cellulaire programmée (apoptose). Seuls survivent et migrent vers les organes lymphoïdes secondaires les lymphocytes qui \textbf{ne réagissent pas (ou faiblement) contre le soi}. C'est le fondement de la \cle{tolérance centrale du soi} : l'organisme ne déclenche normalement aucune réponse immunitaire contre ses propres constituants.
\end{propriete}

\begin{attention}[Vocabulaire précis : « Sélection clonale » vs « Sélection négative »]
Le terme \textbf{sélection clonale} désigne classiquement l'amplification d'un clone lymphocytaire \emph{après} rencontre de son antigène spécifique (réponse immune). Pour l'élimination des auto-réactifs \emph{pendant} la maturation, on parle de \textbf{sélection négative} (ou tolérance centrale). Ne confonds pas les deux étapes.
\end{attention}

Cette sélection explique aussi pourquoi chaque individu possède, au repos, une immense bibliothèque de lymphocytes aux récepteurs tous différents (répertoire immunitaire \textasciitilde $10^7$ à $10^8$ spécificités distinctes), prêts à reconnaître presque n'importe quel antigène étranger : c'est la base de la réponse immunitaire spécifique que nous étudierons dans les sections suivantes. Lorsque cette tolérance est rompue (défaillance de la sélection négative, mimétisme moléculaire, exposition d'antigènes cachés...), des lymphocytes activés attaquent le soi : c'est le mécanisme des \cle{maladies auto-immunes} (vitiligo, diabète de type 1 ou juvénile, myasthénie), que nous aborderons dans la section 6.

\courssub{Application médicale : compatibilité tissulaire et greffes d'organes}

\begin{exemplebox}[Pourquoi cherche-t-on un donneur compatible, souvent un frère ou une sœur, pour une greffe de rein ?]
Au Centre Hospitalier Universitaire d'Essos à Yaoundé comme dans tout service de néphrologie, un patient atteint d'insuffisance rénale terminale peut être sauvé par une greffe de rein. Mais le rein greffé porte le CMH du donneur : si ce CMH diffère trop de celui du receveur, les lymphocytes T CD8$^+$ et CD4$^+$ du receveur reconnaissent le greffon comme \textbf{non-soi} (reconnaissance directe du CMH étranger ou indirecte via peptides du donneur) et le détruisent : c'est le \cle{rejet de greffe} (aigu ou chronique). Or un frère ou une sœur a une probabilité de \SI{25}{\%} d'avoir hérité du même haplotype HLA (même combinaison d'allèles sur le chromosome 6) que le patient, puisque les enfants d'un même couple partagent leurs chromosomes 6 parentaux. On recherche donc en priorité un donneur dans la fratrie, après un \textbf{typage HLA} (analyse moléculaire des allèles HLA-A, -B, -DR sur les leucocytes). Lorsque la compatibilité est imparfaite (donneur cadavérique ou apparenté non identique), le patient reçoit à vie des \cle{immunosuppresseurs} (ciclosporine, tacrolimus, mycophénolate...), médicaments qui inhibent l'activation des lymphocytes T (blocage de l'IL-2, inhibition de la calcineurine) pour limiter le rejet — au prix, hélas, d'une plus grande vulnérabilité aux infections opportunistes et de néphrotoxicité.
\end{exemplebox}

\begin{exoresolu}[Interpréter des résultats de greffes]
\textbf{Énoncé.} On réalise chez des souris les greffes de peau suivantes : (1) peau d'une souris X greffée sur elle-même ; (2) peau d'une souris X greffée sur une souris Y de lignée différente ; (3) peau d'une souris X greffée sur une souris Y préalablement traitée par un immunosuppresseur (ex. ciclosporine). Résultats : (1) acceptation définitive ; (2) rejet en 12 jours ; (3) acceptation prolongée tant que le traitement se poursuit. Que peut-on conclure ?
\tcblower
\textbf{Solution détaillée.}
\begin{enumerate}
\item \textbf{Comparaison des situations :} Dans les cas (1) et (2), le receveur est immunocompétent ; seule l'identité génétique entre donneur et receveur varie (même individu vs individu allogénique). Dans les cas (2) et (3), le couple donneur/receveur est identique (allogreffe), mais l'activité immunitaire du receveur diffère (normal vs immunosupprimé).
\item \textbf{Interprétation :}
\begin{itemize}
\item L'acceptation de l'autogreffe (1) montre que le système immunitaire tolère les cellules portant le \emph{même} CMH (soi).
\item Le rejet de l'allogreffe (2) montre que des cellules portant un CMH \emph{différent} (non-soi) sont reconnues et détruites par le système immunitaire du receveur.
\item L'acceptation prolongée sous immunosuppresseur (3) démontre que le rejet observé en (2) est bien une réponse \emph{immunitaire} active (dépendante des lymphocytes T) : lorsque ceux-ci sont inhibés pharmacologiquement, le greffon étranger n'est plus détruit.
\end{itemize}
\item \textbf{Conclusion :} Le système immunitaire distingue le soi du non-soi grâce aux marqueurs du CMH ; le rejet de greffe est une réaction immunitaire à médiation cellulaire (lymphocytes T) dirigée contre le CMH étranger, que les immunosuppresseurs permettent de limiter tant qu'ils sont administrés.
\end{enumerate}
\end{exoresolu}

\begin{aretenir}[L'essentiel de la section]
\begin{itemize}
\item Chaque cellule nucléée porte à sa surface des marqueurs du \cle{CMH} (HLA chez l'Homme) : classe I sur toutes les cellules nucléées, classe II sur les cellules présentatrices d'antigènes professionnelles (macrophages, cellules dendritiques, lymphocytes B).
\item Le CMH est unique pour chaque individu (polymorphisme extrême, chromosome 6) : il conditionne l'acceptation ou le rejet des greffes (sauf jumeaux monozygotes).
\item Le \cle{soi} (molécules propres, tolérées) est distingué du \cle{non-soi} (antigènes étrangers) grâce aux récepteurs spécifiques des lymphocytes : \cle{BCR} (lymphocytes B, reconnaissance directe d'épitopes conformationnels sur antigène libre) et \cle{TCR} (lymphocytes T, reconnaissance d'épitopes linéaires \emph{uniquement} présentés par le CMH : classe I $\to$ LT CD8, classe II $\to$ LT CD4).
\item Un anticorps ou un récepteur reconnaît un \cle{épitope} précis, jamais l'antigène entier.
\item La \cle{tolérance centrale du soi} est acquise lors de la maturation des lymphocytes (thymus/moelle osseuse) par \textbf{sélection négative} (apoptose des clones auto-réactifs forts).
\item Conséquence médicale majeure : \textbf{typage HLA} et \textbf{immunosuppresseurs} pour les greffes d'organes ; rupture de la tolérance $\to$ maladies auto-immunes.
\end{itemize}
\end{aretenir}

Maintenant que nous savons \emph{qui} reconnaît (lymphocytes B et T), \emph{quoi} est reconnu (épitopes sur antigènes, présentés ou non par le CMH) et \emph{comment} le soi est toléré (sélection négative), il reste à comprendre \emph{comment} l'organisme riposte concrètement lors d'une première rencontre avec un pathogène. Nous commencerons, dans la section suivante, par les réactions immédiates et stéréotypées de l'immunité non spécifique (innée), avant d'étudier les réponses spécifiques adaptatives à médiation cellulaire et humorale.

\coursec{La réponse immunitaire non spécifique (rappels)}

Dans la section précédente, nous avons vu comment l'organisme \textbf{reconnaît} le non-soi grâce aux récepteurs de type \cle{TLR} (Toll-Like Receptors) sur les cellules sentinelles (macrophages, cellules dendritiques, mastocytes) qui détectent les \cle{MAP} (motifs moléculaires associés aux pathogènes). Cette reconnaissance déclenche immédiatement une riposte rapide, stéréotypée et sans mémoire : c'est la \cle{réponse immunitaire non spécifique} (ou innée). Elle constitue la première ligne de défense efficace dans les premières heures (\SIrange{0}{96}{\hour}) suivant une agression, avant que la réponse spécifique ne se mette en place. Au Cameroun, comme partout sous les tropiques, où la pression infectieuse (paludisme, infections respiratoires, diarrhées, plaies souillées) est forte, cette immunité innée est le rempart quotidien de nos populations.

\courssub{Les barrières naturelles : la première ligne de défense}

Avant même qu'un agent pathogène ne pénètre les tissus, il doit franchir des barrières physiques, chimiques et biologiques constitutives.

\begin{definition}[Réponse immunitaire non spécifique (innée)]
Ensemble des mécanismes de défense présents dès la naissance, qui s'activent immédiatement face à tout agresseur (bactérie, virus, champignon, corps étranger) \textbf{sans reconnaissance antigénique spécifique} et \textbf{sans mémoire}. Elle repose sur des barrières anatomiques, des cellules effectrices (phagocytes, NK) et des médiateurs solubles (complément, cytokines, lysozyme).
\end{definition}

\textbf{1. Barrières physiques et mécaniques.}
La \cle{peau} (épiderme kératinisé, sec, pH acide $\approx$ \num{5,5}) et les \cle{muqueuses} (épithélium jointif, cils vibratiles des voies respiratoires, péristaltisme intestinal, flux urinaire) empêchent l'adhésion et la pénétration des microbes. L'intégrité de ces barrières est cruciale : une brûlure étendue, une plaie profonde ou une sonde urinaire constituent des portes d'entrée majeures pour les infections nosocomiales dans nos hôpitaux (CHU de Yaoundé, Douala, Garoua).

\textbf{2. Barrières chimiques.}
Les sécrétions contiennent des substances antimicrobiennes :
\begin{itemize}
    \item \cle{Lysozyme} (larmes, salive, mucus) : hydrolyse la paroi bactérienne (peptidoglycane) $\rightarrow$ lyse des Gram+.
    \item \cle{Suc gastrique} (pH \numrange{1,5}{2,0}) : détruit la majorité des bactéries ingérées (sauf \emph{Helicobacter pylori} adapté).
    \item \cle{Acides gras} à chaîne courte (peau, sébum) : pH acide, effet bactériostatique.
    \item \cle{Défensines} et \cle{cathélicidines} : peptides antimicrobiens produits par épithéliums et neutrophiles (perturbent membranes microbiennes).
\end{itemize}

\textbf{3. Barrière biologique : la flore commensale (microbiote).}
La peau, le tube digestif, le vagin sont colonisés par une flore résidente (\num{e14} bactéries dans le côlon) qui empêche l'implantation de pathogènes par \cle{compétition nutritionnelle}, \cle{exclusion de niche} et \cle{production de bactériocines}. Une antibiothérapie large spectre mal conduite (fréquente en automédication à Yaoundé ou Douala) détruit cette flore et favorise les surinfections (\emph{Candida albicans}, \emph{Clostridioides difficile}).

\begin{exemplebox}[Contexte camerounais : la diarrhée de l'enfant]
Un enfant de \SI{18}{\month} admis au Centre Mère-Enfant de la Fondation Chantal Biya pour diarrhée aqueuse aiguë. L'examen coproscopique montre une flore fécale appauvrie en bactéries lactiques (bifidobactéries) après un traitement antibiotique empirique non justifié. La barrière biologique est effondrée : une souche de \emph{Salmonella} Typhimurium, habituellement tenue en échec par le microbiote, prolifère et envahit la muqueuse. La réhydratation orale (SRO) et l'apport de probiotiques (yaourt local, sachets de \emph{Saccharomyces boulardii}) visent à restaurer cette barrière.
\end{exemplebox}

\courssub{La phagocytose : le « nettoyage » cellulaire}

Lorsque les barrières sont franchies, les \cle{phagocytes professionnels} — \cle{polynucléaires neutrophiles (PNN)} (premiers arrivés, vie courte \SIrange{1}{2}{\day}) et \cle{macrophages} (issus des monocytes, résidents tissulaires, vie longue, grande capacité digestive) — engloutissent et détruisent l'agresseur.

\begin{definition}[Phagocytose]
Processus cellulaire actif par lequel un phagocyte (PNN, macrophage) adhère à une particule (bactérie, débris cellulaire), l'intériorise dans une vésicule (\cle{phagosome}), la fusionne avec des lysosomes pour former un \cle{phagolysosome} où des enzymes hydrolytiques et des espèces réactives de l'oxygène (ERO) digèrent la particule, puis rejette les résidus indigestes par exocytose.
\end{definition}

\begin{experience}[Mise en évidence des étapes de la phagocytose (modèle de Metchnikoff / observation moderne)]
\textbf{Observation :} Au microscope optique (objectif $\times$40 ou $\times$100), on met en contact des macrophages péritonéaux de souris (ou lignée RAW 264.7) avec des levures \emph{Saccharomyces cerevisiae} (ou billes de latex opsonisées) marquées au vert, et un colorant lysosomial rouge (LysoTracker).
\textbf{Protocole :} Temps 0 : mélange. Incubation à \SI{37}{\celsius}. Observation en microscopie vidéo (time-lapse) toutes les \SI{30}{\second} pendant \SI{30}{\minute}.
\textbf{Résultats attendus (schématisés ci-dessous) :}
\begin{enumerate}
    \item \textbf{Adhérence (\SIrange{0}{2}{\minute}) :} Contact membrane-leurre via récepteurs (scavengers, mannose, Fc$\gamma$R si opsonisation). Étalement de pseudopodes.
    \item \textbf{Endocytose / Formation du phagosome (\SIrange{2}{10}{\minute}) :} Fermeture des pseudopodes $\rightarrow$ vésicule intracellulaire (phagosome vert) distincte du cytoplasme.
    \item \textbf{Maturation / Digestion (\SIrange{10}{30}{\minute}) :} Fusion phagosome-lysosomes (apparition signal rouge superposé au vert = phagolysosome jaune/orange). Acidification (pH \numrange{4,5}{5,0}), activation hydrolases, \cle{éclat respiratoire} (production $\ce{O2^{.-}}$, $\ce{H2O2}$, $\ce{HOCl}$ via NADPH oxydase / MPO). Le signal vert (levure) s'estompe (digestion).
    \item \textbf{Exocytose / Rejet ($> \SI{30}{\minute}$) :} Rejet des résidus indigestes (hème, paroi bactérienne modifiée) hors de la cellule.
\end{enumerate}
\textbf{Interprétation :} La phagocytose est un processus séquentiel, énergétiquement coûteux (ATP pour cytosquelette, éclat respiratoire), aboutissant à la destruction intracellulaire de l'agresseur.
\end{experience}

\begin{popfigure}
\begin{tikzpicture}[scale=0.85, every node/.style={font=\small}]
  % Styles
  \tikzset{
    cell/.style={draw=popdark, thick, fill=popcream, rounded corners=3mm, minimum width=5.5cm, minimum height=3.5cm},
    membrane/.style={draw=popblue, ultra thick},
    cyto/.style={fill=popblueL!50},
    nucleus/.style={draw=poppurple, thick, fill=poppurpleL, ellipse},
    bacteria/.style={draw=poporange, thick, fill=poporangeL, ellipse, minimum width=0.6cm, minimum height=0.35cm},
    lyso/.style={draw=popgreen, thick, fill=popgreenL, rounded corners=1mm},
    phago/.style={draw=poppink, thick, fill=poppinkL, rounded corners=1mm},
    arrow/.style={->, >=Stealth, thick, popdark},
    label/.style={font=\footnotesize, text=popdark, align=center},
    stepbox/.style={font=\bfseries, popblue, anchor=north}
  }

  % Coordinates for 4 panels
  \foreach \i/\xbase in {1/0, 2/7.5, 3/15, 4/22.5} {
    % Cell outline
    \node[cell] (C\i) at (\xbase+2.75, 0) {};
    % Nucleus
    \node[nucleus, minimum width=1.2cm, minimum height=0.8cm] (N\i) at (\xbase+2.75, 1.2) {};
    \node[label] at (N\i) {Noyau};
    % Step title
    \node[stepbox] at (\xbase+2.75, 2.1) {Étape \i};
  }

  % Panel 1: Adhérence
  \node[bacteria, rotate=20] (B1) at (0.8, -0.2) {};
  \draw[poporange, thick, decorate, decoration={snake, amplitude=0.5pt, segment length=2pt}] (B1) -- ++(0.3,0.3) node[label, anchor=west] {Bactérie};
  \draw[arrow] (1.5, -0.2) -- (2.0, -0.2) node[label, midway, above] {Récepteurs} node[label, midway, below] {de surface};
  \draw[membrane] (2.0, -0.6) .. controls (2.3, -0.2) and (2.3, 0.2) .. (2.0, 0.6);
  \draw[popblue, thick, ->] (2.0, 0.0) -- (2.5, 0.0) node[label, right] {Pseudopodes};
  \node[label, align=center] at (2.75, -1.5){Adhérence \&\\ reconnaissance};

  % Panel 2: Endocytose
  \node[bacteria, rotate=-10] (B2) at (7.5+2.75, 0.0) {};
  \draw[phago] (7.5+2.0, -0.5) .. controls (7.5+2.3, -0.1) and (7.5+2.3, 0.1) .. (7.5+2.0, 0.5) .. controls (7.5+1.7, 0.6) and (7.5+1.3, 0.6) .. (7.5+1.0, 0.0) .. controls (7.5+1.3, -0.6) and (7.5+1.7, -0.6) .. (7.5+2.0, -0.5);
  \node[label, poppink] at (7.5+2.75, 0.0) {Phagosome};
  \draw[membrane] (7.5+1.0, -1.0) .. controls (7.5+1.5, -0.8) and (7.5+2.5, -0.8) .. (7.5+3.0, -1.0) .. controls (7.5+3.5, -0.5) and (7.5+4.5, -0.5) .. (7.5+5.0, -1.0);
  \node[label, align=center] at (7.5+2.75, -1.5){Endocytose :\\ formation du phagosome};

  % Panel 3: Digestion
  \node[lyso] (L3a) at (15+2.0, 0.8) {};
  \node[lyso] (L3b) at (15+2.0, -0.8) {};
  \node[phago] (P3) at (15+2.75, 0.0) {};
  \node[label, popgreen] at (15+1.3, 0.8) {Lysosomes};
  \node[label, poppink] at (15+2.75, 0.0) {Phagolysosome};
  \draw[popgreen, thick, ->] (L3a) -- (P3);
  \draw[popgreen, thick, ->] (L3b) -- (P3);
  \node[bacteria, fill=popmuted!50, draw=popmuted, opacity=0.5, rotate=30] at (15+2.75, 0.0) {};
  \node[label, align=center] at (15+2.75, -1.5){Fusion \& digestion\\(enzymes, ERO, pH acide)};

  % Panel 4: Exocytose
  \draw[membrane] (22.5+1.0, -1.0) .. controls (22.5+1.5, -0.5) and (22.5+2.5, -0.5) .. (22.5+3.0, -1.0) .. controls (22.5+3.5, -0.5) and (22.5+4.5, -0.5) .. (22.5+5.0, -1.0);
  \draw[phago] (22.5+3.5, -0.3) .. controls (22.5+3.8, 0.0) and (22.5+3.8, 0.0) .. (22.5+3.5, 0.3) .. controls (22.5+3.2, 0.4) and (22.5+2.8, 0.4) .. (22.5+2.5, 0.0) .. controls (22.5+2.8, -0.4) and (22.5+3.2, -0.4) .. (22.5+3.5, -0.3);
  \node[label, poppink] at (22.5+3.0, 0.0) {Résidus};
  \draw[arrow, poporange] (22.5+4.5, 0.0) -- (22.5+5.5, 0.0) node[label, right] {Rejet};
  \node[label, align=center] at (22.5+2.75, -1.5){Exocytose\\ des déchets};
\end{tikzpicture}
\legende{Les quatre étapes séquentielles de la phagocytose par un macrophage : (1) adhérence de la bactérie aux récepteurs membranaires et émission de pseudopodes ; (2) endocytose et formation du phagosome ; (3) fusion avec les lysosomes $\rightarrow$ phagolysosome, acidification, éclat respiratoire et digestion ; (4) exocytose des résidus indigestes.}
\end{popfigure}

\begin{methode}[Reconnaître les étapes de la phagocytose sur un document microscopique]
\begin{enumerate}
    \item \textbf{Identifier la cellule :} grand noyau souvent lobé (PNN) ou réniforme/excentré (macrophage), cytoplasme abondant.
    \item \textbf{Repérer les vésicules :} cercles clairs (phagosomes) ou colorés (phagolysosomes si coloration lysosomale).
    \item \textbf{Ordonner chronologiquement :} bactérie extracellulaire collée $\rightarrow$ bactérie dans vésicule claire $\rightarrow$ vésicule foncée/colorée (digestion) $\rightarrow$ vésicule résiduelle rejetée.
    \item \textbf{Noter l'opsonisation :} si présence d'anticorps/complément (non visible en microscopie ordinaire, mais inféré si phagocytose très efficace).
\end{enumerate}
\end{methode}

\begin{attention}[Pièges fréquents]
\begin{itemize}
    \item Confondre \textbf{phagocytose} (particules solides $> \SI{0,5}{\micro\meter}$ : bactéries, levures, débris) et \textbf{pinocytose} (fluide, macromolécules, vésicules $< \SI{0,1}{\micro\meter}$).
    \item Oublier que la phagocytose est \textbf{active} (consomme ATP, réarrangement actinomyosine) et non une simple diffusion.
    \item Croire que tous les leucocytes phagocytent : les \textbf{lymphocytes B et T NE phagocytent PAS} (sauf rares sous-populations innées-like).
    \item Négliger l'\textbf{opsonisation} (recouvrement par IgG, C3b) qui augmente l'efficacité 100 à 1000 fois via récepteurs Fc$\gamma$R et CR1.
\end{itemize}
\end{attention}

\begin{exoresolu}[Analyse d'une courbe cinétique de phagocytose]
\textbf{Énoncé :} On incubé des macrophages avec des billes de latex opsonisées. On compte le \% de macrophages ayant au moins une bille à différents temps. Résultats :
\begin{center}
\begin{tabular}{|c|c|c|c|c|c|c|}\hline
Temps (min) & 0 & 5 & 15 & 30 & 60 & 120 \\ \hline
\% Macrophages phagocytants & 0 & 15 & 55 & 85 & 95 & 95 \\ \hline
\end{tabular}
\end{center}
\textbf{1. Tracer la courbe.} \textbf{2. Interpréter la forme.} \textbf{3. Calculer la vitesse initiale (0--15 min).}

\tcblower
\textbf{Solution :}
\textbf{1.} Courbe sigmoïde (en S) : phase de latence (0--5 min), phase exponentielle (5--30 min), plateau de saturation (60 min).
\textbf{2.} \textit{Latence :} temps de reconnaissance, adhérence, émission pseudopodes. \textit{Exponentielle :} phagocytose active, toutes les cellules recrutées. \textit{Plateau :} saturation des capacités phagocytaires (nombre limité de phagosomes/cellule) ou épuisement des billes accessibles.
\textbf{3.} Vitesse initiale $\approx \frac{55-0}{15-0} = 3,67\%/\text{min}$ (ou en valeur absolue si nombre total cellules connu).
\end{exoresolu}

\courssub{La réaction inflammatoire : l'appel des renforts}

La phagocytose locale peut être débordée. La \cle{réaction inflammatoire} est la réponse vasculaire et cellulaire qui recrute massivement les effecteurs du sang vers le foyer infectieux.

\begin{definition}[Réaction inflammatoire]
Réponse stéréotypée, non spécifique, déclenchée par une agression (infection, traumatisme, brûlure), associant quatre signes cardinaux (\emph{rubor, calor, tumor, dolor}) dus à la vasodilatation, l'hyperperméabilité vasculaire et l'infiltration leucocytaire, orchestrée par des médiateurs chimiques (histamine, prostaglandines, cytokines, complément).
\end{definition}

\textbf{Mécanismes moléculaires et cellulaires (démarche d'investigation).}

\begin{experience}[Mise en évidence du rôle de l'histamine et de la diapédèse (modèle de la poche cutanée / fenêtre de Miles)]
\textbf{Observation :} Injection intradermique d'histamine (\SI{10}{\micro\gram}) vs sérum physiologique (témoin) sur l'avant-bras d'un volontaire sain.
\textbf{Résultats (à \SIrange{5}{10}{\minute}) :}
\begin{itemize}
    \item \textbf{Site histamine :} \cle{Papule} (gonflement blanc central = œdème par extravasation plasma), \cle{Érythème} (zone rouge périphérique = vasodilatation artériolaire), \cle{Chaleur} locale, \cle{Douleur} / prurit.
    \item \textbf{Témoin :} Aucune réaction.
\end{itemize}
\textbf{Interprétation :} L'histamine (libérée par \cle{mastocytes} et \cle{basophiles} tissulaires suite à activation TLR, complément C5a, traumatisme) agit sur récepteurs H1 des cellules endothéliales $\rightarrow$ contraction des jonctions intercellulaires $\rightarrow$ \textbf{hyperperméabilité} (sortie plasma $\rightarrow$ œdème/tumor) et sur muscle lisse vasculaire $\rightarrow$ \textbf{vasodilatation} (rubor, calor). La douleur (dolor) est due à la compression des terminaisons nerveuses par l'œdème et à l'action directe de médiateurs (bradykinine, prostaglandines $\ce{PGE2}$).
\end{experience}

\begin{popfigure}
\begin{tikzpicture}[scale=0.9, every node/.style={font=\small}]
  % Styles
  \tikzset{
    vessel/.style={draw=popdark, thick, fill=popred!20, rounded corners=2mm},
    endo/.style={draw=popblue, ultra thick},
    basal/.style={draw=popmuted, thick, dashed},
    tissue/.style={fill=popcream!50},
    mast/.style={draw=poppurple, thick, fill=poppurpleL, circle, minimum size=0.5cm},
    pnn/.style={draw=poporange, thick, fill=poporangeL, circle, minimum size=0.4cm},
    mono/.style={draw=popgreen, thick, fill=popgreenL, ellipse, minimum width=0.6cm, minimum height=0.35cm},
    hist/.style={draw=poppink, very thick, -{Stealth[length=3mm]}, poppink},
    cyto/.style={draw=popblue, thick, -{Stealth[length=2mm]}, popblue},
    label/.style={font=\footnotesize, text=popdark, align=center},
    title/.style={font=\bfseries\small, popdark}
  }

  % Normal vessel (left)
  \begin{scope}[xshift=0cm]
    \node[vessel, minimum width=2.5cm, minimum height=1.8cm] (V1) at (1.25, 0) {};
    \draw[endo] (0, 0.3) -- (2.5, 0.3) node[label, midway, above] {Endothélium jointif};
    \draw[basal] (0, -0.3) -- (2.5, -0.3) node[label, midway, below] {Membrane basale};
    \foreach \y in {0.6, 0.9, 1.2} \draw[popred, thick] (0.3, \y) -- (2.2, \y); % Blood flow
    \node[pnn] at (0.6, 0.6) {}; \node[pnn] at (1.2, 0.9) {}; \node[pnn] at (1.8, 0.6) {};
    \node[mast] (M1) at (1.25, -1.2) {};
    \node[label, align=center] at (M1){Mastocyte\\ (granules intacts)};
    \node[title] at (1.25, 1.5) {Vaisseau sain (repos)};
  \end{scope}

  % Inflamed vessel (right)
  \begin{scope}[xshift=7cm]
    \node[vessel, minimum width=3.0cm, minimum height=2.5cm] (V2) at (1.5, 0) {};
    \draw[endo] (0, 0.5) -- (3.0, 0.5);
    \draw[basal] (0, -0.5) -- (3.0, -0.5);
    % Dilated vessel: wider
    \foreach \y in {0.8, 1.2, 1.6, 2.0} \draw[popred, thick] (0.2, \y) -- (2.8, \y);
    % Leukocytes rolling/adhering
    \node[pnn] at (0.5, 0.8) {}; \node[pnn] at (1.0, 1.2) {}; \node[pnn] at (2.0, 1.6) {};
    \node[mono, rotate=90] at (2.5, 1.2) {};
    % Diapedesis: PNN squeezing between endothelial cells
    \node[pnn] (Pdiap) at (1.5, 0.5) {};
    \draw[poporange, thick, decorate, decoration={zigzag, amplitude=1pt, segment length=2pt}] (Pdiap) -- (1.5, -0.2);
    \node[label, poporange] at (1.5, -0.6) {Diapédèse};
    % Mastocyte degranulating
    \node[mast] (M2) at (1.5, -1.5) {};
    \foreach \angle/\len in {30/0.6, 60/0.7, 90/0.8, 120/0.7, 150/0.6}
      \draw[hist] (M2) -- ++(\angle:\len) node[label, pos=1.2, anchor=south] {Histamine};
    \node[label, align=center] at (M2){Mastocyte\\ (dégranulation)};
    % Chemotaxis arrows
    \draw[cyto] (3.2, 1.0) -- (3.8, 0.5) node[label, right, align=center]{Chimioattractants\\(C5a, LTB4, fMLP)};
    \draw[cyto] (3.2, 1.5) -- (3.8, 2.0);
    % Tissue edema
    \draw[popblue!30, fill=popblue!10] (0.2, -0.5) -- (2.8, -0.5) -- (2.8, -1.3) -- (0.2, -1.3) -- cycle;
    \node[label, popblue, align=center] at (1.5, -0.9){Œdème interstitiel\\(protéines, fluide)};
    \node[title] at (1.5, 2.3) {Foyer inflammatoire aigu};
  \end{scope}
\end{tikzpicture}
\legende{Schéma comparatif : à gauche, capillaire au repos (endothélium jointif, flux rapide, mastocyte granulé) ; à droite, foyer inflammatoire : vasodilatation (calor, rubor), hyperperméabilité (œdème/tumor), dégranulation mastocytaire (histamine), adhérence et diapédèse des polynucléaires neutrophiles (PNN) et monocytes vers le tissu (chimioattractisme).}
\end{popfigure}

\begin{methode}[Relier les signes cliniques de l'inflammation aux mécanismes physiopathologiques]
Sur un document (photo clinique, schéma, tableau de signes) :
\begin{enumerate}
    \item \textbf{Lister les 4 signes cardinaux :} Rougeur (Rubor), Chaleur (Calor), Gonflement (Tumor), Douleur (Dolor). (\emph{Optionnel : Functio laesa = perte de fonction}).
    \item \textbf{Associer chaque signe à sa cause vasculaire/cellulaire :}
    \begin{itemize}
        \item Rubor + Calor $\leftarrow$ Vasodilatation artériolaire (afflux sanguin).
        \item Tumor $\leftarrow$ Hyperperméabilité veineuse $\rightarrow$ extravasation plasma (œdème) + infiltration leucocytaire.
        \item Dolor $\leftarrow$ Compression nerfs par œdème + médiateurs algogènes (bradykinine, $\ce{PGE2}$).
    \end{itemize}
    \item \textbf{Identifier les cellules et médiateurs clés :} Mastocytes (histamine), Endothélium (jonctions), PNN (premiers infiltrés, \SIrange{2}{6}{\hour}), Monocytes/Macrophages (plus tard, \SIrange{24}{48}{\hour}), Chimiokines (gradient).
\end{enumerate}
\end{methode}

\begin{attention}[L'inflammation et la fièvre sont des alliés, pas des ennemis !]
\begin{itemize}
    \item \textbf{Ne pas confondre symptôme et maladie.} La fièvre ($> \SI{38,5}{\celsius}$) inhibe la réplication virale/bactérienne, stimule la prolifération lymphocytaire et l'activité des PNN. La traiter systématiquement (paracétamol) peut \textbf{retarder la guérison} (études sur grippe, varicelle).
    \item \textbf{L'inflammation aiguë est bénéfique :} elle isole le foyer (fibrine), apporte effecteurs (PNN, Ac, complément) et nettoie (macrophages). La bloquer (corticoïdes, AINS puissants) trop tôt ou sans indication précise risque de \textbf{transformer une infection banale en septicémie}.
    \item \textbf{Seul un médecin évalue le rapport bénéfice/risque} (ex. : méningite purulente $\rightarrow$ dexaméthasone adjunctive OBLIGATOIRE avant antibiotiques pour limiter l'œdème cérébral ; arthrite septique $\rightarrow$ drainage chirurgical + ATB, pas seulement AINS).
\end{itemize}
\end{attention}

\begin{savaistu}[Le pus : un cimetière de héros]
Le \textbf{pus} (matière jaune-verdâtre, épaisse) qui s'écoule d'un abcès ou d'un bouton infecté (furoncle) est composé à \textbf{90--95 \% de polynucléaires neutrophiles MORTS} (noyaux lysés, ADN extracellulaire $\rightarrow$ viscosité), de débris bactériens, de tissu nécrosé et de liquide interstitiel. La couleur verte de certaines infections à \emph{Pseudomonas aeruginosa} vient de la \cle{pyocyanine} (pigment bactérien). \textbf{Ne jamais « percer » soi-même un abcès} : risquer de disséminer l'infection en profondeur (phlegmon, septicémie). Au dispensaire, le médecin pratique une \textbf{incision-drainage} stérile + prélèvement pour antibiogramme.
\end{savaistu}

\courssub{Caractéristiques, limites et articulation avec la réponse spécifique}

\begin{propriete}[Caractères définitoires de l'immunité non spécifique]
\begin{itemize}
    \item \textbf{Rapidité :} Secondes à heures (barrières, phagocytose, inflammation) vs jours (spécifique).
    \item \textbf{Universalité / Non-spécificité :} Même réponse quel que soit l'agent (Gram+, Gram-, virus, champignon, éclat de bois). Pas de récepteurs clonaux réarrangés.
    \item \textbf{Absence de mémoire :} Pas de réponse plus rapide/forte lors d'une 2\ieme{} rencontre. Pas de lymphocytes mémoire innés (sauf phénomène récent d'\emph{trained immunity} / mémoire innée via reprogrammation épigénétique des monocytes, hors programme Bac).
    \item \textbf{Héréditaire :} Génome non modifié, identiques chez tous les individus d'une espèce.
\end{itemize}
\end{propriete}

\begin{aretenir}[Limites de l'immunité innée et relais vers l'adaptative]
\begin{enumerate}
    \item \textbf{Dépassement quantitatif :} Nombre de pathogènes $>$ capacité phagocytaire (ex. : septicémie à \num{e8} UFC/mL).
    \item \textbf{Évasion microbienne :} Capsules anti-phagocytaires (\emph{Streptococcus pneumoniae}, \emph{Klebsiella}), toxines lysant les PNN (\emph{Staph. aureus} : leucocidine), survie intracellulaire (\emph{Mycobacterium tuberculosis}, \emph{Salmonella}, \emph{Listeria}).
    \item \textbf{Relais obligatoire :} Les \textbf{macrophages} (et cellules dendritiques) qui ont phagocyté migrent vers le ganglion drainant via lymphatiques. Ils deviennent \textbf{Cellules Présentatrices d'Antigène Professionnelles (CPAP)} : ils présentent des peptides antigéniques sur \cle{CMH de classe II} aux \cle{lymphocytes T CD4$^+$ naïfs} $\rightarrow$ déclenchement réponse spécifique (voir Section 4).
\end{enumerate}
\end{aretenir}

\begin{exemplebox}[Articulation innée / adaptative : le cas de la tuberculose]
Au Centre Jamot de Yaoundé (référence TB), un patient présente une tuberculose pulmonaire. \emph{Mycobacterium tuberculosis} (bacille acido-alcool résistant, BAAR) résiste à la digestion dans le macrophage (inhibition fusion phagosome-lysosome, paroi cireuse mycoliques). L'immunité innée \textbf{contient} mais n'élimine pas : formation de \cle{granulome} (macrophages épithélioïdes, géants multinucléés, couronne lymphocytaire). Seule la réponse spécifique (Th1 $\rightarrow$ IFN-$\gamma$ $\rightarrow$ activation macrophage « tueur » + cytolytique CD8) peut stériliser le foyer. Le test tuberculinique (IDR) et l'IGRA (dosage IFN-$\gamma$ sanguin) mesurent \textbf{cette mémoire spécifique}, pas l'innée.
\end{exemplebox}

\begin{exercice}[Entraînement Bac : Analyse d'un document sur l'inflammation]
\textbf{Document :} Coupe histologique d'un ganglion lymphatique drainant un abcès cutané à \emph{Staphylococcus aureus} (coloration H\&E). On observe : 1/ Zone corticale : follicles lymphoïdes réactifs (centres germinatifs). 2/ Zone paracorticale : forte hyperplasie, nombreux lymphoblastes, mitoses. 3/ Sinus médullaires : distendus, remplis de polynucléaires neutrophiles (PNN) et de macrophages chargées de débris. 4/ Vaisseaux à haute endothélias (HEV) : endothélium cuboïde, adhérence de lymphocytes.

\textbf{Questions :}
1. Identifier les signes histologiques de la réaction inflammatoire aigüe dans les sinus.
2. Expliquer la présence massive de PNN dans les sinus médullaires (origine, chemin, molécules d'adhérence).
3. Relier l'hyperplasie paracorticale et les centres germinatifs à la réponse immunitaire spécifique naissante.
4. Quel est le rôle des HEV dans ce contexte ?
\end{exercice}

\begin{corrige}[Exercice n°1]
1. \textbf{Signes inflammatoires aigus :} Infiltrat massif de PNN (noyaux lobés, cytoplasme clair) dans les sinus $\rightarrow$ marqueur de l'inflammation aiguë bactérienne (diapédèse, chimioattractisme IL-8, C5a, LTB4). Œdème (sinus distendus). Macrophages phagocytaires (cytoplasme abondant, débris).
2. \textbf{Origine \& chemin :} PNN $\leftarrow$ Moelle osseuse (hématopoïèse d'urgence stimulée par G-CSF) $\rightarrow$ Sang (marge $\rightarrow$ roulement sur endothélium activé via selectines E/P $\rightarrow$ activation intégrines $\beta_2$ (LFA-1/Mac-1) sur ICAM-1/2 endothélial $\rightarrow$ arrêt ferme $\rightarrow$ diapédèse $\rightarrow$ chimioattractisme vers ganglion via lymphatiques afférents ou voie sanguine).
3. \textbf{Hyperplasie paracorticale / centres germinatifs :} Zone T (paracortex) : prolifération clonale lymphocytes T CD4$^+$ activés par DCs présentatrices (migrées du foyer). Centres germinatifs (zone B) : prolifération clonale lymphocytes B, hypermutation somatique, commutation de classe $\rightarrow$ production IgG anti-staphylocoque (toxines, protéines de surface).
4. \textbf{HEV (High Endothelial Venules) :} Porte d'entrée spécifique des lymphocytes naïfs (CD62L / L-selectine $\leftrightarrow$ adressines PNAd sur HEV) $\rightarrow$ CCR7 $\leftrightarrow$ CCL19/21 $\rightarrow$ entrée dans ganglion pour rencontre antigène. Preuve du recrutement actif de la réponse adaptative.
\end{corrige}

\coursec{La réponse immunitaire spécifique : médiation cellulaire et humorale}

Après avoir vu comment l'immunité non spécifique constitue une première ligne de défense rapide mais peu discriminante, nous abordons maintenant la réponse immunitaire \textbf{spécifique} (ou adaptative). Elle se caractérise par sa capacité à reconnaître précisément un antigène donné, à amplifier une réponse ciblée et à conserver une \textbf{mémoire} permettant une réaction plus rapide et plus intense lors d'une seconde rencontre. Cette réponse se déploie selon \textbf{deux voies effectrices} qui collaborent : la médiation cellulaire et la médiation humorale.

\courssub{Deux voies de réponse spécifique}

\begin{definition}[Réponse à médiation cellulaire]
Réponse immunitaire spécifique dont les effecteurs sont des \cle{lymphocytes T cytotoxiques (LT8)} qui détruisent directement les cellules de l'organisme présentant un antigène étranger associé à leur molécule de \cle{CMH de classe I} (cellules infectées par un virus, cellules cancéreuses, greffon allogénique).
\end{definition}

\begin{definition}[Réponse à médiation humorale]
Réponse immunitaire spécifique dont les effecteurs sont des \cle{anticorps} (immunoglobulines) sécrétés par les \cle{plasmocytes}, issus de la différenciation des \cle{lymphocytes B (LB)} activés. Les anticorps agissent dans les liquides extracellulaires (sang, lymphe, liqueurs biologiques) contre les antigènes libres ou portés par des agents pathogènes extracellulaires (bactéries, toxines, virus avant pénétration cellulaire).
\end{definition}

\begin{definition}[Plasmocyte]
Cellule effectrice de la réponse humorale, issue de la différenciation terminale d'un lymphocyte B activé. C'est une \textbf{usine à anticorps} : son cytoplasme est riche en réticulum endoplasmique rugueux et en appareil de Golgi, lui permettant de synthétiser et de sécréter jusqu'à \num{10000} molécules d'anticorps par seconde.
\end{definition}

\begin{definition}[Anticorps (immunoglobuline)]
Glycoprotéine en Y, composée de deux chaînes lourdes et deux chaînes légères, dont la région variable (extrémités des bras du Y) assure la reconnaissance spécifique d'un \cle{épitope} antigénique, et la région constante (pied du Y) détermine la classe (IgG, IgM, IgA, IgE, IgD) et les fonctions effectrices (fixation du complément, liaison aux récepteurs Fc des phagocytes, passage placentaire).
\end{definition}

\courssub{Mise en évidence expérimentale de la médiation cellulaire : le rejet de greffe}

L'expérience classique qui a démontré le rôle des \textbf{cellules} (et non des anticorps) dans le rejet de greffe est due à \textbf{Mitchison (1954)} et confirmée par de nombreux travaux ultérieurs.

\begin{experience}[Transfert du rejet de greffe : cellules vs sérum]
\textbf{Protocole :}
\begin{enumerate}
    \item Des souris de souche A reçoivent une greffe de peau de souche B (allogreffe). Elles rejettent la greffe en \SIrange{10}{12}{\day} (rejet de premier contact).
    \item On prélève des \textbf{lymphocytes} (ganglions, rate) et du \textbf{sérum} (sang décoagulé, sans cellules) chez ces souris « rejetées » (immunisées).
    \item On transfère \textbf{séparément} :
    \begin{itemize}
        \item des lymphocytes de souris rejetées $\to$ souris naïves (non immunisées) de souche A (récepteurs) qui reçoivent une greffe de souche B ;
        \item du sérum de souris rejetées $\to$ autres souris naïves de souche A recevant une greffe de souche B.
    \end{itemize}
    \item On surveille le sort des greffes chez les récepteurs.
\end{enumerate}
\textbf{Observations :}
\begin{itemize}
    \item Les souris ayant reçu des \textbf{lymphocytes} rejettent la greffe en \SIrange{5}{6}{\day} (rejet accéléré, signe de mémoire).
    \item Les souris ayant reçu du \textbf{sérum} rejettent la greffe en \SIrange{10}{12}{\day} (délai normal de premier contact, pas d'accélération).
\end{itemize}
\textbf{Interprétation :} Seuls les lymphocytes (cellules) transfèrent la capacité de rejet accéléré. Le sérum, qui contient les anticorps, ne transfère pas cette immunité. La réponse de rejet de greffe est donc une \textbf{réponse à médiation cellulaire}.
\legende{Expérience de transfert adoptif démontrant que l'immunité de rejet de greffe est transférée par les cellules (lymphocytes T) et non par le sérum (anticorps).}
\end{experience}

\begin{attention}[Piège fréquent]
Ne pas confondre « réponse cellulaire » (effecteurs = cellules T cytotoxiques) et « immunité cellulaire » (terme ancien parfois utilisé pour désigner toute réponse impliquant des cellules). Au Bac, on parle de \textbf{réponse à médiation cellulaire} pour la voie LT8, et de \textbf{réponse à médiation humorale} pour la voie anticorps/LB/plasmocytes.
\end{attention}

\courssub{Mécanisme de la médiation cellulaire : destruction des cellules cibles}

La réponse à médiation cellulaire vise à éliminer les \textbf{cellules altérées du soi} (infectées, tumorales, greffées). Elle repose sur la reconnaissance d'un complexe \textbf{antigène peptidique + CMH de classe I} par le récepteur du lymphocyte T cytotoxique (LT8).

\begin{popfigure}
\begin{tikzpicture}[
    scale=0.9,
    every node/.style={font=\small},
    cell/.style={draw=popdark, thick, fill=popcream, rounded corners=0.3cm, minimum width=2.5cm, minimum height=1.8cm},
    lymp/.style={draw=popblue, thick, fill=popblueL, circle, minimum size=1.2cm},
    arrow/.style={->, thick, >=Stealth, popdark},
    mol/.style={circle, fill=poporange, minimum size=0.35cm, inner sep=0},
    perforin/.style={draw=poppink, thick, fill=poppinkL, rectangle, rounded corners=0.15cm, minimum width=0.6cm, minimum height=0.35cm}
]
% Cellule cible infectée
\node[cell] (cible) at (0,0) {};
\node[align=center] at (0,1.2) {\textbf{Cellule cible infectée}};
% Noyau
\draw[fill=popmuted!30] (0,0.2) ellipse (0.6 and 0.4);
\node at (0,0.2) {Noyau};
% CMH I + peptide viral
\foreach \angle/\r in {30/0.9, 90/0.9, 150/0.9, -30/0.9, -90/0.9, -150/0.9} {
    \node[mol] (cmh\angle) at (\angle:\r) {};
    \draw[poporange, thick] (\angle:\r) -- (\angle:1.15);
}
\node[poporange, align=center] at (0,-1.4) {CMH I + peptide viral};
% LT8
\node[lymp] (lt8) at (4,0) {};
\node[align=center] at (4,0.9) {\textbf{LT8 (CD8$^+$)}};
\node at (4,0) {TCR};
\node at (4,-0.4) {CD8};
% Flèche reconnaissance
\draw[arrow] (lt8.west) -- ++(-0.8,0) node[midway, above] {Reconnaissance} -- (cible.east);
% Synapse immunologique
\draw[poppurple, thick, dashed] (1.2,-0.3) rectangle (2.8,0.3);
\node[poppurple, align=center] at (2,0.6) {Synapse immunologique};
% Activation : LT4 aide (schématisé à côté)
\node[draw=popgreen, thick, fill=popgreenL, circle, minimum size=1cm] (lt4) at (4,2.2) {};
\node[align=center] at (4,3) {\textbf{LT4 (CD4$^+$)}};
\node at (4,2.2) {TCR};
\node at (4,1.8) {CD4};
\draw[arrow, popgreen] (lt4.south) -- (lt8.north) node[midway, right] {IL-2};
% Multiplication clonale
\node[draw=popblue, thick, fill=popblueL, rounded corners=0.2cm, minimum width=3cm, minimum height=1cm] (clone) at (7.5,2.2) {};
\node[align=center] at (7.5,2.2) {Multiplication clonale\\(expansion des LT8 spécifiques)};
\draw[arrow] (lt8.east) -- (clone.west);
% Effecteurs LT8 vers cible
\draw[arrow] (clone.south) |- (3.5,-1.5) -| (cible.south) node[pos=0.25, right] {LT8 effecteurs};
% Perforines / Granzymes
\node[perforin] (perf) at (2.5,-1.5) {};
\node[align=center] at (2.5,-2.1) {Perforines};
\node[perforin] (granz) at (1.2,-1.5) {};
\node[align=center] at (1.2,-2.1) {Granzymes};
\draw[arrow, poppink] (perf) -- (0,-1);
\draw[arrow, poppink] (granz) -- (0,-1);
% Apoptose
\node[draw=poppink, thick, fill=poppinkL, rounded corners=0.2cm, minimum width=2.5cm, minimum height=0.7cm] (apop) at (0,-2.8) {};
\node[align=center] at (0,-2.8) {Apoptose de la cellule cible};
\draw[arrow, poppink] (0,-1.7) -- (apop.north);
% Légende
\node[align=left, anchor=north west] at (-4.5,1.3) {
\textbf{Légende :}\\
1. LT8 reconnaît complexe Ag + CMH I via TCR + CD8.\\
2. LT4 (reconnaissant Ag + CMH II sur CPA) sécrète IL-2 $\to$ prolifération LT8.\\
3. LT8 effecteurs libèrent perforines + granzymes.\\
4. Apoptose (mort programmée) de la cellule infectée.
};
\end{tikzpicture}
\legende{Mécanisme de la réponse à médiation cellulaire : reconnaissance de la cellule infectée (CMH I + peptide viral) par le LT8, activation avec aide des LT4 (IL-2), expansion clonale, destruction de la cellule cible par perforines et granzymes (apoptose).}
\end{popfigure}

\begin{propriete}[Sélection et expansion clonale (théorie de Burnet, 1957 — savoir admis)]
La spécificité de la réponse immunitaire repose sur la \cle{sélection clonale} : parmi le répertoire préexistant de lymphocytes aux récepteurs divers (générés par recombinaison V(D)J), \textbf{seuls ceux dont le récepteur reconnaît l'antigène} sont sélectionnés, activés, et subissent une \textbf{multiplication clonale} massive (mitoses successives) pour produire une population de cellules effectrices (LT8, plasmocytes) et de cellules mémoire, toutes porteuses du même récepteur spécifique.
\end{propriete}

\begin{aretenir}[Étapes clés de la médiation cellulaire]
\begin{enumerate}
    \item \textbf{Reconnaissance :} Le TCR du LT8 (co-récepteur CD8) lie un peptide antigénique présenté par le CMH de classe I à la surface d'une cellule cible (cellule infectée, cellule tumorale, greffon).
    \item \textbf{Activation (avec aide des LT4) :} La reconnaissance seule est insuffisante. Le LT4 spécifique du même antigène (présenté par le CMH II d'une cellule présentatrice d'antigène) sécrète de l'\textbf{IL-2} (interleukine-2) qui stimule la prolifération du LT8.
    \item \textbf{Expansion clonale :} Le LT8 activé se divise rapidement $\to$ armée de LT8 effecteurs (cytotoxiques) + LT8 mémoire.
    \item \textbf{Effet effecteur :} Le LT8 effecteur forme une synapse immunologique avec la cellule cible et libère des \textbf{perforines} (forment des pores dans la membrane) et des \textbf{granzymes} (sérine-protéases qui déclenchent la cascade des caspases) $\to$ \textbf{apoptose} de la cellule cible.
\end{enumerate}
\end{aretenir}

\courssub{Mise en évidence expérimentale de la médiation humorale : la sérothérapie}

L'histoire de la médecine fournit la preuve classique du rôle des \textbf{anticorps} (facteurs solubles du sérum) dans la protection contre les toxines et bactéries extracellulaires.

\begin{experience}[Sérothérapie anti-diphtérique et anti-tétanique (von Behring \ \& Kitasato, 1890 ; prix Nobel 1901)]
\textbf{Contexte :} La diphtérie et le tétanos sont dus à des \textbf{toxines} bactériennes solubles qui diffusent dans l'organisme.
\textbf{Protocole :}
\begin{enumerate}
    \item Immunisation d'animaux (chevaux, moutons) par injections répétées de toxines atténuées (anatoxines) $\to$ production d'anticorps anti-toxines.
    \item Prélèvement du \textbf{sérum} de ces animaux immunisés (sérum antivenimeux, antidiphtérique, antitétanique).
    \item Injection de ce sérum à des animaux (ou humains) non immunisés, puis exposition à la toxine mortelle.
\end{enumerate}
\textbf{Observation :} Les animaux receveurs du sérum sont \textbf{protégés} (neutralisation de la toxine). La protection est \textbf{immédiate} mais \textbf{temporaire} (quelques semaines, durée de vie des anticorps IgG $\approx$ \SIrange{21}{28}{\day}).
\textbf{Interprétation :} Le sérum (donc les anticorps) transfère la protection. C'est une \textbf{immunité passive} (anticorps préformés, pas de mémoire induite chez le receveur). La réponse humorale est à médiation \textbf{anticorps}.
\legende{La sérothérapie historique démontre que les anticorps (contenus dans le sérum) confèrent une protection contre les toxines extracellulaires.}
\end{experience}

\begin{exemplebox}[Sérothérapie antivenimeuse au Cameroun]
Au Cameroun, les morsures de serpents venimeux (\emph{Echis ocellatus}, \emph{Naja nigricollis}, \emph{Bitis gabonica}) constituent une urgence médicale fréquente en zone rurale (régions du Nord, de l'Adamaoua, de l'Extrême-Nord). Le traitement de référence est l'injection intraveineuse de \textbf{sérum antivenimeux polyvalent} (ex. \emph{Inoserp} ou \emph{FAV-Afrique}), produit par immunisation de chevaux avec des venins de plusieurs espèces. Ce sérum contient des anticorps (principalement IgG) qui \textbf{neutralisent} les toxines du venin (neurotoxines, hémorragiques, nécrosantes). C'est une application directe de la médiation humorale : transfert d'anticorps préformés pour une protection immédiate. \textbf{Attention :} ce sérum peut provoquer des réactions allergiques (choc anaphylactique, maladie du sérum) $\to$ administration en milieu hospitalier sous surveillance.
\end{exemplebox}

\courssub{Mécanisme de la médiation humorale : des LB aux plasmocytes}

La réponse humorale cible les \textbf{antigènes libres} ou les \textbf{pathogènes extracellulaires}. Elle nécessite l'activation des lymphocytes B (LB).

\begin{popfigure}
\begin{tikzpicture}[
    scale=0.9,
    every node/.style={font=\small},
    cell/.style={draw=popdark, thick, fill=popcream, rounded corners=0.3cm, minimum width=2.2cm, minimum height=1.6cm},
    lymp/.style={draw=popblue, thick, fill=popblueL, circle, minimum size=1.1cm},
    plasm/.style={draw=poppink, thick, fill=poppinkL, ellipse, minimum width=1.5cm, minimum height=1cm},
    arrow/.style={->, thick, >=Stealth, popdark},
    ag/.style={draw=poporange, thick, fill=poporangeL, circle, minimum size=0.4cm},
    ab/.style={draw=poppurple, thick, fill=poppurpleL, minimum width=0.6cm, minimum height=0.35cm, rounded corners=0.1cm}
]
% Antigène libre
\node[ag] (ag1) at (-3,2) {};
\node[ag] (ag2) at (-2.4,2.4) {};
\node[ag] (ag3) at (-2.4,1.6) {};
\node[align=center] at (-3,3) {Antigène libre};
% LB naïf
\node[lymp] (lb) at (-3,0) {};
\node[align=center] at (-3,-0.9) {\textbf{LB naïf}};
\node at (-3,0) {BCR (IgM/IgD)};
% Internalisation et présentation
\draw[arrow] (ag1) -- (lb.north);
\node[draw=poporange, thick, fill=poporangeL, circle, minimum size=0.3cm] (pep) at (-3.5,0.5) {};
\node[poporange] at (-4,0.5) {Peptide};
\draw[arrow, poporange] (lb.north) |- (-4,1) -| (-1,1) node[midway, above] {Présentation CMH II};
% LT4
\node[draw=popgreen, thick, fill=popgreenL, circle, minimum size=1.1cm] (lt4) at (1,2) {};
\node[align=center] at (1,2.9) {\textbf{LT4 (CD4$^+$)}};
\node at (1,2) {TCR};
\node at (1,1.6) {CD4};
\draw[arrow] (-1.2,1.2) -- (lt4.west);
% Activation LB : signaux 1 et 2
\draw[arrow, popgreen] (lt4.south) |- (-3,0.5) node[midway, left, align=center]{IL-4, IL-5, IL-6\\CD40L};
% Différenciation
\node[plasm] (plasm) at (-3,-2.5) {};
\node[align=center] at (-3,-3.4) {\textbf{Plasmocyte}};
\node at (-3,-2.5) {Usine à Ac};
\draw[arrow] (lb.south) -- (plasm.north) node[midway, right] {Différenciation};
% Multiplication clonale (schématisée)
\node[draw=popblue, thick, fill=popblueL, rounded corners=0.2cm, minimum width=2.5cm, minimum height=0.8cm] (clone) at (-5.5,-2.5) {};
\node[align=center] at (-5.5,-2.5) {Expansion clonale};
\draw[arrow] (lb.south west) -- (clone.north east);
% Sécrétion anticorps
\foreach \y in {-2.8, -2.5, -2.2} {
    \node[ab] (ab\y) at (-1.8,\y) {};
}
\draw[arrow, poppurple] (plasm.east) -- (-1.5,-2.5) node[midway, above] {Sécrétion};
\node[align=center, poppurple] at (-0.5,-2.5) {Anticorps spécifiques};
% Actions des Ac
\node[draw=poppurple, thick, fill=poppurpleL, rounded corners=0.2cm, minimum width=3.5cm, minimum height=2.5cm] (actions) at (3,-1) {};
\node[align=left] at (3,0.2) {
\textbf{Modes d'action des Ac :}\\
$\bullet$ Neutralisation (blocage fixation)\\
$\bullet$ Agglutination / Précipitation\\
$\bullet$ Opsonisation (facilite phagocytose)\\
$\bullet$ Activation complément (lyse)
};
\draw[arrow, poppurple] (-0.5,-2.5) -- (actions.west);
% Légende
\node[align=left, anchor=north west] at (-6.5,3) {
\textbf{Légende :}\\
1. LB capte Ag via BCR $\to$ internalisation, présentation peptide + CMH II.\\
2. LT4 spécifique reconnaît peptide + CMH II $\to$ activation (CD40L + cytokines).\\
3. LB activé $\to$ multiplication clonale + différenciation en plasmocytes.\\
4. Plasmocytes sécrètent anticorps spécifiques.\\
5. Ac agissent dans les liquides (neutralisation, opsonisation, complément...).
};
\end{tikzpicture}
\legende{Mécanisme de la réponse à médiation humorale : activation du LB par l'antigène avec aide des LT4 (double reconnaissance), expansion clonale, différenciation en plasmocytes, sécrétion d'anticorps et leurs modes d'action.}
\end{popfigure}

\begin{aretenir}[Étapes clés de la médiation humorale]
\begin{enumerate}
    \item \textbf{Reconnaissance et capture :} Le \textbf{BCR} (récepteur du LB, IgM/IgD membranaires) lie l'antigène natif (conformationnelle). Le complexe BCR-Ag est internalisé par endocytose.
    \item \textbf{Présentation et aide des LT4 :} Le LB traite l'antigène et présente des peptides sur son \textbf{CMH de classe II}. Un \textbf{LT4} spécifique du même antigène (activé par une CPA professionnelle) reconnaît ce complexe via son TCR + CD4. Cette \textbf{double reconnaissance} (LT4 voit Ag + CMH II sur le LB) déclenche l'activation du LT4 qui exprime \textbf{CD40L} et sécrète des cytokines (\textbf{IL-4, IL-5, IL-6, IL-21}) $\to$ signaux co-stimulateurs indispensables pour le LB (réponse \textbf{T-dépendante}).
    \item \textbf{Activation, expansion clonale et différenciation :} Le LB activé prolifère (expansion clonale) et se différencie en :
    \begin{itemize}
        \item \textbf{Plasmocytes} (cellules effectrices de courte durée, \SIrange{3}{5}{\day}, usines à anticorps) ;
        \item \textbf{Lymphocytes B mémoire} (longue durée, réponse secondaire rapide).
    \end{itemize}
    \item \textbf{Action des anticorps :} Sécrétés dans le sang et la lymphe, ils agissent par plusieurs mécanismes (voir ci-dessous).
\end{enumerate}
\end{aretenir}

\courssub{Modes d'action des anticorps}

Les anticorps ne tuent pas directement ; ils \textbf{marquent} l'antigène pour faciliter son élimination par d'autres composants du système immunitaire.

\begin{propriete}[Quatre modes d'action principaux des anticorps]
\begin{enumerate}
    \item \textbf{Neutralisation :} L'anticorps masque le site actif d'une toxine ou le domaine de fixation d'un virus sur son récepteur cellulaire $\to$ empêche l'action pathogène. Exemple : anticorps anti-toxine tétanique, anticorps neutralisants anti-VIH (recherchés pour les vaccins).
    \item \textbf{Agglutination / Précipitation :} Les anticorps (notamment IgM pentamères, IgG) lient plusieurs antigènes entre eux (agglutinines) ou précipitent des antigènes solubles (précipitines) $\to$ forment des complexes immuns volumineux, facilement phagocytés.
    \item \textbf{Opsonisation :} La région constante (Fc) de l'IgG se lie aux \textbf{récepteurs Fc$\gamma$} des phagocytes (macrophages, neutrophiles) $\to$ \textbf{facilite la phagocytose} de la bactérie opsonisée. C'est un mécanisme majeur contre les bactéries encapsulées (\emph{Streptococcus pneumoniae}, \emph{Haemophilus influenzae}).
    \item \textbf{Activation de la voie classique du complément :} La fixation de C1q sur les régions Fc d'IgG ou d'IgM agrégées déclenche la cascade du complément $\to$ formation du \textbf{Complexe d'Attaque Membranaire (MAC, C5b-9)} $\to$ \textbf{lyse} de la bactérie (surtout Gram-négatif) ou de la cellule cible. Le C3b déposé agit aussi comme opsonine.
\end{enumerate}
\end{propriete}

\begin{attention}[Erreur fréquente au Bac]
\begin{itemize}
    \item \textbf{Les anticorps n'entrent PAS dans les cellules.} Ils agissent exclusivement dans les \textbf{milieux extracellulaires} (sang, lymphe, liquide interstitiel, muqueuses pour IgA). Ils ne peuvent pas atteindre un virus déjà à l'intérieur d'une cellule infectée $\to$ c'est le rôle des LT8 (médiation cellulaire).
    \item Ne pas confondre \textbf{opsonisation} (Ac + récepteur Fc sur phagocyte) et \textbf{activation du complément} (cascade enzymatique $\to$ lyse + C3b opsonine). Les deux facilitent la phagocytose mais par des récepteurs différents (Fc$\gamma$R vs CR1 pour C3b).
    \item L'IgM est le premier anticorps produit (réponse primaire), pentamère, très efficace pour activer le complément et agglutiner, mais de faible affinité. L'IgG domine la réponse secondaire (affinité mature, passe le placenta, opsonisation forte).
\end{itemize}
\end{attention}

\courssub{Coopération cellulaire : le rôle central des LT4 (CD4$^+$)}

C'est le point d'articulation \textbf{clé} entre les deux voies et la clé de compréhension de la pathogénie du VIH/Sida.

\begin{popfigure}
\begin{tikzpicture}[
    scale=0.95,
    every node/.style={font=\small},
    apc/.style={draw=popdark, thick, fill=popcream, ellipse, minimum width=2.5cm, minimum height=1.5cm},
    lymp/.style={draw=popblue, thick, fill=popblueL, circle, minimum size=1.1cm},
    lt4/.style={draw=popgreen, thick, fill=popgreenL, circle, minimum size=1.1cm},
    arrow/.style={->, thick, >=Stealth, popdark},
    cyto/.style={->, thick, >=Stealth, poppink}
]
% CPA (ex: cellule dendritique)
\node[apc] (cpa) at (0,2) {};
\node[align=center] at (0,2.9) {\textbf{Cellule présentatrice d'antigène (CPA)}};
\node at (0,2.3) {CMH II + peptide};
\node at (0,1.7) {CMH I + peptide};
\node at (0,1.1) {B7 (CD80/86)};
\node at (0,0.5) {Cytokines (IL-12)};
% LT4
\node[lt4] (lt4) at (-3.5,0) {};
\node[align=center] at (-3.5,-1) {\textbf{LT4 (CD4$^+$)}};
\node at (-3.5,0) {TCR + CD4};
\node at (-3.5,-0.4) {CD28};
% LB
\node[lymp] (lb) at (3.5,0) {};
\node[align=center] at (3.5,-1) {\textbf{LB}};
\node at (3.5,0) {BCR};
\node at (3.5,-0.4) {CMH II + peptide};
% LT8
\node[lymp] (lt8) at (0,-2.5) {};
\node[align=center] at (0,-3.5) {\textbf{LT8 (CD8$^+$)}};
\node at (0,-2.5) {TCR + CD8};
% Flèches activation LT4 par CPA
\draw[arrow] (cpa.south west) -- (lt4.north east) node[midway, left, align=center] {Signal 1 : TCR\\CMH II + peptide\\Signal 2 : CD28-B7};
\draw[arrow] (cpa.south) -- ++(0,-0.3) -| (lt4.north) node[midway, above] {IL-12 $\to$ Th1};
% LT4 activé aide LB
\draw[arrow, popgreen] (lt4.east) -- node[midway, above] {CD40L + IL-4, IL-5, IL-6, IL-21} (lb.west);
% LT4 activé aide LT8
\draw[arrow, popgreen] (lt4.south) -- node[midway, right] {IL-2} (lt8.north);
% LB présente Ag à LT4 (double reconnaissance)
\draw[arrow, poporange] (lb.north west) -- node[midway, above, sloped] {Présentation Ag + CMH II} (lt4.north east);
% Légende
\node[align=left, anchor=north west] at (-6,3) {
\textbf{Coopération LT4 :}\\
$\bullet$ LT4 activé par CPA (Signaux 1+2+cytokines).\\
$\bullet$ LT4 $\xrightarrow{\text{CD40L, cytokines}}$ aide LB $\to$ plasmocytes + mémoire.\\
$\bullet$ LT4 $\xrightarrow{\text{IL-2}}$ aide LT8 $\to$ prolifération cytotoxiques.\\
$\bullet$ \textbf{VIH détruit les LT4} $\to$ effondrement des DEUX voies.
};
\end{tikzpicture}
\legende{Schéma de la coopération cellulaire : le LT4 (CD4$^+$), activé par une CPA, apporte l'aide nécessaire (CD40L, IL-2, IL-4, etc.) à la fois au LB (voie humorale) et au LT8 (voie cellulaire). La destruction des LT4 par le VIH paralyse les deux bras de l'immunité spécifique.}
\end{popfigure}

\begin{propriete}[Rôle pivot des LT4 (Th = T helper)]
Les lymphocytes T auxiliaires (LT4, CD4$^+$) sont les \textbf{chefs d'orchestre} de l'immunité adaptative. Après leur activation par une CPA (signaux 1, 2, cytokines polarisantes), ils :
\begin{itemize}
    \item expriment \textbf{CD40L} (CD154) qui lie CD40 sur le LB et la CPA (licensing) ;
    \item sécrètent des \textbf{cytokines} différenciées selon le profil Th1/Th2/Th17/Tfh :
    \begin{itemize}
        \item \textbf{IL-2} : facteur de croissance autocrine pour LT4 et paracrine pour LT8 (prolifération) ;
        \item \textbf{IL-4, IL-5, IL-6, IL-21} (profil Th2/Tfh) : activation, commutation de classe (IgG, IgE, IgA), différenciation en plasmocytes (voie humorale) ;
        \item \textbf{IFN-$\gamma$} (profil Th1) : activation macrophagique, commutation vers IgG opsonisantes, soutien LT8.
    \end{itemize}
\end{itemize}
\textbf{Conséquence majeure :} La déplétion des LT4 par le VIH entraîne une \textbf{immunodéficience combinée} : défaut de réponse humorale (pas d'aide aux LB) ET défaut de réponse cellulaire (pas d'IL-2 pour LT8).
\end{propriete}

\courssub{Tableau comparatif : médiation cellulaire vs médiation humorale}

\begin{center}
\begin{tabularx}{\linewidth}{|l|X|X|}
\hline
\rowcolor{popblueL}
\textbf{Critère} & \textbf{Médiation cellulaire (LT8)} & \textbf{Médiation humorale (LB $\to$ Plasmocytes / Ac)} \\
\hline
\textbf{Effecteur principal} & Lymphocyte T cytotoxique (LT8, CD8$^+$) & Plasmocyte (usine) $\to$ Anticorps (IgG, IgM, IgA, IgE) \\
\hline
\textbf{Cible} & \textbf{Cellules altérées du soi} : infectées (virus, bactéries intracellulaires), tumorales, greffon allogénique & \textbf{Antigènes libres / Pathogènes extracellulaires} : bactéries, toxines, virus libres, parasites extracellulaires \\
\hline
\textbf{Reconnaissance} & TCR + CD8 lie \textbf{peptide + CMH I} (sur toute cellule nucléée) & BCR (IgM/IgD membranaire) lie \textbf{antigène natif} (conformationnel) ; puis présentation peptide + CMH II au LT4 \\
\hline
\textbf{Activation} & Nécessite aide des LT4 (IL-2) après double reconnaissance (CPA + LT4) & Nécessite aide des LT4 (CD40L, IL-4, IL-5, IL-6, IL-21) après double reconnaissance (LB présente au LT4) \\
\hline
\textbf{Mécanisme effecteur} & \textbf{Contact direct} : synapse immunologique $\to$ perforines + granzymes $\to$ \textbf{apoptose} de la cellule cible & \textbf{À distance} (liquides) : neutralisation, agglutination, opsonisation (phagocytose), activation complément (lyse) \\
\hline
\textbf{Transfert adoptif} & \textbf{Transférable par les cellules} (lymphocytes) — \emph{expérience de Mitchison} & \textbf{Transférable par le sérum} (anticorps) — \emph{sérothérapie von Behring} \\
\hline
\textbf{Immunité passive} & Non transférable durablement (cellules rejetées si allogéniques) & Transférable : sérum immune, colostrum (IgA), IgG placentaires \\
\hline
\textbf{Exemple pathologique} & Rejet de greffe, contrôle infections virales, surveillance tumorale & Protection anti-toxines (tétanos, diphtérie), bactéries encapsulées, neutralisation virale \\
\hline
\end{tabularx}
\end{center}

\begin{methode}[Attribuer une immunité à la médiation cellulaire ou humorale]
Pour déterminer si une immunité est à médiation cellulaire ou humorale, on réalise des \textbf{transferts adoptifs} chez un animal naïve (syngénique ou immunosupprimé) :
\begin{enumerate}
    \item Prélever \textbf{séparément} les \textbf{lymphocytes} (purifiés sur gradient Ficoll ou tri magnétique CD3$^+$) et le \textbf{sérum} d'un animal immunisé.
    \item Transférer \textbf{uniquement les lymphocytes} $\to$ groupe 1 ; \textbf{uniquement le sérum} $\to$ groupe 2 ; \textbf{les deux} $\to$ groupe 3 ; \textbf{témoin} (sérum/lymphocytes naïfs) $\to$ groupe 4.
    \item Mettre au contact de l'antigène / pathogène / greffon.
    \item \textbf{Interpréter :}
    \begin{itemize}
        \item Protection par les lymphocytes seulement $\to$ \textbf{médiation cellulaire}.
        \item Protection par le sérum seulement $\to$ \textbf{médiation humorale}.
        \item Protection par les deux $\to$ \textbf{coopération des deux voies} (cas fréquent).
    \end{itemize}
\end{enumerate}
\end{methode}

\begin{exoresolu}[Analyse d'une expérience de transfert adoptif]
\textbf{Énoncé :} On immunise des souris de souche BALB/c contre un virus intracellulaire. Après 14 jours, on prélève rate (lymphocytes) et sérum. On transfère :
\begin{itemize}
    \item Groupe A : lymphocytes de souris immunisées $\to$ souris naïves BALB/c.
    \item Groupe B : sérum de souris immunisées $\to$ souris naïves BALB/c.
    \item Groupe C : lymphocytes + sérum $\to$ souris naïves BALB/c.
    \item Groupe Témoin : lymphocytes + sérum de souris naïves.
\end{itemize}
Tous les groupes sont ensuite infectés par le virus. On mesure la charge virale dans la rate à J+5.
\textbf{Résultats :} Charge virale (log$_{10}$ PFU/g) : Témoin = 7,2 ; Groupe A = 3,1 ; Groupe B = 6,9 ; Groupe C = 2,8.
\textbf{Questions :}
\begin{enumerate}
    \item Quelle voie de l'immunité spécifique est principalement responsable de la protection contre ce virus ? Justifier.
    \item Pourquoi le sérum seul (Groupe B) ne protège-t-il pas efficacement ?
    \item Quel type de lymphocytes est transféré dans le Groupe A pour conférer la protection ?
\end{enumerate}

\tcblower

\textbf{Solution détaillée :}
\begin{enumerate}
    \item \textbf{Médiation cellulaire.} Le Groupe A (lymphocytes seuls) réduit la charge virale de \num{4,1} log (soit $\approx$ \num{12500} fois moins de virus) par rapport au témoin, alors que le Groupe B (sérum seul) ne réduit que de \num{0,3} log (non significatif). Le transfert de la protection par les cellules et non par le sérum signe une immunité à médiation cellulaire.
    \item Le virus est \textbf{intracellulaire}. Les anticorps (contenus dans le sérum) n'ont pas accès au virus à l'intérieur des cellules infectées. Ils ne peuvent neutraliser que les virions libres extracellulaires, ce qui est insuffisant pour contrôler une infection établie. Seuls les LT8 peuvent reconnaître et tuer les cellules infectées (complexe peptide viral + CMH I).
    \item Les \textbf{lymphocytes T cytotoxiques (LT8, CD8$^+$)} spécifiques du virus. Ce sont les effecteurs de la médiation cellulaire. Le transfert inclut aussi des LT4 mémoire qui fourniront l'IL-2 nécessaire à l'expansion des LT8 chez le receveur.
\end{enumerate}
\end{exoresolu}

\begin{savaistu}[Le VIH cible le chef d'orchestre]
Le VIH (Virus de l'Immunodéficience Humaine) infecte préférentiellement les cellules exprimant le récepteur \textbf{CD4} : les \textbf{LT4}, mais aussi les macrophages et les cellules dendritiques. En détruisant massivement les LT4 (par lyse virale, apoptose, mort par syncytia, élimination par LT8 anti-VIH), le VIH \textbf{prive l'organisme de l'aide indispensable} aux deux bras de l'immunité spécifique :
\begin{itemize}
    \item Pas d'IL-2 $\to$ pas d'expansion clonale des LT8 $\to$ infections virales opportunistes, cancers (lymphomes, sarcome de Kaposi).
    \item Pas de CD40L / IL-4 / IL-21 $\to$ pas d'activation des LB, pas de commutation de classe, pas de plasmocytes $\to$ hypogammaglobulinémie paradoxale (malgré hypergammaglobulinémie polyclonale non spécifique), défaut de réponse aux nouveaux antigènes (vaccins inefficaces), infections bactériennes récidivantes.
\end{itemize}
C'est pourquoi le Sida (stade terminal) se définit par une \textbf{immunodéficience combinée sévère} (cellulaire + humorale). La thérapie antirétrovirale (ARV) vise à bloquer la réplication virale pour permettre la reconstitution des LT4 (réparation immunitaire).
\end{savaistu}

\begin{exercice}[Entraînement Bac — Médiation cellulaire vs humorale]
Un patient présente une infection à \emph{Mycobacterium tuberculosis} (bacille de Koch, bactérie intracellulaire facultative, se multiplie dans les macrophages). On isole des lymphocytes de ce patient et on les met en culture avec des macrophages infectés par le BCG (vaccin atténué). On observe une activation des lymphocytes et une destruction des macrophages infectés.
\begin{enumerate}
    \item Quelle voie de l'immunité spécifique est mise en jeu ici ? Justifier par la nature de l'agent pathogène et le mécanisme effecteur observé.
    \item Préciser les molécules de restriction (CMH) et les co-récepteurs impliqués dans la reconnaissance.
    \item Expliquer pourquoi l'injection de sérum d'un patient guéri ne protégerait pas un sujet naïf contre la tuberculose.
\end{enumerate}
\end{exercice}

\begin{corrige}[Exercice n°1]
\begin{enumerate}
    \item \textbf{Médiation cellulaire.} \emph{M. tuberculosis} est une bactérie \textbf{intracellulaire} (vit dans les macrophages). La destruction observée des macrophages infectés par les lymphocytes du patient est le signe direct de l'action des \textbf{LT8 cytotoxiques} qui tuent la cellule hôte infectée.
    \item Reconnaissance : \textbf{TCR du LT8 + CD8} lie le \textbf{peptide antigénique (origine bactérienne) présenté par le CMH de classe I} à la surface du macrophage infecté. Le macrophage agit ici comme cellule cible (et comme CPA pour l'activation initiale via CMH II vers LT4).
    \item Les anticorps (sérum) agissent dans le milieu extracellulaire. Ils ne peuvent pas pénétrer dans le macrophage pour atteindre la bactérie, ni tuer la cellule infectée. La protection contre les bactéries intracellulaires repose sur l'activation des macrophages (par IFN-$\gamma$ des LT4 Th1) et la destruction des cellules infectées par les LT8. La sérothérapie est inefficace contre la tuberculose.
\end{enumerate}
\end{corrige}

\coursec{La mémoire immunitaire et la vaccination}

Après avoir analysé les mécanismes effecteurs de la réponse immunitaire spécifique — médiation cellulaire par les lymphocytes T cytotoxiques (LTc) et médiation humorale par les anticorps produits par les plasmocytes — nous comprenons comment l'organisme élimine un agent pathogène lors d'une primo-infection. Mais une question fondamentale demeure : pourquoi une seconde rencontre avec le même agent est-elle le plus souvent silencieuse, sans symptômes cliniques ? C'est le mystère de la \cle{mémoire immunitaire}, que la vaccination exploite brillamment pour protéger les populations, y compris au Cameroun où le Programme Élargi de Vaccination (PEV) sauve chaque année des milliers de vies.

\courssub{L'observation fondatrice : une protection plus rapide et plus forte}

\emph{Observation clinique.} Un enfant qui a contracté la rougeole guérit généralement et ne développera plus jamais la maladie, même s'il est à nouveau exposé au virus des années plus tard. De même, dans les zones d'endémie palustre du Sud-Cameroun, les adultes présentent une \cle{immunité prémunitionnelle} : ils portent des plasmodies (\emph{Plasmodium falciparum}) à bas bruit, sans fièvre ni symptômes graves, alors que l'enfant non immunisé fait un paludisme grave. Cette protection relative, non stérilisante, repose sur une mémoire immunitaire entretenue par des réinfections fréquentes.

\begin{experience}[Mise en évidence de la mémoire immunitaire (expérience classique sur modèle animal)]
\textbf{Matériel :} Lapins, antigène (ex. albumine sérique bovine - BSA), adjuvant (alum), seringues, kits dosage anticorps anti-BSA (technique ELISA).
\textbf{Protocole :}
\begin{enumerate}
    \item \textbf{Primo-injection (J0) :} Injection sous-cutanée d'antigène + adjuvant chez un groupe de lapins.
    \item \textbf{Suivi primaire :} Prélèvements sanguins réguliers (J7, J14, J21, J30, J60) pour doser les anticorps anti-BSA (IgM puis IgG).
    \item \textbf{Second injection (J60) :} Injection de la même dose d'antigène (sans adjuvant nécessaire) aux mêmes lapins.
    \item \textbf{Suivi secondaire :} Prélèvements aux J62, J65, J70, J80, J100.
    \item \textbf{Témoin :} Groupe de lapins naïfs recevant une seule injection au J60 (contrôle de la réponse primaire chez un sujet non pré-immunisé).
\end{enumerate}
\textbf{Observations :} La courbe de titration des anticorps (Figure 1) montre deux phases distinctes.
\textbf{Interprétation :} La première injection déclenche une \cle{réponse primaire} lente, de faible amplitude, dominée par les IgM puis les IgG de faible affinité. La deuxième injection, plusieurs semaines plus tard, déclenche une \cle{réponse secondaire} (ou réponse d'anamnèse) plus rapide, plus intense (taux d'IgG 10 à 100 fois supérieurs), d'affinité plus élevée et plus durable. L'organisme a « gardé le souvenir » du premier antigène.
\textbf{Conclusion :} L'exposition à un antigène induit la formation de cellules à longue durée de vie (lymphocytes mémoires) capables de répondre plus efficacement lors d'une réexposition.
\end{experience}

\begin{popfigure}
\begin{tikzpicture}
\begin{axis}[
    width=\linewidth,
    height=9cm,
    xlabel={Temps (jours)},
    ylabel={Taux d'anticorps (UA/mL)},
    xmin=0, xmax=120,
    ymin=0, ymax=1100,
    xtick={0,14,30,60,70,90,120},
    xticklabels={0, 14, 30, 60 (2\textsuperscript{e} inj.), 70, 90, 120},
    ytick={0,200,400,600,800,1000},
    grid=major,
    grid style={popline},
    axis line style={popdark},
    tick style={popdark},
    legend pos=north west,
    legend cell align=left,
    clip=false
]
% Primary response curve
\addplot[popcurve, popblue, very thick, domain=0:60, samples=200] 
    ({x}, {800/(1+exp(-(x-18)/4)) * exp(-(x-18)/60) + 20*exp(-x/10)});
\addlegendentry{Réponse primaire (1\textsuperscript{re} injection J0)}

% Secondary response curve
\addplot[popcurve, poporange, very thick, domain=60:120, samples=200] 
    ({x}, {1000/(1+exp(-(x-63)/2)) * exp(-(x-63)/80) + 50});
\addlegendentry{Réponse secondaire (2\textsuperscript{e} injection J60)}

% Control primary response at J60 (naive animal)
\addplot[popcurve, popgreen, dashed, thick, domain=60:120, samples=200] 
    ({x}, {800/(1+exp(-(x-78)/4)) * exp(-(x-78)/60) + 20});
\addlegendentry{Témoin : réponse primaire chez animal naïf (inj. J60)}

% Vertical line for second injection
\draw[popdark, dashed, thick] (axis cs:60,0) -- (axis cs:60,1100) node[above, pos=0.95, popdark] {2\textsuperscript{e} injection};
% Annotations
\node[popblue, anchor=south east] at (axis cs:25,450) {Latence $\approx$ 7-10 j};
\node[popblue, anchor=north east] at (axis cs:25,400) {Pic modéré};
\node[poporange, anchor=south west] at (axis cs:68,800) {Latence $\approx$ 1-3 j};
\node[poporange, anchor=north west] at (axis cs:68,750) {Pic élevé \& durable};
\end{axis}
\end{tikzpicture}
\legende{Cinétique de la production d'anticorps lors d'une primo-infection (réponse primaire) et d'une réinfection ou rappel vaccinal (réponse secondaire). La réponse secondaire est plus précoce, plus intense, d'affinité plus forte et plus persistante. UA = Unités Arbitraires.}
\end{popfigure}

\begin{methode}[Exploiter une courbe de cinétique des anticorps (classique au Bac)]
\textbf{Objectif :} Comparer la réponse primaire et la réponse secondaire pour caractériser la mémoire immunitaire.
\textbf{Étapes :}
\begin{enumerate}
    \item \textbf{Identifier les phases :} Repérer sur l'axe des abscisses le moment des injections (flèches verticales). Délimiter la réponse primaire (après 1\textsuperscript{re} injection) et la réponse secondaire (après 2\textsuperscript{e} injection).
    \item \textbf{Mesurer la latence :} Temps entre l'injection et le début de la montée significative du taux d'anticorps (seuil de détection).
        \begin{itemize}
            \item Réponse primaire : latence longue (généralement 7 à 14 jours chez le lapin, 10-14 jours chez l'homme).
            \item Réponse secondaire : latence courte (1 à 3 jours).
        \end{itemize}
    \item \textbf{Comparer l'intensité (taux maximal) :} Lire l'ordonnée du pic.
        \begin{itemize}
            \item Réponse primaire : taux maximal faible à modéré (souvent $<$ 100-200 UA/mL ici).
            \item Réponse secondaire : taux maximal élevé (souvent 10 à 100 fois supérieur, $>$ 800 UA/mL ici).
        \end{itemize}
    \item \textbf{Comparer la durée / persistance :} Observer la pente de la décroissance (phase de catabolisme des anticorps).
        \begin{itemize}
            \item Réponse primaire : décroissance rapide (demi-vie courte des IgM $\approx$ 5 j, puis IgG sans rappel $\approx$ 21-28 j).
            \item Réponse secondaire : décroissance lente, maintien d'un taux protecteur pendant des mois/années (plasmocytes de vie longue en moelle osseuse + réactivation mémoires).
        \end{itemize}
    \item \textbf{Comparer la qualité (isotype/affinité) :} La réponse primaire est dominée par les IgM puis IgG de faible affinité ; la réponse secondaire est dominée précocement par des IgG de haute affinité (maturation d'affinité).
    \item \textbf{Conclure :} La réponse secondaire est plus précoce, plus intense, plus durable, de meilleure qualité $\rightarrow$ existence d'une mémoire immunitaire.
\end{enumerate}
\textbf{Attention :} Ne pas confondre « latence » (pas d'anticorps détectables) et « phase de latence clinique » (maladie silencieuse). Préciser l'unité du temps (jours/semaines) et du taux (mg/L, UI/mL, ou UA).
\end{methode}

\begin{exoresolu}[Analyse d'une courbe de vaccination antitétanique]
\textbf{Énoncé :} Un patient reçoit une primo-vaccination antitétanique (J0) puis un rappel à J30. La figure ci-contre (similaire à la Figure 1) montre l'évolution du taux d'anticorps antitétaniques. On considère qu'un taux $\ge$ 0,1 UI/mL est protecteur.
\begin{enumerate}
    \item Préciser la nature de l'immunité acquise (active/passive, naturelle/artificielle).
    \item Déterminer la latence, le pic et la durée de protection pour chaque réponse.
    \item Expliquer la différence de cinétique au niveau cellulaire.
\end{enumerate}
\textbf{Solution détaillée :}
\begin{enumerate}
    \item \textbf{Immunité active artificielle :} Le patient reçoit des antigènes (anatoxine tétanique) qui stimulent \emph{ses propres} lymphocytes B et T. Ce n'est pas un transfert d'anticorps (sérothérapie = immunité passive).
    \item \textbf{Réponse primaire (J0-J30) :} Latence $\approx$ 10-14 j. Pic $\approx$ 0,5-1 UI/mL vers J21 (principalement IgM puis IgG). Taux retombe sous le seuil protecteur (0,1 UI/mL) vers J50-J60 sans rappel. Protection temporaire et incomplète.
    \item \textbf{Réponse secondaire (après J30) :} Latence $\approx$ 2-3 j. Pic $\approx$ 10-50 UI/mL vers J37-J40 (IgG haute affinité). Taux reste $>$ 0,1 UI/mL pendant plusieurs années (souvent $>$ 10 ans). Protection durable.
    \item \textbf{Explication cellulaire :} Lors de la primo-vaccination, les lymphocytes B naïfs spécifiques de l'anatoxine s'activent (avec aide des LTh), prolifèrent (expansion clonale) et se différencient en \textbf{plasmocytes} (usines à anticorps, vie courte pour la majorité, vie longue pour une minorité en moelle osseuse) et en \textbf{lymphocytes B mémoires} (vie longue, ne sécrètent pas d'anticorps au repos, BCR d'affinité mature). Lors du rappel, les lymphocytes B mémoires, plus nombreux, à seuil d'activation plus bas et à affinité mature, se différencient massivement et rapidement en plasmocytes à haute productivité, sans passer par la phase de sélection clonale lente des centres germinatifs.
\end{enumerate}
\end{exoresolu}

\courssub{Base cellulaire de la mémoire immunitaire : les lymphocytes mémoires}

\begin{definition}[Mémoire immunitaire]
Capacité du système immunitaire adaptatif de conserver, après une primo-rencontre avec un antigène, une \cle{population de lymphocytes mémoires} (B et T) spécifiques de cet antigène, permettant une réponse secondaire plus rapide, plus intense, de meilleure qualité (affinité) et plus durable lors d'une réexposition.
\end{definition}

\begin{definition}[Lymphocyte mémoire (B ou T)]
Lymphocyte issu de la différenciation d'un lymphocyte naïf activé lors de la réponse primaire (dans les centres germinatifs pour les B). Il ne sécrète pas d'effecteurs (anticorps, cytokines cytotoxiques) au repos. Il se caractérise par :
\begin{itemize}
    \item Une \textbf{longue durée de vie} (années, voire toute la vie pour les B mémoires et plasmocytes moelle osseuse ; décennies pour les T mémoires).
    \item Un \textbf{seuil d'activation abaissé} : il répond à de plus faibles concentrations d'antigène et nécessite moins de signaux de costimulation (CD28/B7).
    \item Une \textbf{réactivité accrue} : prolifération et différenciation plus rapides en cellules effectrices (plasmocytes pour B, LTc/LTh pour T).
    \item Pour les B mémoires : des récepteurs (BCR) d'\textbf{affinité mature} (résultat de l'hypermutation somatique + sélection d'affinité dans les centres germinatifs) et commutation de classe (IgG, IgA, IgE).
    \item Pour les T mémoires : distinction entre \textbf{T mémoires centraux ($T_{CM}$)} (homing ganglions, fort potentiel prolifératif) et \textbf{T mémoires effecteurs ($T_{EM}$)} (homing tissus périphériques, fonction effectrice immédiate).
\end{itemize}
\end{definition}

\emph{Mécanisme établi :} Lors de la réponse primaire dans les organes lymphoïdes secondaires (ganglions, rate, amygdales, plaques de Peyer), l'activation clonale des lymphocytes B et T naïfs (reconnaissance antigène + costimulation + cytokines) donne deux lignées parallèles :
\begin{itemize}
    \item \textbf{Lignée effectrice de courte durée (majoritaire) :} Plasmocytes (sécrétion massive d'anticorps, meurent par apoptose en quelques jours/semaines, sauf minorité moelle osseuse) et Lymphocytes T effecteurs (cytotoxiques ou auxiliaires, meurent par apoptose après élimination du pathogène = \emph{contraction}).
    \item \textbf{Lignée mémoire de longue durée (minoritiaire mais cruciale) :} Lymphocytes B mémoires (recirculent sang/lymphes, résident en zone marginale rate/ganglion) et Lymphocytes T mémoires ($T_{CM}$ organes lymphoïdes, $T_{EM}$ tissus).
\end{itemize}
C'est la persistance de cette \textbf{deuxième lignée} qui constitue le substrat cellulaire de la mémoire.

\begin{popfigure}
\begin{tikzpicture}[
    node distance=1.2cm and 1.5cm,
    box/.style={rectangle, draw=popdark, thick, rounded corners, fill=popcream, minimum height=1cm, minimum width=2.8cm, align=center, font=\small},
    arrow/.style={-Stealth, thick, popdark},
    branchB/.style={-Stealth, thick, popblue},
    branchT/.style={-Stealth, thick, poporange}
]
\node[box, fill=popblueL, align=center] (Ag){Antigène \\ (Primo-infection / Vaccination)};
\node[box, below=of Ag, align=center] (Naive){Lymphocyte naïf spécifique \\ (B ou T CD4+/CD8+)};
\node[box, fill=poporangeL, below left=of Naive, xshift=-1.2cm, align=center] (Effecteur){Lignée effectrice \\ \textbf{Courte durée (semaines)} \\ Plasmocytes / LT effecteurs \\ $\rightarrow$ Élimination pathogène};
\node[box, fill=popgreenL, below right=of Naive, xshift=1.2cm, align=center] (Memoire){Lignée mémoire \\ \textbf{Longue durée (années/vie)} \\ B mémoires / T mémoires ($T_{CM}$, $T_{EM}$) \\ $\rightarrow$ Surveillance immunologique};
\node[box, below=of Memoire, yshift=-0.5cm, align=center] (Reexpo){Réexposition \\ (Même antigène)};
\node[box, fill=poporangeL, below left=of Reexpo, xshift=-1.2cm, align=center] (Effecteur2){Réponse secondaire : \\ Différenciation massive \\ et rapide en effecteurs};
\node[box, fill=popgreenL, below right=of Reexpo, xshift=1.2cm, align=center] (Memoire2){Maintien / Expansion \\ du pool mémoires};

\draw[arrow] (Ag) -- (Naive);
\draw[branchB] (Naive) -- node[left, pos=0.5, font=\small] {Activation / Prolifération / Différenciation} (Effecteur);
\draw[branchT] (Naive) -- node[right, pos=0.5, font=\small] {Activation / Prolifération / Différenciation} (Memoire);
\draw[arrow] (Memoire) -- (Reexpo);
\draw[branchB] (Reexpo) -- (Effecteur2);
\draw[branchT] (Reexpo) -- (Memoire2);
\end{tikzpicture}
\legende{Sort des lymphocytes naïfs activés lors de la primo-infection : bifurcation vers une lignée effectrice (élimination immédiate, mort programmée) et une lignée mémoire (protection future, longue survie).}
\end{popfigure}

\courssub{La vaccination : simuler une primo-infection sans risque}

\begin{definition}[Vaccination]
Administration délibérée d'un \cle{vaccin} (préparation antigénique) chez un sujet sain pour induire une \cle{immunité active artificielle} : stimulation d'une réponse primaire (production d'anticorps et génération de lymphocytes mémoires) sans provoquer la maladie, afin de protéger contre une infection future.
\end{definition}

\emph{Principe :} Le vaccin mime l'antigène du pathogène. Il déclenche la réponse primaire (latence, pic modéré, génération mémoires) et laisse en place le pool de lymphocytes mémoires. Lors de l'infection naturelle, la réponse secondaire neutralise l'agent avant qu'il ne cause des dégâts cliniques.

\begin{center}
\begin{tabularx}{\linewidth}{|l|X|X|X|}
\hline
\rowcolor{popblueL}
\textbf{Type de vaccin} & \textbf{Composition / Principe} & \textbf{Exemples (Calendrier Cameroun PEV 2024)} & \textbf{Avantages / Inconvénients} \\
\hline
\cle{Vaccin vivant atténué} & Micro-organisme vivant dont la virulence a été abolie (passages en culture, mutations) mais garde pouvoir réplicatif limité. & \textbf{BCG} (tuberculose - naissance), \textbf{Rougeole-Rubéole} (9 \& 15 mois), \textbf{Fièvre jaune} (9 mois), \textbf{Rotavirus} (6 \& 10 sem), \textbf{Polio oral (VPO - campagnes)}. & \textbf{+} Immunité forte (humorale + cellulaire CD8), durable, souvent 1-2 doses. Mimé infection naturelle. \textbf{-} Contre-indiqué immunodéprimés (VIH avancé), grossesse ; risque rare de réversion virulente (VPO) ; chaîne du froid stricte. \\
\hline
\cle{Vaccin inactivé (tué)} & Micro-organisme entier tué (chaleur, formol, $\beta$-propiolactone). Ne se réplique pas. & \textbf{Polio injectable (VPI)} (6, 10, 14 sem), \textbf{Coqueluche (acellulaire - dans Penta)}, \textbf{Hépatite A}, \textbf{Rage}. & \textbf{+} Sûr (immunodéprimés, grossesse), stable. \textbf{-} Immunité plus faible (surtout humorale Th2, peu CD8), nécessite \textbf{rappels} + \textbf{adjuvant} obligatoire (Alum). \\
\hline
\cle{Vaccin sous-unitaire / anatoxine} & Fraction purifiée (protéine, polysaccharide) ou toxine inactivée (anatoxine = toxine détoxiquée gardant immunogénicité). & \textbf{Tétanos} (anatoxine), \textbf{Diphtérie} (anatoxine), \textbf{Hépatite B} (protéine S recombinante - levure), \textbf{HPV} (VLP - particules pseudo-virales), \textbf{Méningocoque} (polysaccharide conjugué). & \textbf{+} Très sûr, pas de composant vivant, production standardisée. \textbf{-} Faiblement immunogène $\rightarrow$ adjuvant + \textbf{rappels} indispensables. \textbf{Conjugaison} (polysaccharide + protéine porteuse tétanique/diphtérique) essentielle pour réponse T-dépendante, commutation de classe et mémoire chez nourrisson ($<$ 2 ans). \\
\hline
\cle{Vaccin à ARNm / vecteur viral} & Acide nucléique codant l'antigène (ARNm lipidique) ou vecteur viral non réplicatif (adénovirus chimpanzé/humain) l'exprimant dans cellules hôtes. & \textbf{COVID-19} (Pfizer/BioNTech, Moderna = ARNm ; AstraZeneca = vecteur). & \textbf{+} Développement rapide, pas de culture pathogène, forte immunité cellulaire (CD8) + humorale, pas d'adjuvant classique (ARNm = adjuvant propre). \textbf{-} Chaîne du froid ultra-stricte (ARNm : -80°C à -20°C), surveillance pharmacovigilance long terme (myocardite rare). \\
\hline
\end{tabularx}
\end{center}

\emph{Le calendrier vaccinal au Cameroun (PEV 2024) :} Il protège contre 12 maladies cibles : Tuberculose (BCG à la naissance), Polio (VPI aux 6, 10, 14 sem + VPO campagnes), Diphtérie-Tétanos-Coqueluche-Hépatite B-Hib (Pentavalent aux 6, 10, 14 sem), Pneumocoque (PCV13 même calendrier), Rotavirus (2 doses 6 \& 10 sem), Rougeole-Rubéole (9 et 15 mois), Fièvre jaune (9 mois), HPV (filles 9-14 ans, 2 doses à 6 mois d'intervalle), COVID-19 (adultes). \textbf{Les rappels (DTCoq à 18 mois, 6 ans, 12 ans, puis tous les 10-20 ans ; Fièvre jaune tous les 10 ans ; Tétanos femmes enceintes VAT 1 à 5) sont cruciaux} pour réactiver la mémoire, maintenir le taux d'anticorps protecteur et élargir le répertoire clonale.

\begin{attention}[Pièges fréquents sur la vaccination (Bac)]
\begin{itemize}
    \item \textbf{Confondre vaccination et sérothérapie :} La vaccination induit une immunité \emph{active} (le sujet produit ses anticorps + mémoire). La sérothérapie apporte une immunité \emph{passive} (anticorps tout faits, pas de mémoire). Voir section suivante.
    \item \textbf{Oublier l'adjuvant :} Pour les vaccins inactivés/sous-unitaires/anatoxines, l'adjuvant (sels d'aluminium, émulsions MF59, AS01) est \emph{indispensable} pour activer l'immunité innée (recrutement/activation cellules dendritiques $\rightarrow$ costimulation CD80/86, cytokines IL-12) et permettre la réponse adaptative T-dépendante. Sans adjuvant $\rightarrow$ tolérance ou réponse nulle/faible.
    \item \textbf{Croire qu'un vaccin vivant atténué « donne la maladie » :} Il peut donner une forme \emph{atténuée, asymptomatique ou pauci-symptomatique} (ex. fièvre modérée post-Rougeole/BCG, adénopathie BCGite), mais pas la maladie grave sauvage. Contre-indiqué chez l'immunodéprimé sévère (risque de dissémination vaccinale).
    \item \textbf{Négliger les rappels :} La mémoire immunitaire n'est pas éternelle sans stimulation antigénique (naturelle ou vaccinale). Le taux d'anticorps chute (catabolisme). Le rappel (boost) réactive les mémoires $\rightarrow$ nouveau pic + expansion du pool mémoires + maturation d'affinité supplémentaire. \emph{Exemple :} L'immunité anti-tétanique chute sous seuil protecteur (0,1 UI/mL) $\approx$ 10 ans sans rappel.
    \item \textbf{Confondre « efficacité vaccinale » (infection) et « efficacité clinique » (maladie grave) :} Un vaccin peut ne pas empêcher l'infection/portage (ex. COVID-19 variant Omicron, Coqueluche acellulaire) mais empêcher la maladie grave/hospitalisation. La protection stérilisante (empêche toute infection) est rare (ex. vaccin rougeole, VPI).
\end{itemize}
\end{attention}

\courssub{Vaccination et sérothérapie : deux stratégies, deux natures d'immunité}

Il est capital de distinguer ces deux approches thérapeutiques/prophylactiques, souvent confondues par les élèves.

\begin{popfigure}
\begin{tikzpicture}[
    node distance=0.8cm and 1.5cm,
    mainbox/.style={rectangle, draw=popdark, thick, rounded corners, fill=popcream, minimum height=1.1cm, text width=3.8cm, align=center, font=\small},
    titlebox/.style={rectangle, draw=popdark, thick, rounded corners, fill=popblue, text=white, minimum height=0.8cm, text width=4cm, align=center, font=\bfseries\small},
    arrow/.style={-Stealth, thick, popdark},
]
% Column 1: Vaccination
\node[titlebox] (Vtitle) {VACCINATION (Immunité Active Artificielle)};
\node[mainbox, below=of Vtitle, align=center] (V1){Injection d'\textbf{ANTIGÈNES} \\ (vivants atténués, inactivés, sous-unitaires, ARNm)};
\node[mainbox, below=of V1, align=center] (V2){Reconnaissance par lymphocytes B/T naïfs du \textbf{sujet receveur} \\ (APC $\rightarrow$ MHC + Costimulation)};
\node[mainbox, below=of V2, align=center] (V3){Activation, prolifération clonale, différenciation \\ \textbf{Plasmocytes + Lymphocytes Mémoires (B \& T)}};
\node[mainbox, below=of V3, fill=popgreenL, align=center] (V4){Production d'anticorps \textbf{par le sujet} \\ + \textbf{Mémoire immunitaire durable}};
\node[mainbox, below=of V4, align=center] (V5){\textbf{Délai protection :} Jours à semaines \\ \textbf{Durée :} Années (vie entière) \\ \textbf{Rappels :} Nécessaires pour maintenir taux};

% Column 2: Serotherapy
\node[titlebox, right=4.5cm of Vtitle] (Stitle) {SÉROTHÉRAPIE (Immunité Passive Artificielle)};
\node[mainbox, below=of Stitle, align=center] (S1){Injection d'\textbf{ANTICORPS} tout faits \\ (Sérum hyperimmune, Immunoglobulines poly/monoclonales) \\ \emph{Origine : donneur immunisé (humain/cheval) ou recombinant}};
\node[mainbox, below=of S1, align=center] (S2){Pas de stimulation du système immunitaire \\ du receveur (pas d'antigène injecté, pas d'APC)}; 
\node[mainbox, below=of S2, align=center] (S3){Anticorps circulent immédiatement \\ \textbf{Aucun lymphocyte mémoire produit} \\ \textbf{Aucune maturation d'affinité}};
\node[mainbox, below=of S3, fill=poporangeL, align=center] (S4){Protection \textbf{immédiate mais transitoire} \\ \textbf{Pas de mémoire immunitaire}};
\node[mainbox, below=of S4, align=center] (S5){\textbf{Délai protection :} Immédiat (heures) \\ \textbf{Durée :} Semaines (demi-vie IgG $\approx$ 21-28 j) \\ \textbf{Rappels :} Inutiles (pas de mémoire), réinjection si besoin};

% Arrows
\draw[arrow] (V1) -- (V2);
\draw[arrow] (V2) -- (V3);
\draw[arrow] (V3) -- (V4);
\draw[arrow] (V4) -- (V5);
\draw[arrow] (S1) -- (S2);
\draw[arrow] (S2) -- (S3);
\draw[arrow] (S3) -- (S4);
\draw[arrow] (S4) -- (S5);

% Comparison labels
\node[font=\small\bfseries, popblue, align=center] at ($(Vtitle.north west)!0.5!(V5.south west) + (-1.8,0)$) {Le sujet \textbf{apprend} \\ à se défendre};
\node[font=\small\bfseries, poporange, align=center] at ($(Stitle.north east)!0.5!(S5.south east) + (1.8,0)$) {On \textbf{donne} les armes \\ toutes faites};
\end{tikzpicture}
\legende{Comparatif schématique : Vaccination (immunité active, mémoire, délai, durée) vs Sérothérapie (immunité passive, pas de mémoire, immédiat, bref).}
\end{popfigure}

\begin{definition}[Sérothérapie (Immunothérapie passive)]
Administration d'immunoglobulines (anticorps) purifiées, issues du sang de donneurs immunisés (sérum hyperimmune) ou produites par génie génétique (anticorps monoclonaux), pour conférer une \cle{protection immédiate mais temporaire} (immunité passive artificielle). Aucune mémoire immunitaire n'est induite chez le receveur.
\end{definition}

\begin{exemplebox}[Indications cliniques au Cameroun]
\begin{itemize}
    \item \textbf{Vaccination (Prophylaxie préventive) :} PEV nourrissons (BCG, Penta, Pneumo, Rota, VPI, RR, FJ), HPV adolescentes (9-14 ans), Fièvre jaune (voyageurs/endémie, 1 dose vie entière), COVID-19 (groupes cibles), Tétanos (femmes enceintes : protocole VAT 1, 2, 3, 4, 5 pour protéger la mère et le nouveau-né via transfert transplacentaire d'IgG).
    \item \textbf{Sérothérapie (Urgence curative / Prophylaxie post-exposition immédiate) :}
    \begin{itemize}
        \item \textbf{Envenimation serpent (vipère, cobra, mamba - Grand Nord, Est, Adamaoua) :} Antivenin polyvalent africain (ex. SAV-Aventis/Inosan, FAV-Afrique) $\rightarrow$ injection IV urgente d'anticorps F(ab')2 anti-venin. Sauve des vies, prévient nécrose/coagulopathie.
        \item \textbf{Exposition au virus de la rage (morsure/griffure animal suspect) :} Immunoglobulines anti-rabiques humaines (RIG) ou équines (ERIG) infiltrées \textbf{autour de la plaie} (dose restante IM) \textbf{+ Vaccin antirabique} (protocole Zagreb 2-1-1 ou Essen 4-5 doses). Le sérum comble le délai de 7-14 j de la vaccination active.
        \item \textbf{Tétanos maternel/néonatal (accouchement insalubre, mère non vaccinée) :} Immunoglobulines anti-tétaniques humaines (TIG) IM à la mère et/ou au nouveau-né \textbf{+ Vaccin anatoxine tétanique} (site distinct).
        \item \textbf{Diphtérie (cas confirmés/contacts) :} Antitoxine diphtérique (SAD) IV/IM en urgence (neutralise toxine libre) \textbf{+ antibiotiques} (érythromycine/azithromycine) \textbf{+ vaccination} (rappel DTC).
    \end{itemize}
\end{itemize}
\end{exemplebox}

\begin{attention}[La sérothérapie ne confère AUCUNE mémoire -- Point clé Bac]
C'est le point le plus sanctionné. Les anticorps injectés (IgG majoritairement) sont catabolisés avec une demi-vie d'environ \textbf{21 à 28 jours}. Au bout de 3 à 4 mois, ils ont disparu. Le receveur n'a \textbf{pas} produit de lymphocytes mémoires car il n'a pas reçu d'antigène (ou quantité négligeable dans le sérum). Si on veut une protection durable après une sérothérapie (ex. post-morsure rage, tétanos), \textbf{il faut OBLIGATOIREMENT associer une vaccination} (immunité active) qui prendra le relais quand les anticorps passifs auront disparu. 
\textbf{Ne jamais écrire :} « La sérothérapie permet de se souvenir de l'antigène » ou « La sérothérapie induit une mémoire immunitaire. » $\rightarrow$ \textbf{Faux.}
\end{attention}

\begin{savaistu}[Edward Jenner (1796) : la naissance de la vaccinologie]
En 1796, le médecin anglais \textbf{Edward Jenner} observe que les laitières contaminées par la \textbf{vaccine} (variole de la vache, \emph{cowpox}), maladie bénigne des bovins transmissible à l'homme (pustules sur les mains), ne contractent pas la \textbf{variole} humaine (mortelle à 30\%, défigurante). Il inocule du pus de vaccine au jeune James Phipps (8 ans) : celui-ci développe une pustule locale, guérit. Six semaines plus tard, Jenner l'inocule avec du pus de variole humaine (variolisation, pratique risquée) : l'enfant ne tombe pas malade. \textbf{Preuve de concept :} Un agent apparenté, non pathogène pour l'homme, protège contre le pathogène meurtrier. Le terme « \emph{vaccination} » vient de \emph{vacca} (vache en latin). La variole est la seule maladie humaine éradiquée (OMS, 1980) grâce à la vaccination mondiale (dernier cas naturel 1977, Somalie).
\end{savaistu}

\begin{aretenir}[Synthèse : Mémoire et Vaccination]
\begin{itemize}
    \item La \textbf{réponse primaire} (primo-infection/vaccination) est lente (latence 1-2 sem), de faible amplitude, dominée par IgM puis IgG faible affinité, et génère des \textbf{lymphocytes B et T mémoires} à longue durée de vie + plasmocytes moelle osseuse.
    \item La \textbf{réponse secondaire} (réinfection/rappel) est rapide (latence 1-3 j), intense (taux d'anticorps $\times$ 10-100), durable, dominée par IgG haute affinité, grâce à la réactivation des lymphocytes mémoires (seuil bas, nombre élevé, affinité mature).
    \item La \textbf{vaccination} utilise des antigènes (atténués, inactivés, sous-unitaires, ARNm) pour provoquer une réponse primaire \textbf{sans maladie} et installer la mémoire. Les \textbf{rappels} entretiennent le taux d'anticorps protecteur, élargissent le pool mémoires et améliorent l'affinité.
    \item La \textbf{sérothérapie} apporte des anticorps tout faits (immunité passive) : protection \textbf{immédiate} mais \textbf{brève} (semaines, demi-vie IgG $\approx$ 21-28 j), \textbf{sans mémoire}. Indiquée en urgence vitale (rage, tétanos, envenimation, diphtérie) \textbf{toujours associée à la vaccination} pour la protection durable.
    \item Au Cameroun, le PEV, la riposte aux épidémies (choléra, méningite, COVID, fièvre jaune, polio) et la prise en charge des urgences toxiques/infectieuses reposent sur cette compréhension rigoureuse de la mémoire immunitaire.
\end{itemize}
\end{aretenir}

\coursec{Les dysfonctionnements : maladies auto-immunes}

Après avoir étudié comment le système immunitaire construit une mémoire protectrice grâce à la vaccination, nous devons aborder une situation paradoxale : que se passe-t-il lorsque ce même système, au lieu de défendre l'organisme, se retourne contre lui ? C'est l'objet des \textbf{maladies auto-immunes}, une rupture de la tolérance immunologique qui touche environ 5 à 8 \% de la population mondiale, avec une prévalence croissante, y compris au Cameroun où les services d'endocrinologie et de neurologie des hôpitaux centraux (Yaoundé, Douala) en suivent de nombreux cas.

\courssub{Définition et mécanisme général : la rupture de la tolérance au soi}

\begin{definition}[Maladie auto-immune, auto-anticorps, rupture de tolérance]
Une \cle{maladie auto-immune} est une pathologie résultant d'une \textbf{rupture de la tolérance immunologique au soi} : le système immunitaire reconnaît anormalement certains constituants de l'organisme (protéines, récepteurs, ADN, etc.) comme du \emph{non-soi} et déclenche une réponse immune (production d'\cle{auto-anticorps} et/ou activation de \cle{lymphocytes T auto-réactifs}) dirigée contre ses propres tissus, entraînant des lésions inflammatoires et une dysfonction d'organe.
\end{definition}

Normalement, lors de la maturation dans le thymus (pour les LT) et la moelle osseuse (pour les LB), les lymphocytes dont les récepteurs reconnaissent fortement les antigènes propres (\emph{self}) sont éliminés par \textbf{sélection négative} (apoptose) ou rendus anergiques (inactifs). C'est la \textbf{tolérance centrale}. Une \textbf{tolérance périphérique} (régulation par les LT régulateurs $T_{reg}$, ignorance antigénique) prend le relais pour les clones échappés.

\begin{propriete}[Mécanismes de la rupture de tolérance]
La survenue d'une maladie auto-immune résulte de la conjonction de trois facteurs (triade étiologique) :
\begin{enumerate}
    \item \textbf{Facteurs génétiques :} Polymorphismes du \cle{CMHC} (HLA chez l'homme) prédisposant à présenter certains auto-antigènes (ex : HLA-DR3/DR4 pour le diabète type 1, HLA-B27 pour la spondylarthrite ankylosante).
    \item \textbf{Facteurs environnementaux déclenchants :} Infections virales/bactériennes (\emph{mimétisme moléculaire} : épitopes pathogènes ressemblant à des épitopes propres), expositions chimiques, stress, modifications du microbiote.
    \item \textbf{Défaillance des mécanismes de régulation :} Déficit en $T_{reg}$ (FoxP3+), défaut d'apoptose des lymphocytes activés (Fas/FasL), inflammation chronique maintenant l'activation.
\end{enumerate}
\end{propriete}

\begin{experience}[Mise en évidence du mimétisme moléculaire (modèle expérimental)]
\textbf{Observation :} Après une infection à \emph{Streptococcus pyogenes} (angine), certains patients développent un rhumatisme articulaire aigu (maladie auto-immune touchant cœur, articulations, peau).
\textbf{Hypothèse :} Des antigènes bactériens partagent des épitopes avec des protéines humaines (myosine cardiaque, protéines articulaires).
\textbf{Expérience :} On immunise des souris avec une protéine M de \emph{S. pyogenes}. On isole leurs anticorps et on teste leur réactivité sur des coupes de cœur de souris.
\textbf{Résultat :} Les anticorps anti-protéine M se fixent sur la myosine cardiaque des souris (immunofluorescence positive).
\textbf{Interprétation :} La réponse immune anti-bactérienne \textbf{cible par mimétisme} un antigène propre. Les lymphocytes B/T activés contre la bactérie cross-réagissent avec le soi.
\textbf{Conclusion :} Le mimétisme moléculaire est un mécanisme majeur de rupture de tolérance périphérique.
\end{experience}

\begin{attention}[Piège fréquent : Auto-immunité vs Immunodéficience]
Ne confondez jamais :
\begin{itemize}
    \item \textbf{Maladie auto-immune :} Immunité \textbf{excessive} et \textbf{maldirigée} contre le soi (ex : diabète type 1, polyarthrite rhumatoïde). Le système immunitaire est \emph{hyperactif} mais \emph{mal éduqué}.
    \item \textbf{Immunodéficience (ex : Sida) :} Immunité \textbf{insuffisante} (déficit en LT4, défaut de réponse). Le système immunitaire est \emph{affaibli} ou \emph{absent}.
\end{itemize}
Au Bac, si on vous demande « Comment le VIH cause-t-il le Sida ? », ne répondez pas « C'est une maladie auto-immune ». Le VIH détruit les LT4 ; il ne déclenche pas d'auto-anticorps pathogènes directs (même s'il y a une activation immune chronique).
\end{attention}

\courssub{Exemples types au programme : mécanismes, signes et traitements}

Nous étudierons trois exemples emblématiques illustrant les deux grandes cibles de l'auto-immunité : les \textbf{cellules} (diabète type 1, vitiligo) et les \textbf{récepteurs membranaires} (myasthénie).

\textbf{1. Le diabète sucré de type 1 (diabète juvénile) : destruction des cellules $\beta$ pancréatiques.}

\begin{experience}[Démonstration de l'origine auto-immune du diabète type 1 (modèle NOD - Non Obese Diabetic mouse)]
\textbf{Observation :} La souche de souris NOD développe spontanément un diabète type 1 vers 12-20 semaines. L'infiltrat pancréatique (\emph{insulite}) est riche en LT4 et LT8.
\textbf{Expérience de transfert adoptif :} On prélève des lymphocytes de rate de souris NOD diabétiques et on les injecte à des souris NOD jeunes (prédiabétiques) ou à des souris immunodéficientes (nu/nu).
\textbf{Résultat :} Les souris receveuses développent rapidement le diabète et une insulite. Le transfert de sérum seul (anticorps) ne transfère pas la maladie efficacement.
\textbf{Interprétation :} La maladie est \textbf{transférable par des cellules lymphoïdes (LT)}, prouvant le rôle central des \textbf{lymphocytes T auto-réactifs (CD8+ cytotoxiques et CD4+ Th1)} dans la destruction des cellules $\beta$. Les auto-anticorps (anti-GAD, anti-IA2, anti-insuline, anti-ZnT8) sont des \textbf{marqueurs diagnostiques} précoces mais pas les effecteurs principaux de la destruction.
\end{experience}

\begin{propriete}[Physiopathologie du diabète type 1]
\begin{itemize}
    \item \textbf{Cible :} Cellules $\beta$ des îlots de Langerhans (productrices d'insuline).
    \item \textbf{Mécanisme :} Infiltration par des LT8+ cytotoxiques (reconnaissance peptide-HLA I) et LT4+ Th1 (secrétion d'IFN-$\gamma$, TNF-$\alpha$, IL-2) $\rightarrow$ apoptose des cellules $\beta$ (perforine/granzyme, Fas/FasL, cytokines pro-inflammatoires).
    \item \textbf{Seuil clinique :} La destruction est souvent \textbf{supérieure à 80-90 \%} au moment de l'hyperglycémie symptomatique (polydipsie, polyurie, amaigrissement, cétoacidose).
    \item \textbf{Marqueurs biologiques :} Auto-anticorps anti-GAD65 (glutamate décarboxylase), anti-IA2, anti-insuline (IAA), anti-ZnT8. Présents des mois/années avant le diagnostic.
\end{itemize}
\end{propriete}

\begin{popfigure}
\begin{tikzpicture}[scale=0.9, every node/.style={font=\small}]
    % Couleurs
    \colorlet{betahealthy}{popgreen!80!popdark}
    \colorlet{betadestroyed}{poporange!80!popdark}
    \colorlet{lympho}{poppurple}
    \colorlet{insulin}{popblue}
    \colorlet{blood}{poppink!70!popdark}
    \colorlet{capillary}{popmuted}
    
    % --- PANCREAS SAIN (Gauche) ---
    \begin{scope}[xshift=-7cm]
        % Capillaire
        \draw[capillary, line width=2pt] (-1.5,0) -- (1.5,0);
        \draw[capillary, line width=2pt] (-1.5,-0.1) -- (1.5,-0.1);
        \node[capillary, anchor=north] at (0,0.2) {Capillaire sanguin};
        
        % Îlot de Langerhans
        \draw[popdark, thick, fill=popcream] (0,-2) ellipse (1.2cm and 0.8cm);
        \node[popdark, anchor=north] at (0,-1.1) {Îlot de Langerhans sain};
        
        % Cellules bêta (majoritaires au centre)
        \foreach \angle/\r in {30/0.5, 90/0.4, 150/0.5, 210/0.4, 270/0.5, 330/0.4, 0/0.3, 60/0.3, 120/0.3, 180/0.3, 240/0.3, 300/0.3} {
            \pgfmathsetmacro{\x}{\r*cos(\angle)}
            \pgfmathsetmacro{\y}{\r*sin(\angle)}
            \draw[betahealthy, fill=betahealthy!30, thick] (\x,\y-2) circle (0.12cm);
            % Granules d'insuline
            \foreach \i in {1,2,3} {
                \pgfmathsetmacro{\ang}{rand*360}
                \pgfmathsetmacro{\rad}{0.02+rand*0.05}
                \fill[insulin] ($(\x,\y-2)+(\ang:\rad)$) circle (0.015cm);
            }
        }
        \node[betahealthy, anchor=north] at (0,-3.2) {Cellules $\beta$ (insuline +++)};
        
        % Cellules alpha (périphérie)
        \foreach \angle in {45, 135, 225, 315} {
            \pgfmathsetmacro{\x}{1.0*cos(\angle)}
            \pgfmathsetmacro{\y}{0.6*sin(\angle)}
            \draw[popgold, fill=popgold!30, thick] (\x,\y-2) circle (0.1cm);
        }
        \node[popgold, anchor=south] at (0,-0.9) {Cellules $\alpha$ (glucagon)};
        
        % Flèche sécrétion insuline
        \draw[insulin, ->, thick, >=Stealth] (0,-2.8) -- (0,-1.2) node[midway, right] {Sécrétion d'insuline};
    \end{scope}
    
    % Flèche transition
    \draw[popdark, ->, very thick, >=Stealth] (-5.5,-3.5) -- (-3.5,-3.5) node[midway, above, popdark] {\textbf{Auto-immunité}};
    \draw[popdark, ->, very thick, >=Stealth] (-5.5,-4.0) -- (-3.5,-4.0) node[midway, above, popdark] {\textbf{(LT8+ cytotoxiques)}};
    
    % --- PANCREAS DIABÉTIQUE (Droite) ---
    \begin{scope}[xshift=3.5cm]
        % Capillaire
        \draw[capillary, line width=2pt] (-1.5,0) -- (1.5,0);
        \draw[capillary, line width=2pt] (-1.5,-0.1) -- (1.5,-0.1);
        \node[capillary, anchor=north] at (0,0.2) {Capillaire sanguin};
        
        % Îlot détruit
        \draw[popdark, thick, fill=popcream!50!poporangeL] (0,-2) ellipse (1.2cm and 0.8cm);
        \node[popdark, anchor=north] at (0,-1.1) {Îlot de Langerhans détruit (insulite)};
        
        % Débris cellulaires / fibrose
        \draw[poporange, fill=poporange!20, thick] (0,-2) ellipse (0.8cm and 0.5cm);
        \node[poporange, anchor=north] at (0,-2.0) {Fibrose / Débris};
        
        % Lymphocytes T infiltrants
        \foreach \i in {1,...,15} {
            \pgfmathsetmacro{\x}{-1.0+rand*2.0}
            \pgfmathsetmacro{\y}{-2.5+rand*1.0}
            \draw[lympho, fill=lympho!20, thick] (\x,\y) circle (0.07cm);
            % Noyau
            \fill[lympho!80!black] (\x,\y) circle (0.03cm);
        }
        \node[lympho, anchor=north] at (0,-3.3) {Infiltrat lymphocytaire (LT8+ auto-réactifs)};
        
        % Cellules bêta résiduelles (rares, en apoptose)
        \foreach \angle/\r in {90/0.6, 270/0.6} {
            \pgfmathsetmacro{\x}{\r*cos(\angle)}
            \pgfmathsetmacro{\y}{\r*sin(\angle)}
            \draw[betadestroyed, fill=betadestroyed!30, thick, dashed] (\x,\y-2) circle (0.12cm);
            % Noyau fragmenté (apoptose)
            \draw[betadestroyed, thick] (\x-0.04,\y-2) -- (\x+0.04,\y-2);
            \draw[betadestroyed, thick] (\x,\y-2.04) -- (\x,\y-1.96);
        }
        \node[betadestroyed, anchor=south] at (0,-3.7) {Cellules $\beta$ résiduelles (apoptose) $< 20\%$};
        
        % Pas d'insuline sécrétée
        \draw[insulin, ->, thick, >=Stealth, dashed] (0,-2.8) -- (0,-1.2) node[midway, right, opacity=0.5] {Sécrétion nulle};
    \end{scope}
    
    % Légendes communes
    \node[align=center, font=\bfseries\small] at (-7, 1.5) {Pancréas \textbf{SAIN}};
    \node[align=center, font=\bfseries\small] at (3.5, 1.5) {Pancréas \textbf{DIABÉTIQUE TYPE 1}};
\end{tikzpicture}
\legende{Comparaison histologique schématique d'un îlot de Langerhans sain (gauche) et détruit par l'insulite auto-immune (droite) dans le diabète de type 1. À gauche : cellules $\beta$ riches en granules d'insuline (bleu), cellules $\alpha$ périphériques (or). À droite : infiltration massive par des lymphocytes T CD8+ cytotoxiques (violet), fibrose, rares cellules $\beta$ résiduelles en apoptose (pointillés). L'absence d'insuline entraîne l'hyperglycémie.}
\end{popfigure}

\begin{exemplebox}[Prise en charge au Cameroun]
Au Cameroun, le diagnostic repose sur l'hyperglycémie à jeun $\ge 1,26$ g/L (ou $7,0$ mmol/L) ou $> 2,00$ g/L à toute heure + signes cliniques, confirmée par la présence d'auto-anticorps (anti-GAD, anti-IA2) si disponible (labos de référence : Centre Pasteur, Hôpitaux Centraux/Universitaires). Le traitement est \textbf{vital et à vie} : injections d'insuline (souvent schémas basal-bolus : insuline lente + rapide aux repas). L'éducation thérapeutique (autosurveillance glycémique, alimentation, reconnaissance de l'hypoglycémie) est cruciale. L'accès aux bandelettes et à l'insuline (chaîne du froid) reste un défi de santé publique dans les zones rurales.
\end{exemplebox}

\textbf{2. Le vitiligo : destruction des mélanocytes.}

\begin{propriete}[Vitiligo]
\begin{itemize}
    \item \textbf{Cible :} \textbf{Mélanocytes} (cellules productrices de mélanine) de l'épiderme basal et des follicules pileux.
    \item \textbf{Mécanisme :} Auto-immunité \textbf{humorale} (auto-anticorps anti-mélanocytes : anti-tyrosinase, anti-gp100, anti-MART-1) \textbf{ET cellulaire} (LT8+ cytotoxiques spécifiques d'antigènes mélanocytaires, infiltration lymphocytaire au bord des plaques). Stress oxydatif intrinsèque des mélanocytes favorisant l'exposition d'antigènes cachés.
    \item \textbf{Signe clinique :} \textbf{Dépigmentation} par plaques achromiques (blanc laiteux), symétriques souvent, zones de frottement, visage, mains, orifices. Leukotrichie (poils blancs) possible.
    \item \textbf{Évolution :} Imprévisible, poussées et rémissions. Impact psychologique majeur (stigmatisation).
    \item \textbf{Traitement :} Corticoïdes topiques, inhibiteurs de calcineurine topiques (tacrolimus), photothérapie UVB à bande étroite (311 nm), greffe de mélanocytes (stade stable). Pas de traitement systémique curatif standard.
\end{itemize}
\end{propriete}

\textbf{3. La myasthénie auto-immune (myasthénie grave) : blocage des récepteurs à l'acétylcholine.}

C'est l'exemple type de la \textbf{maladie auto-immune à anticorps anti-récepteur} (type II de Gell et Coombs).

\begin{popfigure}
\begin{tikzpicture}[scale=0.95, every node/.style={font=\small}]
    % Couleurs
    \colorlet{axon}{popblue}
    \colorlet{muscle}{popred!70!popdark}
    \colorlet{ach}{popgreen}
    \colorlet{rach}{poppurple}
    \colorlet{ab}{poporange}
    \colorlet{mito}{popgold}
    
    % --- JONCTION NORMALE (Haut) ---
    \begin{scope}[yshift=4cm]
        % Terminaison axonique
        \draw[axon, fill=axon!10, thick] (-3,0) rectangle (3,1);
        \node[axon, anchor=south] at (0,1.1) {Terminaison motoneurone (axone)};
        % Vésicules synaptiques
        \foreach \i in {1,...,12} {
            \pgfmathsetmacro{\x}{-2.5+rand*4.5}
            \pgfmathsetmacro{\y}{0.15+rand*0.6}
            \draw[ach, fill=ach!30] (\x,\y) circle (0.06cm);
        }
        \node[ach, anchor=north] at (0,-0.2) {Vésicules d'acétylcholine (ACh)};
        % Mitochondries
        \foreach \i in {1,...,4} {
            \pgfmathsetmacro{\x}{-2.5+rand*4.5}
            \pgfmathsetmacro{\y}{0.15+rand*0.6}
            \draw[mito, fill=mito!30, thick] (\x,\y) ellipse (0.08cm and 0.04cm);
        }
        
        % Fente synaptique
        \draw[popline, thick] (-3,-0.1) -- (3,-0.1);
        \draw[popline, thick] (-3,-0.6) -- (3,-0.6);
        \node[popmuted, anchor=north] at (0,-0.35) {Fente synaptique ($\sim 50$ nm)};
        % ACh libérée
        \foreach \i in {1,...,8} {
            \pgfmathsetmacro{\x}{-2+rand*4}
            \pgfmathsetmacro{\y}{-0.2+rand*0.3}
            \node[ach, scale=0.7] at (\x,\y) {ACh};
        }
        
        % Membrane musculaire (plissée)
        \draw[muscle, thick, decorate, decoration={snake, amplitude=0.1cm, segment length=0.3cm}] (-3,-0.6) -- (3,-0.6);
        \node[muscle, anchor=north] at (0,-0.7) {Membrane musculaire (plissée)};
        
        % Récepteurs AChR (normaux)
        \foreach \x in {-2.5, -1.5, -0.5, 0.5, 1.5, 2.5} {
            \draw[rach, fill=rach!30, thick] (\x,-0.75) rectangle ++(0.6, -0.4);
            % Site de fixation
            \node[rach, scale=0.6] at (\x+0.3, -0.95) {AChR};
        }
        \node[rach, anchor=north] at (0,-1.3) {Récepteurs à l'ACh (AChR) fonctionnels};
        
        % Flèche dépolarisation
        \draw[popgreen, ->, very thick, >=Stealth] (0,-1.4) -- (0,-1.8) node[midway, right] {Ouverture canal Na$^+$ $\rightarrow$ Dépolarisation (PPP) $\rightarrow$ Potentiel d'action musculaire $\rightarrow$ \textbf{Contraction}};
    \end{scope}
    
    % Séparateur
    \draw[popline, dashed, thick] (-4, 2.8) -- (4, 2.8);
    
    % --- JONCTION MYASTHÉNIQUE (Bas) ---
    \begin{scope}[yshift=-1.5cm]
        % Terminaison axonique (identique)
        \draw[axon, fill=axon!10, thick] (-3,0) rectangle (3,1);
        \node[axon, anchor=south] at (0,1.1) {Terminaison motoneurone (normale)};
        \foreach \i in {1,...,12} {
            \pgfmathsetmacro{\x}{-2.5+rand*4.5}
            \pgfmathsetmacro{\y}{0.15+rand*0.6}
            \draw[ach, fill=ach!30] (\x,\y) circle (0.06cm);
        }
        
        % Fente synaptique
        \draw[popline, thick] (-3,-0.1) -- (3,-0.1);
        \draw[popline, thick] (-3,-0.6) -- (3,-0.6);
        \node[popmuted, anchor=north] at (0,-0.35) {Fente synaptique};
        % ACh libérée (normale)
        \foreach \i in {1,...,8} {
            \pgfmathsetmacro{\x}{-2+rand*4}
            \pgfmathsetmacro{\y}{-0.2+rand*0.3}
            \node[ach, scale=0.7] at (\x,\y) {ACh};
        }
        
        % Membrane musculaire
        \draw[muscle, thick, decorate, decoration={snake, amplitude=0.1cm, segment length=0.3cm}] (-3,-0.6) -- (3,-0.6);
        \node[muscle, anchor=north] at (0,-0.7) {Membrane musculaire};
        
        % Récepteurs AChR BLOQUÉS / DÉTRUITS
        \foreach \x in {-2.5, -1.5, -0.5, 0.5, 1.5, 2.5} {
            % Récepteur
            \draw[rach, fill=rach!30, thick] (\x,-0.75) rectangle ++(0.6, -0.4);
            % Auto-anticorps fixés (IgG)
            \draw[ab, fill=ab!30, thick] (\x+0.1,-0.75) -- (\x+0.25,-0.55) -- (\x+0.5,-0.55) -- (\x+0.35,-0.65) -- (\x+0.5,-0.75) -- (\x+0.25,-0.75) -- cycle; % Forme Y simplifiée
            \draw[ab, fill=ab!30, thick] (\x+0.35,-0.75) -- (\x+0.2,-0.55) -- (\x-0.05,-0.55) -- (\x+0.1,-0.65) -- (\x-0.05,-0.75) -- (\x+0.2,-0.75) -- cycle; % Second bras
            % Corps IgG
            \fill[ab] (\x+0.3,-0.65) circle (0.04cm);
        }
        \node[ab, anchor=north] at (0,-1.3) {Auto-anticorps anti-AChR (IgG1/IgG3) fixés};
        \node[rach, anchor=north] at (0,-1.55) {Récepteurs \textbf{bloqués}, \textbf{modulés} (internalisés) et \textbf{détruits} (complément)};
        
        % Flèche ÉCHEC
        \draw[poporange, ->, very thick, >=Stealth] (0,-1.7) -- (0,-2.1) node[midway, right, align=center] {Fixation ACh \textbf{réduite} $\rightarrow$ PPP \textbf{insuffisante} $\rightarrow$ \textbf{Pas de potentiel d'action} $\rightarrow$ \textbf{Fatigabilité musculaire}};
        
        % Légende anticorps
        \draw[ab, thick] (-3.5, -2.5) -- (-2.5, -2.5) node[right, ab] {Auto-anticorps IgG anti-AChR};
        \draw[rach, thick] (-3.5, -2.8) -- (-2.5, -2.8) node[right, rach] {Récepteur AChR (cible)};
        \draw[ach, thick] (-3.5, -3.1) -- (-2.5, -3.1) node[right, ach] {Acétylcholine};
    \end{scope}
    
    % Titres
    \node[align=center, font=\bfseries\normalsize] at (-4.5, 4.5) {Jonction neuromusculaire \textbf{NORMALE}};
    \node[align=center, font=\bfseries\normalsize] at (-4.5, 1.0) {Jonction neuromusculaire \textbf{MYASTHÉNIQUE}};
\end{tikzpicture}
\legende{Mécanisme de la myasthénie auto-immune. Haut : Jonction normale. L'acétylcholine (ACh, vert) libérée se fixe sur ses récepteurs (AChR, violet) $\rightarrow$ ouverture canaux Na$^+$ $\rightarrow$ contraction. Bas : Jonction myasthénique. Des auto-anticorps IgG (orange) se fixent sur les AChR. Conséquences : 1) \textbf{Blocage compétitif} de la fixation de l'ACh. 2) \textbf{Modulation antigénique} : internalisation et dégradation accélérée des AChR. 3) \textbf{Activation du complément} (C3b, C5b-9) $\rightarrow$ destruction de la membrane post-synaptique (plis disparus). Résultat : amplitude de la potentielle de plaque (PPP) réduite, souvent sous le seuil de déclenchement du potentiel d'action musculaire $\rightarrow$ faiblesse fatigable.}
\end{popfigure}

\begin{propriete}[Myasthénie auto-immune (Myasthénie grave)]
\begin{itemize}
    \item \textbf{Cible :} \textbf{Récepteur nicotinique à l'acétylcholine (AChR)} de la plaque motrice (jonction neuromusculaire). Parfois anti-MuSK (Muscle-specific Kinase) ou anti-LRP4.
    \item \textbf{Mécanisme (3 effets synergiques des auto-Ac IgG1/IgG3) :}
    \begin{enumerate}
        \item \textbf{Blocage compétitif :} L'auto-Ac empêche la fixation de l'ACh sur son site de liaison.
        \item \textbf{Modulation antigénique (cross-linking) :} L'auto-Ac lie 2 AChR $\rightarrow$ internalisation par endocytose $\rightarrow$ dégradation lysosomale $\rightarrow$ \textbf{diminution de la densité des AChR} (de $\sim 10^4$/$\mu$m$^2$ à $< 10^3$/$\mu$m$^2$).
        \item \textbf{Activation du complément (voie classique) :} Fixation de C1q sur la portion Fc des IgG $\rightarrow$ formation du Complexe d'Attaque Membranaire (MAC, C5b-9) $\rightarrow$ lyse de la membrane post-synaptique $\rightarrow$ \textbf{simplification des plis} (perte de surface).
    \end{enumerate}
    \item \textbf{Physiopathologie :} La Potentielle de Plaque Post-synaptique (PPP) a une amplitude réduite. Normalement $\gg$ seuil. En crise myasthénique, PPP $<$ seuil $\rightarrow$ pas de potentiel d'action musculaire $\rightarrow$ pas de contraction.
    \item \textbf{Signe majeur :} \textbf{Fatigabilité musculaire} (aggravation à l'effort, amélioration au repos). Atteinte fréquente : oculomoteurs (ptôsis, diplopie), bulaires (dysphagie, dysphonie), respiratoire (crise myasthénique = urgence vitale).
    \item \textbf{Diagnostic :} Test à la néostigmine (inhibiteur de l'acétylcholinestérase $\rightarrow$ augmentation [ACh] fente $\rightarrow$ amélioration transitoire), dosage auto-Ac anti-AChR (sensibilité 85-90 \%), électromyogramme (réponse décrémentaire à la stimulation répétitive bas fréquence 3 Hz).
    \item \textbf{Traitement :} Inhibiteurs de l'acétylcholinestérase (pyridostigmine), immunosuppresseurs (corticoïdes, azathioprine, mycophénolate), immunoglobulines IV / plasmaphérèse (crises), thymectomie (si thymome ou hyperplasie thymique, fréquente chez les jeunes).
\end{itemize}
\end{propriete}

\courssub{Interprétation de documents cliniques et biologiques}

\begin{methode}[Face à un cas clinique suspect d'auto-immunité : identifier la cible]
\begin{enumerate}
    \item \textbf{Analyser les symptômes :} Quel organe/fonction est touché ? (Glycémie $\rightarrow$ pancréas endocrine ; Peau blanche $\rightarrow$ mélanocytes ; Fatigabilité musculaire $\rightarrow$ jonction neuromusculaire).
    \item \textbf{Identifier le marqueur biologique spécifique :} Quel auto-anticorps ou quelle population lymphocytaire est dosée/observée ? (Anti-GAD/IA2, Anti-AChR, Anti-mélanocytes, Infiltrat LT8+).
    \item \textbf{Relier cible $\rightarrow$ mécanisme $\rightarrow$ symptôme :}
    \begin{itemize}
        \item Cible = Cellule sécrétrice (cellule $\beta$) $\rightarrow$ Destruction cellulaire (LT8+) $\rightarrow$ Déficit de sécrétion (insuline) $\rightarrow$ Trouble métabolique (hyperglycémie).
        \item Cible = Récepteur membranaire (AChR) $\rightarrow$ Blocage/Destruction récepteur (Auto-Ac + Complément) $\rightarrow$ Déficit de transduction signal (dépolarisation) $\rightarrow$ Trouble fonctionnel (faiblesse).
        \item Cible = Cellule pigmentaire (mélanocyte) $\rightarrow$ Destruction (Auto-Ac + LT8+) $\rightarrow$ Perte fonction (mélanogénèse) $\rightarrow$ Signe visible (dépigmentation).
    \end{itemize}
    \item \textbf{Proposer le traitement logique :} Substitution (insuline), Symptomatique + Immunosuppression (myasthénie), Immunosuppression locale/systémique + Photothérapie (vitiligo).
\end{enumerate}
\end{methode}

\begin{exoresolu}[Interprétation d'un dossier myasthénie]
\textbf{Énoncé :} Une patiente de 32 ans consulte pour une ptôsis bilatérale s'aggravant en fin de journée et une diplopie. L'examen neurologique montre une fatigabilité des muscles oculomoteurs et des membres supérieurs à la répétition. Le test à la néostigmine (1 mg IV) est positif (amélioration nette à 3 min). Le dosage des auto-anticorps anti-récepteurs à l'acétylcholine (anti-AChR) revient à 12 nmol/L (valeurs de référence $< 0,5$ nmol/L). L'EMG de stimulation répétitive (3 Hz) du nerf cubital montre une décrémentation de 18 \% de l'amplitude de la 5ème réponse par rapport à la 1ère. Un scanner thoracique révèle une hyperplasie thymique.
\\
\textbf{Solution détaillée :}
\begin{enumerate}
    \item \textbf{Diagnostic :} Myasthénie auto-immune généralisée (forme oculobulaire + membres) avec hyperplasie thymique associée. Arguments : tableau clinique typique (fatigabilité), test pharmacologique positif, auto-anticorps anti-AChR fortement positifs (spécificité $> 95\%$), signe électrophysiologique caractéristique (décrémentation $> 10\%$).
    \item \textbf{Mécanisme lésionnel :} Les auto-Ac anti-AChR (IgG1/3) fixés sur les AChR de la plaque motrice entraînent : (1) blocage compétitif de la fixation de l'ACh ; (2) cross-linking $\rightarrow$ internalisation/dégradation des AChR (modulation antigénique) ; (3) activation du complément $\rightarrow$ destruction des plis post-synaptiques. $\rightarrow$ Baisse de la densité des AChR $\rightarrow$ Amplitude de la PPP réduite $\rightarrow$ Échec de la transmission neuromusculaire à l'effort (fatigabilité).
    \item \textbf{Rôle du thymus :} Le thymus est le site d'initiation/entretien de la réponse auto-immune (présence de myoïdes exprimant l'AChR, défaut de sélection négative, hyperplasie germinative). La thymectomie est indiquée ici (patiente jeune, hyperplasie, anticorps positifs) pour améliorer la rémission.
    \item \textbf{Traitement proposé :} Pyridostigmine (inhibiteur AChE) en 1ère intention symptomatique. Corticothérapie (prednisone) + immunosuppresseur d'épargne (azathioprine/mycophénolate) pour induire une rémission. Thymectomie programmée. Éducation : signes de crise myasthénique (dyspnée, dysphagie sévère) = urgence réanimation (plasmaphérèse/IgIV).
\end{enumerate}
\end{exoresolu}

\courssub{Traitements : principes, limites et effets secondaires}

\begin{propriete}[Stratégies thérapeutiques dans les maladies auto-immunes]
\begin{center}
\begin{tabularx}{\linewidth}{|l|X|X|X|}\hline
\rowcolor{popblueL}
\textbf{Stratégie} & \textbf{Principe / Exemples} & \textbf{Indications types} & \textbf{Limites / Effets secondaires majeurs} \\ \hline
\textbf{Substitution} & Remplacer la fonction perdue par l'organe détruit. \newline Ex : \textbf{Insuline} (diabète T1), Thyroxine (thyroïdite de Hashimoto), Facteur VIII (hémophilie acquise auto-immune). & Maladies à \textbf{destruction cellulaire} avec déficit hormonal/enzymatique. & Ne traite pas la cause immune. Contraintes : injections, surveillance (glycémie), coût, chaîne du froid (insuline). \\ \hline
\textbf{Immunosuppression / Immunomodulation globales} & \textbf{Corticoïdes} (prednisone) : anti-inflammatoire puissant, lymphopénie. \newline \textbf{Immunosuppresseurs classiques :} Azathioprine, Cyclophosphamide, Mycophénolate mofétil (inhibent prolifération lymphocytaire). \newline \textbf{Inhibiteurs calcineurine :} Ciclosporine, Tacrolimus (blocage activation LT). & Poussées aiguës, maladies systémiques (lupus, polyarthrite), myasthénie, greffes. & \textbf{Infections opportunistes}, toxicité rénale/hépatique, hypertension, diabète induit, ostéoporose, risque cancérogène (lymphomes), tératogénicité. Nécessite surveillance biologique stricte (NFS, fonction rénale, taux résiduel). \\ \hline
\textbf{Biothérapies (ciblées)} & \textbf{Anti-TNF$\alpha$} (Infliximab, Adalimumab, Etanercept). \newline \textbf{Anti-CD20} (Rituximab : déplétion LB). \newline \textbf{Anti-IL-6R} (Tocilizumab). \newline \textbf{Anti-CTLA4 / PD-1} (cancer mais risque auto-immunité paradoxale). & Polyarthrite rhumatoïde, spondylarthrite, lupus réfractaire, myasthénie anti-MuSK (Rituximab). & Coût très élevé (accès limité pays en développement), risque infections (réactivation tuberculose pour anti-TNF), réactions perfusion, anticorps anti-médicament (perte efficacité). \\ \hline
\textbf{Techniques d'épuration} & \textbf{Plasmaphérèse / Échange plasmatique :} Retrait auto-Ac circulants. \newline \textbf{Immunoglobulines IV (IgIV) :} Modulation FcR, inhibition complément, neutralisation auto-Ac, régulation $T_{reg}$. & Crises aiguës sévères : crise myasthénique, syndrome de Guillain-Barré, poussée lupus néphritique, purpura thrombopénique. & Effet transitoire (semaines), invasif (cathéter veineux central), coût, risques hémodynamiques/allergiques. \\ \hline
\textbf{Restauration de la tolérance (Recherche/Expérimental)} & Vaccination à l'auto-antigène (peptides), thérapie cellulaire ($T_{reg}$ autologues expandis), anticorps monoclonaux co-stimulateurs (Abatacept : CTLA4-Ig). & Essais cliniques (diabète T1 stade précoce, polyarthrite). & Pas encore standard. Risque d'exacerber l'auto-immunité si mal dosé. \\ \hline
\end{tabularx}
\end{center}
\end{propriete}

\begin{attention}[Surveillance au long cours : l'exemple du diabète type 1 au Cameroun]
Un jeune diabétique type 1 suivi à l'Hôpital Général de Douala ou au Centre des Maladies Métaboliques de Yaoundé doit bénéficier d'une surveillance trimestrielle : HbA1c (cible $< 7,5\%$ chez l'adolescent/adulte jeune), fonction rénale (microalbuminurie, créatininémie), fond d'œil (rétinopathie), pieds (neuropathie/artériopathie), profil lipidique. L'éducation thérapeutique (comptage glucides, ajustement doses insuline, gestion hypoglycémie/hyperglycémie/cétoacidose) est le \textbf{seul levier} pour prévenir les complications chroniques invalidantes. L'accès aux capteurs glucose en continu (FreeStyle Libre, Dexcom) reste limité mais change la donne là où il est disponible.
\end{attention}

\begin{savaistu}[Le thymus, organe central de l'auto-immunité ?]
Dans la myasthénie, 65 à 85 \% des patients ont une \textbf{hyperplasie thymique} (follicules germinatifs actifs) et 10-15 \% un \textbf{thymome} (tumeur thymique). Le thymus exprime l'AChR sur des cellules myoïdes. C'est un site de \textbf{sélection négative défaillante} et/ou d'activation aberrante de LT auto-réactifs. La \textbf{thymectomie} (chirurgie vidéo-assistée ou sternotomie) induit une rémission complète ou une amélioration significative chez $\sim 30-40\%$ des patients anticorps+ sans thymome, et améliore la réponse aux immunosuppresseurs. C'est l'une des rares maladies auto-immunes où l'ablation d'un organe lymphoïde primaire est thérapeutique.
\end{savaistu}

\begin{exercice}[Cas clinique : Diabète type 1 révélé par une cétoacidose]
Un adolescent de 14 ans, originaire de Bafoussam, est admis en urgence pour coma. L'interrogatoire rétrospectif retrouve : polyuro-polydipsie depuis 3 semaines, amaigrissement de 5 kg, fatigue. À l'admission : conscience altérée, déshydratation cutanée majeure, tachycardie 110/min, TA 90/60 mmHg, respiration de Kussmaul (ample, rapide). Glycémie capillaire : 4,80 g/L (26,6 mmol/L). Gaz du sang artériel : pH 7,18 ; $p\ce{CO2}$ 25 mmHg ; $\ce{HCO3-}$ 10 mmol/L ; Lactates 3 mmol/L. Urines : glycosurie +++, cétonurie +++. Le dosage des auto-anticorps anti-GAD65 revient positif à 120 UI/mL ($< 5$ UI/mL). Le peptide C (C-peptide) basal est à 0,15 nmol/L (valeurs de référence jeun : 0,3 - 0,9 nmol/L).
\\
\textbf{Questions :}
\begin{enumerate}
    \item Qualifier le trouble de la conscience et expliquer son mécanisme physiopathologique (lien glycémie/pH/osmolarité).
    \item Interpréter les résultats biologiques (glycémie, gaz du sang, C-peptide, auto-Ac) pour confirmer le diagnostic étiologique.
    \item Expliquer pourquoi le dosage du C-peptide est plus pertinent que celui de l'insuline chez ce patient non encore traité.
    \item Proposer la prise en charge initiale d'urgence (premières 6-12h) et le traitement de fond.
\end{enumerate}
\end{exercice}

\begin{corrige}[Exercice n]
\begin{enumerate}
    \item \textbf{Coma cétoacidosique diabétique.} Mécanisme : Hyperglycémie sévère $\rightarrow$ hyperosmolarité plasmatique $\rightarrow$ déshydratation intracellulaire (neurones) + acidose métabolique (accumulation corps cétoniques $\beta$-hydroxybutyrate, acétoacétate) $\rightarrow$ dépression du SNC (coma). Respiration de Kussmaul = compensation respiratoire (hyperventilation $\rightarrow$ baisse $p\ce{CO2}$) pour tamponner l'acidose.
    \item \textbf{Diagnostic : Diabète sucré de type 1 révélé par une cétoacidose.}
    \begin{itemize}
        \item Glycémie 4,80 g/L $>>$ 1,26 g/L : diabète avéré.
        \item pH 7,18 + $\ce{HCO3-}$ 10 mmol/L + cétonurie +++ + glycémie élevée : \textbf{Cétoacidose diabétique (CAD) modérée-sévère} (pH $< 7,3$, $\ce{HCO3-} < 15$).
        \item Auto-Ac anti-GAD65 ++ : \textbf{Étiologie auto-immune} confirmée.
        \item C-peptide très bas (0,15 nmol/L) : \textbf{Réserve sécrétoire insulinique quasi nulle} (destruction $> 90\%$ cellules $\beta$).
    \end{itemize}
    \item \textbf{C-peptide vs Insuline :} L'insuline exogène (traitement) n'a pas encore été administrée, mais l'insuline endogène a une demi-vie très courte ($\sim 5-6$ min) et est extraite hépatiquement (50 \%). Le C-peptide (peptide de liaison de la proinsuline) est sécrété en \textbf{quantité équimolaire} à l'insuline, n'est \textbf{pas extrait par le foie}, a une demi-vie plus longue ($\sim 20-30$ min) et n'est \textbf{pas présent dans les préparations d'insuline thérapeutique}. Son dosage reflète fidèlement la sécrétion endocytaire résiduelle, indispensable pour distinguer type 1 (C-peptide bas) de type 2 (C-peptide normal/élevé) ou diabète secondaire.
    \item \textbf{Prise en charge urgence (premières 6-12h) :}
    \begin{enumerate}
        \item \textbf{Réhydratation IV :} Serum salé isotonique (\ce{NaCl} 0,9\%) 10-20 mL/kg/h (1er litre), puis adapter selon hémodynamique/hémoglobine/hématocrite. But : restaurer volémie, clairance rénale glucose/cétones.
        \item \textbf{Insulinothérapie IV :} Insuline rapide (Actrapid\textsuperscript{\textregistered} ou équivalent) en pompe électrique : bolus 0,1 UI/kg puis 0,1 UI/kg/h continu. Adapter pour faire baisser glycémie de $\sim 0,5-1,0$ g/L/h (éviter chute brutale $\rightarrow$ œdème cérébral). Arrêter perfusion insuline IV quand pH $> 7,3$ et cétonémie normalisée $\rightarrow$ relais SC (basal-bolus).
        \item \textbf{Correction potassique :} CAD = déplétion potassique totale (pertes urinaires + sortie cellulaire par acidose). L'insuline + correction acidose $\rightarrow$ rentrée K$^+$ intracellulaire $\rightarrow$ \textbf{hypokaliémie sévère risquée (arythmies)}. Surveiller kaliémie H1-H2. Supplémenter KCl dès que diurèse reprise et kaliémie $< 5,0$ mmol/L (ex: 20-40 mmol/L dans perfusions).
        \item \textbf{Surveillance continue :} Glycémie H1, Ionogramme/Gaz du sang H2-H4, Neurologie (risque œdème cérébral : céphalées, agitation, troubles pupillaires).
    \end{enumerate}
    \textbf{Traitement de fond :} Éducation thérapeutique (patient + famille), Schéma basal-bolus (insuline basale 1-2/j + rapide 3-4/j aux repas), Autosurveillance glycémique (4-6/j), Objectif HbA1c $< 7,5\%$, Suivi complications (néphro/rétrino/neuropathie) annuel, Vaccination grippe/pneumocoque/COVID à jour. Soutien psychologique (acceptation maladie chronique).
\end{enumerate}
\end{corrige}

\begin{aretenir}[Synthèse : Les maladies auto-immunes]
\begin{itemize}
    \item \textbf{Définition :} Atteinte de tissus propres par auto-anticorps et/ou lymphocytes T auto-réactifs $\rightarrow$ rupture de la tolérance au soi.
    \item \textbf{Triade étiologique :} Génétique (HLA) + Environnement (mimétisme, infections) + Défaillance régulation ($T_{reg}$).
    \item \textbf{3 modèles types au programme :}
    \begin{itemize}
        \item \textbf{Diabète type 1 :} Cible = Cellule $\beta$ pancréatique. Effecteur = LT8+ cytotoxiques. Conséquence = Déficit insuline absolu. Traitement = Substitution insuline (vitale).
        \item \textbf{Vitiligo :} Cible = Mélanocyte. Effecteur = Auto-Ac + LT8+. Conséquence = Dépigmentation. Traitement = Immunosuppression locale + UVB.
        \item \textbf{Myasthénie :} Cible = Récepteur AChR (jonction neuromusculaire). Effecteur = Auto-Ac IgG (blocage, modulation, complément). Conséquence = Fatigabilité musculaire. Traitement = Anti-AChE + Immunosuppresseurs $\pm$ Thymectomie.
    \end{itemize}
    \item \textbf{Diagnostic :} Association Clinique + Auto-anticorps spécifiques + Histologie (infiltrat) + Tests fonctionnels (EMG, test pharmacologique).
    \item \textbf{Traitements :} Substitution (si déficit), Immunosuppression (corticoïdes, conventionnels, biothérapies), Épurations (plasmaphérèse, IgIV). Surveillance infectieuse et métabolique indispensable.
\end{itemize}
\end{aretenir}

\coursec{Le VIH/Sida : mécanismes, prévention et solidarité}

Après avoir étudié les dysfonctionnements où le système immunitaire s'attaque à lui-même (maladies auto-immunes), nous abordons maintenant une situation où un agent pathogène extérieur, le \cle{VIH} (Virus de l'Immunodéficience Humaine), détruit spécifiquement le chef d'orchestre de la réponse immunitaire spécifique : le lymphocyte T4 (LT4). Cette destruction progressive conduit, en l'absence de traitement, au stade \cle{Sida} (Syndrome d'Immunodéficience Acquise), ouvrant la porte aux infections opportunistes. Comprendre les mécanismes moléculaires de ce virus, interpréter l'évolution biologique de l'infection et maîtriser les moyens de prévention sont des compétences essentielles, tant sur le plan scientifique que citoyen, notamment dans notre contexte camerounais où la riposte au VIH/Sida reste une priorité de santé publique.

\courssub{Le VIH : un rétrovirus ciblant le lymphocyte T4}

\begin{definition}[VIH et Rétrovirus]
Le \textbf{VIH} (Virus de l'Immunodéficience Humaine) est un \cle{rétrovirus} à enveloppe, dont le matériel génétique est constitué de deux brins d'\textbf{ARN} simple brin (+). Sa particularité majeure est de posséder une enzyme spécifique, la \cle{transcriptase inverse}, qui permet de synthétiser de l'ADN à partir de son ARN génomique, inversant ainsi le flux classique de l'information génétique (ADN $\to$ ARN $\to$ Protéine). Deux types existent : VIH-1 (le plus répandu mondialement et au Cameroun) et VIH-2 (plus fréquent en Afrique de l'Ouest, moins pathogène).
\end{definition}

La cible privilégiée du VIH est le \textbf{lymphocyte T4} (ou lymphocyte T helper CD4+), reconnaissable à la présence à sa surface du récepteur \cle{CD4}. Ce récepteur sert de « porte d'entrée » principale au virus. La glycoprotéine virale \textbf{gp120}, portée par l'enveloppe, se lie avec une forte affinité au récepteur CD4, puis à un co-récepteur (généralement CCR5 ou CXCR4), déclenchant la fusion de l'enveloppe virale avec la membrane cellulaire.

\begin{experience}[Mise en évidence du tropisme du VIH pour les cellules CD4+]
\textbf{Observation :} En culture cellulaire, le VIH infecte efficacement les lignées de lymphocytes T exprimant le CD4 (ex: lignée SupT1, cellules PBMC activées), mais n'infecte pas les lignées dépourvues de CD4 (ex: lignées de lymphocytes B, cellules épithéliales HeLa), sauf si on les transfecte artificiellement pour exprimer le CD4.
\textbf{Interprétation :} Le récepteur CD4 est nécessaire et suffisant (avec un co-récepteur) pour l'entrée du virus. Cela explique la sélectivité de la destruction des LT4 \emph{in vivo}.
\textbf{Conclusion :} La déplétion des LT4 est une conséquence directe du cycle de réplication viral dans ces cellules.
\end{experience}

\courssub{Le cycle de multiplication du VIH dans le lymphocyte T4}

Le cycle de réplication suit les étapes classiques des rétrovirus, avec des cibles thérapeutiques majeures pour les antirétroviraux (ARV).

\begin{popfigure}
\begin{tikzpicture}[
    scale=0.9, transform shape,
    node distance=1.2cm and 1.5cm,
    box/.style={rectangle, draw=popdark, thick, fill=popcream, rounded corners, minimum width=2.5cm, minimum height=1cm, align=center, font=\small},
    arrow/.style={-Stealth, thick, popdark},
    virus/.style={circle, draw=popred, thick, fill=poppinkL, minimum size=0.8cm},
    nucleus/.style={ellipse, draw=popblue, thick, fill=popblueL, minimum width=2.5cm, minimum height=1.5cm},
    enzyme/.style={diamond, draw=poporange, thick, fill=poporangeL, aspect=2, inner sep=2pt, font=\tiny}
]
% Cell membrane
\draw[popdark, very thick] (-5,-3) -- (5,-3) node[midway, below=0.2cm, font=\small] {Membrane plasmique (CD4 + Co-récepteur)};
% Virus outside
\node[virus, label=above:\textcolor{popred}{Virion VIH}] (virion) at (-4, -1) {};
% 1 Fixation/Entree
\node[box, right=of virion, align=center] (fix){1. Fixation gp120/CD4\\+ Co-récepteur $\to$ Fusion};
\draw[arrow] (virion) -- (fix);
% 2 Entree capside
\node[box, below=of fix, align=center] (entree){2. Pénétration \&\\Décapsidage\\Libération ARN + Transcriptase inverse};
\draw[arrow] (fix) -- (entree);
% 3 Transcription inverse
\node[enzyme, right=of entree, align=center] (ti){Transcriptase\\ inverse};
\node[box, right=of ti, align=center] (rt){3. Transcription inverse\\ARN viral $\xrightarrow{\text{TI}}$ ADN double brin (ADNv)};
\draw[arrow] (entree) -- (ti);
\draw[arrow] (ti) -- (rt);
% 4 Integration
\node[nucleus, above right=1.5cm and 1cm of rt, align=center] (nuc){Noyau\\ADN cellulaire};
\node[box, above=of rt, align=center] (int){4. Intégration\\ADNv $\xrightarrow{\text{Intégrase}}$ ADN hôte (Provirus)};
\draw[arrow] (rt) -- (int);
\draw[arrow, dashed] (int) -- (nuc) node[midway, left, font=\tiny, text=poppurple] {Transport actif};
% 5 Transcription/Traduction
\node[box, right=of nuc, align=center] (trans){5. Transcription/Traduction\\ARNm viraux $\to$ Protéines virales\\ (Polyprotéines Gag/Pol/Env)};
\draw[arrow] (nuc) -- (trans);
% 6 Assemblage/Budding
\node[box, below=of trans, align=center] (ass){6. Assemblage \& Bourgeonnement\\Nouveaux virions immatures};
\draw[arrow] (trans) -- (ass);
\node[box, below=of ass, align=center] (mat){7. Maturation\\Protéase $\to$ Virions infectieux};
\draw[arrow] (ass) -- (mat);
% Release
\node[virus, right=of mat, label=above:\textcolor{popred}{Virions fils}] (virion2) {};
\draw[arrow] (mat) -- (virion2);
% ARV targets
\node[draw=popgreen, thick, fill=popgreenL, rounded corners, align=left, font=\tiny, anchor=north west] at (6, 2) {
\textbf{Cibles des ARV :}\\
$\bullet$ Inhibiteurs d'entrée (Maraviroc, Enfuvirtide)\\
$\bullet$ Inhibiteurs Transcriptase Inverse (TDF, 3TC, EFV...)\\
$\bullet$ Inhibiteurs Intégrase (DTG, RAL)\\
$\bullet$ Inhibiteurs Protéase (LPV/r, DRV/r)
};
\end{tikzpicture}
\legende{Cycle de multiplication du VIH dans un lymphocyte T4. Le virus fixe le récepteur CD4 et un co-récepteur, fusionne, libère son ARN et sa transcriptase inverse. L'ADN viral (provirus) s'intègre au génome hôte. La machinerie cellulaire produit les composants viraux qui s'assemblent et bourgeonnent. La maturation par la protéase virale rend les nouveaux virions infectieux. Les classes d'antirétroviraux (ARV) bloquent chacune une étape clé.}
\end{popfigure}

\begin{aretenir}[Étapes clés et cibles thérapeutiques]
\begin{enumerate}
    \item \textbf{Fixation / Entrée :} gp120 $\leftrightarrow$ CD4 + Co-récepteur (CCR5/CXCR4) $\to$ Fusion. \textit{Cible : Inhibiteurs d'entrée.}
    \item \textbf{Transcription inverse :} ARN viral $\xrightarrow{\text{Transcriptase inverse}}$ ADN double brin. \textit{Cible : INRTI (Ténofovir, Lamivudine) et INNRTI (Éfavirenz, Névirapine).}
    \item \textbf{Intégration :} ADN viral $\xrightarrow{\text{Intégrase}}$ ADN hôte (Provirus). \textit{Cible : Inhibiteurs d'intégrase (Dolutégravir -- recommandé en 1ère ligne au Cameroun).}
    \item \textbf{Expression / Assemblage / Maturation :} Transcription $\to$ Traduction $\to$ Assemblage $\xrightarrow{\text{Protéase}}$ Virions matures. \textit{Cible : Inhibiteurs de protéase (Lopinavir/r, Darunavir/r).}
\end{enumerate}
\end{aretenir}

\courssub{Conséquences immunologiques : de l'infection au Sida}

Le rôle central des LT4 (activation des LT8 cytotoxiques, aide aux LB pour la réponse humorale, activation des macrophages via IFN-$\gamma$, régulation via cytokines) explique pourquoi leur destruction entraîne un effondrement global de l'immunité spécifique.

\begin{propriete}[Évolution de l'infection VIH non traitée en trois phases]
L'infection par le VIH évolue classiquement en trois phases cliniques et biologiques distinctes, caractérisées par l'évolution opposée du \textbf{taux de LT4} (cellules/$\mu$L) et de la \cle{charge virale} (copies ARN/mL) :
\begin{enumerate}
    \item \textbf{Primo-infection (semaines 1--12) :} Réplication virale explosive (charge virale $> 10^6$ copies/mL), chute brutale et transitoire des LT4 ($< 500$/$\mu$L), syndrome pseudo-grippal ou mononucléosique (fièvre, adénopathies, éruption cutanée). Séroconversion (apparition anticorps anti-VIH) vers 3--6 semaines.
    \item \textbf{Latence clinique (années 1--10+) :} Équilibre précaire. Charge virale stabilisée à un « point de consigne » (ex: $10^4$--$10^5$ copies/mL). Déclin lent et continu des LT4 ($\approx 50$--$80$ cellules/$\mu$L/an). Patient asymptomatique mais contagieux.
    \item \textbf{Stade Sida :} Effondrement des LT4 $< 200$/$\mu$L (seuil de risque majeur). Charge virale en remontée. Survenue d'\textbf{infections opportunistes} (tuberculose, pneumocystose, toxoplasmose, méningites cryptococciques, cytomégalovirus) et de cancers (sarcome de Kaposi, lymphomes). Sans traitement, décès en 1--2 ans.
\end{enumerate}
\end{propriete}

\begin{popfigure}
\begin{tikzpicture}[trim axis left, trim axis right]
% Axis 1: Viral Load (Left, Log scale)
\begin{axis}[
    width=\linewidth, height=10cm,
    xlabel={Temps écoulé depuis l'infection (années)},
    ylabel={Charge virale (log$_{10}$ copies/mL)},
    ymode=log,
    log basis y=10,
    ymin=1, ymax=1000000,
    xmin=0, xmax=12,
    axis y line*=left,
    axis x line=bottom,
    tick align=outside,
    grid=major,
    grid style={popline},
    ytick={1,10,100,1000,10000,100000,1000000},
    yticklabels={$10^0$,$10^1$,$10^2$,$10^3$,$10^4$,$10^5$,$10^6$},
    legend style={at={(0.98,0.98)}, anchor=north east, font=\small, fill=popcream, draw=popdark},
    clip=false
]
\addplot[popred, very thick, smooth, mark=none] coordinates {
    (0, 100) (0.1, 500000) (0.3, 800000) (0.5, 300000) (1, 50000) (3, 30000) (5, 50000) (8, 80000) (10, 200000) (12, 500000)
};
\addlegendentry{Charge virale (copies/mL)}

\node[popred, font=\small, rotate=90, anchor=south] at (axis cs:0.2, 1000000) {Pic primo-infection};
\node[popred, font=\small, anchor=north] at (axis cs:5, 30000) {Point de consigne (Latence)};
\node[popred, font=\small, anchor=west] at (axis cs:10, 200000) {Remontée (Sida)};
\end{axis}

% Axis 2: CD4 Count (Right, Linear scale)
\begin{axis}[
    width=\linewidth, height=10cm,
    ylabel={Taux de LT4 (cellules/$\mu$L)},
    ymin=0, ymax=1300,
    xmin=0, xmax=12,
    axis y line*=right,
    axis x line=none,
    tick align=outside,
    ytick={0,200,400,600,800,1000,1200},
    legend style={at={(0.98,0.85)}, anchor=north east, font=\small, fill=popcream, draw=popdark},
    clip=false
]
\addplot[popblue, very thick, smooth, mark=none] coordinates {
    (0, 1000) (0.1, 600) (0.5, 500) (1, 700) (3, 600) (5, 450) (8, 250) (10, 150) (12, 50)
};
\addlegendentry{Taux de LT4 (cellules/$\mu$L)}

\draw[poporange, dashed, thick] (axis cs:0,200) -- (axis cs:12,200) node[pos=0.5, above, font=\tiny, poporange] {Seuil Sida (200/$\mu$L)};
\draw[popgreen, dashed, thick] (axis cs:0,500) -- (axis cs:12,500) node[pos=0.5, below, font=\tiny, popgreen] {Seuil normal bas (500/$\mu$L)};

\fill[popred!10] (axis cs:0,0) rectangle (axis cs:0.5,1300);
\fill[popgold!10] (axis cs:0.5,0) rectangle (axis cs:8,1300);
\fill[poppurple!10] (axis cs:8,0) rectangle (axis cs:12,1300);

\node[font=\small\bfseries, popdark] at (axis cs:0.25, 1200) {Primo-infection};
\node[font=\small\bfseries, popdark] at (axis cs:4, 1200) {Latence clinique};
\node[font=\small\bfseries, popdark] at (axis cs:10, 1200) {Stade Sida};
\end{axis}
\end{tikzpicture}
\legende{Évolution conjointe de la charge virale (courbe rouge, échelle logarithmique à gauche) et du taux de lymphocytes T4 (courbe bleue, échelle linéaire à droite) au cours des trois phases de l'infection VIH non traitée. La primo-infection se traduit par un pic viral massif et une chute aiguë des LT4. La latence clinique montre un équilibre instable (point de consigne viral, déclin lent des LT4). Le stade Sida correspond à l'effondrement des LT4 $< 200/\mu$L et à la remontée de la charge virale.}
\end{popfigure}

\begin{attention}[Séropositivité $\neq$ Sida]
\textbf{Être séropositif} signifie que des anticorps anti-VIH sont détectés dans le sang : l'infection est établie. \textbf{Avoir le Sida} signifie que l'infection a progressé jusqu'au stade terminal (LT4 $< 200/\mu$L ou infection opportuniste définissante). \textbf{Une personne séropositive sous traitement efficace (ARV) ne développe pas le Sida}, maintient un taux de LT4 normal et une charge virale indétectable. Ne jamais confondre le statut sérologique et le stade clinique.
\end{attention}

\courssub{Modes de transmission et prévention}

\begin{definition}[Modes de transmission du VIH]
Le VIH se transmet par trois voies principales, nécessitant un contact avec des liquides biologiques contaminés (sang, sperme, sécrétions vaginales, lait maternel) en quantité suffisante :
\begin{enumerate}
    \item \textbf{Via sexuelle (majoritaire au Cameroun, $\approx 90\%$ des cas) :} Rapports vaginaux, anaux ou oraux non protégés avec une personne infectée (charge virale détectable). Le risque est plus élevé en cas d'IST non traitées (lésions génitales).
    \item \textbf{Via sanguine :} Transfusion de sang non dépisté (rare aujourd'hui grâce au dépistage systématique), partage de seringues/aiguilles/objets tranchants (rasoirs, lames de scarification, matériel de tatouage/piercing non stérilisé), accidents d'exposition au sang (AES) chez le personnel soignant.
    \item \textbf{Mère-enfant (transmission verticale) :} Pendant la grossesse (transplacentaire), l'accouchement (contact sang/sécrétions), ou l'allaitement (lait maternel). Sans intervention, risque $15$--$45\%$ ; avec ARV + accouchement encadré + substituts laitiers, risque $< 2\%$ (élimination validée au Cameroun pour plusieurs sites).
\end{enumerate}
\end{definition}

\begin{attention}[Idées reçues : gestes qui NE transmettent PAS le VIH]
Le VIH est un virus fragile hors de l'organisme (inactivé par la chaleur, le savon, l'eau de Javel, le séchage). \textbf{Il ne se transmet PAS par :}
\begin{itemize}
    \item Les contacts quotidiens : poignées de main, embrassades (baisers sociaux), câlins, partage des repas, des verres, des couverts, des toilettes, des draps, des vêtements.
    \item Les piqûres d'insectes (moustiques, tiques) : le virus ne se réplique pas dans le moustique et la quantité de sang résiduaire sur la trompe est insuffisante.
    \item La salive, les larmes, la sueur, l'urine, les selles (sauf s'ils sont visiblement souillés de sang en grande quantité).
    \item L'air, l'eau, les aliments.
\end{itemize}
Lutter contre ces fausses croyances est la base de la lutte contre la stigmatisation.
\end{attention}

\begin{methode}[Stratégies de prévention combinée (Approche « Combinaison »)]
La prévention efficace repose sur l'association de plusieurs leviers adaptés à la situation de chacun :
\begin{enumerate}
    \item \textbf{Abstinence / Différer les premiers rapports :} Protection à 100\% (seule méthode absolue).
    \item \textbf{Fidélité mutuelle au sein du couple :} Après dépistage conjoint négatif.
    \item \textbf{Préservatif (masculin ou féminin) :} Utilisé correctement à \textbf{chaque} rapport (vaginal, anal, oral). Seule méthode protégeant aussi des IST et grossesses non désirées. Disponible gratuitement dans les centres de santé et CEC (Centres d'Écoute et de Conseil) au Cameroun.
    \item \textbf{Dépistage volontaire (DV) / Dépistage Proposé par le Fournisseur (DPF) :} Connaître son statut sérologique (et celui de son partenaire) est la porte d'entrée de la prévention et de la prise en charge. Test rapide (résultat $< 30$ min) ou ELISA (labo).
    \item \textbf{PrEP (Prophylaxie Pré-Exposition) :} Prise d'ARV (TDF/3TC ou TDF/FTC) par une personne séronégative à haut risque (couples sérodifférents, travailleurs du sexe, HSH). Efficacité $> 95\%$ si observance parfaite. Disponible au Cameroun dans certains sites pilotes.
    \item \textbf{PEP (Prophylaxie Post-Exposition) :} Traitement d'urgence (28 jours) à démarrer \textbf{le plus vite possible, $< 48$h (idéalement $< 4$h)} après une exposition à risque (viol, accident d'aiguille, rapport non protégé avec partenaire de statut inconnu/positif non traité).
    \item \textbf{Traitement comme prévention (TasP / I=I) :} Personne vivant avec le VIH (PVVIH) sous ARV efficace $\to$ Charge virale indétectable ($< 50$ copies/mL) $\to$ \textbf{Risque de transmission sexuelle nul (Indétectable = Intransmissible)}.
    \item \textbf{Prévention Mère-Enfant (PTME) :} Dépistage systématique à la 1ère CPN, ARV à la mère (option B+ : à vie), accouchement sécurisé, prophylaxie ARV au nouveau-né (NVP/AZT 6 semaines), conseil sur l'alimentation du nourrisson (lait maternel exclusif 6 mois sous ARV maternel OU lait de substitution si conditions AFASS réunies).
    \item \textbf{Sécurité transfusionnelle \& Précautions standard :} Dépistage obligatoire de chaque don de sang, stérilisation matériel médical, usage unique des seringues, port de gants.
\end{enumerate}
\end{methode}

\courssub{Diagnostic : du dépistage à la confirmation}

\begin{experience}[Principe du test de dépistage rapide (TDR) / ELISA]
\textbf{Document :} Notice d'un test rapide de détection des anticorps anti-VIH-1/2 (ex: Determine, SD Bioline, First Response utilisés au Cameroun).
\textbf{Principe :} Immunochromatographie sur bandelette. L'échantillon (sang total, sérum, plasma) migre. Si des anticorps anti-VIH sont présents, ils se lient à l'antigène viral (gp41, gp120, gp36) conjugué à des particules d'or (couleur rouge/rose) $\to$ complexe Ag-Ac. Ce complexe est capturé par des anticorps anti-humains immobilisés sur la bande (ligne T = Test). Un excès de conjugué est capturé en ligne C (Contrôle) validant le test.
\textbf{Interprétation :}
\begin{itemize}
    \item Ligne C seule visible $\to$ \textbf{Négatif} (pas d'anticorps détectés).
    \item Lignes C et T visibles $\to$ \textbf{Positif} (anticorps détectés).
    \item Ligne C absente $\to$ \textbf{Invalide} (refaire le test).
\end{itemize}
\textbf{Confirmation :} Tout test rapide positif doit être confirmé par un deuxième test rapide de marque différente (algorithme national) ou par Western Blot / Test de confirmation en laboratoire (spécificité $\approx 100\%$).
\end{experience}

\begin{definition}[Fenêtre sérologique]
Période entre la contamination et l'apparition des anticorps détectables par les tests sérologiques standards. \textbf{Durée moyenne : 3 à 6 semaines (jusqu'à 3 mois dans de rares cas).} Pendant cette période, le test est \textbf{faux négatif} alors que la personne est contaminée et très contagieuse (charge virale très élevée). D'où l'importance de \textbf{répéter le test 3 mois après la dernière exposition à risque} pour lever le doute. La recherche d'antigène p24 (test combiné Ag/Ac) ou de l'ARN viral (PCR) raccourcit cette fenêtre à $\approx 10$--$15$ jours.
\end{definition}

\begin{exemplebox}[Exemple chiffré : Dépistage au Cameroun]
Au Cameroun, le dépistage du VIH est \textbf{gratuit} dans tous les centres de santé publics et agréés (hôpitaux, centres de santé, CEC, certains laboratoires privés conventionnés).
\begin{itemize}
    \item \textbf{Test rapide :} Résultat en $15$--$30$ minutes sur sang total (prélèvement au bout du doigt).
    \item \textbf{Algorithme national :} Test A (ex: Determine) $\to$ Si positif, Test B (ex: SD Bioline) $\to$ Si positif, Test C (ex: First Response ou Stat-Pak) pour confirmation. Si discordance A+/B- $\to$ Test C (arbitre).
    \item \textbf{Prise en charge :} Tout test confirmé positif ouvre droit à la mise sous ARV immédiate (stratégie « Test and Treat » / « Dépister et Traiter » recommandée par l'OMS et appliquée au Cameroun depuis 2017), quel que soit le taux de LT4.
\end{itemize}
\end{exemplebox}

\courssub{Conséquences socioculturelles et compétences psychosociales}

L'épidémie de VIH/Sida n'est pas seulement biomédicale ; elle est profondément sociale. Au Cameroun, comme ailleurs, la peur, l'ignorance et les préjugés moraux ont généré une stigmatisation lourde de conséquences.

\begin{savaistu}[Stigmatisation et Droits Humains au Cameroun]
La \textbf{Loi n°2021/013 du 09 juillet 2021} relative à la prévention, la prise en charge et le contrôle du VIH/Sida au Cameroun consacre :
\begin{itemize}
    \item Le droit au consentement libre et éclairé pour le dépistage (interdiction du dépistage obligatoire à l'embauche, au mariage, à l'école).
    \item La confidentialité absolue du statut sérologique (secret médical).
    \item La non-discrimination : accès aux soins, à l'emploi, à l'éducation, au logement sans distinction.
    \item La protection des personnes vulnérables (femmes, enfants, minorités sexuelles, usagers de drogues, détenus).
\end{itemize}
Pourtant, la stigmatisation persiste : rejet familial, exclusion scolaire, licenciement abusif, violences conjugales après annonce de la séropositivité. Cela pousse au silence, empêche le dépistage, brise l'observance thérapeutique et alimente l'épidémie.
\end{savaistu}

\begin{aretenir}[Compétences psychosociales : Vivre avec le VIH / Vivre avec les PVVIH]
\textbf{Pour la personne vivant avec le VIH (PVVIH) :}
\begin{itemize}
    \item \textbf{Acceptation \& Estime de soi :} Accepter son statut sans culpabilité ni honte. Le VIH est une maladie chronique comme le diabète ou l'hypertension.
    \item \textbf{Observance stricte (Adhérence) :} Prendre ses ARV \textbf{tous les jours, à heure fixe, à vie}. C'est la condition \textit{sine qua non} de l'efficacité (charge virale indétectable) et de la prévention de la résistance. Outils : boîtes à pilules, alarmes téléphone, soutien pair-éducateur.
    \item \textbf{Hygiène de vie :} Alimentation équilibrée, sommeil, activité physique, arrêt tabac/alcool/drogues, vaccination à jour (hépatite B, HPV, grippe, pneumocoque), prévention du paludisme (MILDA).
    \item \textbf{Suivi médical régulier :} Consultation trimestrielle, contrôle charge virale (cible $< 50$ copies/mL) et LT4 tous les 6 mois.
    \item \textbf{Vie affective et sexuelle épanouie :} Information du partenaire (obligation légale de protection/prévention), usage du préservatif ou confiance dans l'I=I (si charge virale indétectable confirmée $> 6$ mois), projet de grossesse sécurisé (PTME).
\end{itemize}

\textbf{Pour l'entourage / la communauté (Solidarité) :}
\begin{itemize}
    \item \textbf{Empathie \& Non-jugement :} Écouter sans moraliser. Séparer la maladie de la personne.
    \item \textbf{Confidentialité :} Ne jamais divulguer le statut sérologique d'un tiers sans son accord explicite.
    \item \textbf{Soutien pratique :} Rappel des rendez-vous, aide au transport pour la clinique, partage des repas (aucune transmission), garde d'enfants.
    \item \textbf{Lutte contre la stigmatisation :} Dénoncer les discriminations, corriger les idées reçues (moustiques, repas), promouvoir le langage respectueux (« personne vivant avec le VIH » et non « sidéen », « contaminé », « victime »).
    \item \textbf{Engagement citoyen :} Participer aux campagnes de sensibilisation (Journée Mondiale du Sida 1er décembre), soutenir les associations de PVVIH (ex: Réseau Camerounais des Associations de Personnes Vivant avec le VIH - RECAP+).
\end{itemize}
\end{aretenir}

\begin{methode}[Produire un outil de sensibilisation (Affiche, Sketch, Dépliant, Spot radio)]
\textbf{Compétence évaluée :} « Produire des outils de sensibilisation dans le cadre de la prévention du VIH/Sida ».
\begin{enumerate}
    \item \textbf{Analyser la cible :} Jeunes scolarisés ? Femmes enceintes ? Travailleurs du sexe ? Population générale ? Langue (français, anglais, langues locales : fulfulde, ewondo, bassa, pidgin) ?
    \item \textbf{Définir UN message clé clair et unique :} Ex: « Dépistez-vous ensemble avant d'arrêter le préservatif » / « VIH indétectable = Intransmissible » / « Allaitement maternel exclusif 6 mois + ARV = Bébé sain ».
    \item \textbf{Choisir le support et le ton :} Affiche visuelle (peu de texte, images fortes), Sketch (humour, situations vécues, dialogue), Dépliant (infos pratiques : où, quand, comment), Spot radio (court, percutant, en langue locale).
    \item \textbf{Valider le contenu scientifique :} Vérifier l'exactitude (modes transmission, fenêtre sérologique, efficacité ARV, gratuité). \textbf{Zéro stigmatisation} dans les visuels/mots.
    \item \textbf{Appel à l'action concret :} « Viens te faire tester gratuitement au Centre de Santé X, du lundi au vendredi, 7h-15h » / « Appelle le 117 (Numérique Vert Santé) pour info anonyme ».
    \item \textbf{Tester \& Diffuser :} Montrer à un petit groupe cible, corriger, puis diffuser (lycée, marché, église, mosquée, réseaux sociaux, radio communautaire).
\end{enumerate}
\end{methode}

\begin{exoresolu}[Analyse d'une situation clinique : Suivi thérapeutique]
\textbf{Énoncé :} Marie, 28 ans, diagnostiquée VIH+ en 2020 (LT4 = 350/$\mu$L, CV = 50 000 copies/mL). Elle démarre un traitement ARV (TDF/3TC/DTG) en 2020. En 2024, ses résultats sont : LT4 = 720/$\mu$L, CV = $< 20$ copies/mL (indétectable). Elle souhaite avoir un enfant avec son partenaire séronégatif.
\begin{enumerate}
    \item Interpréter l'évolution biologique de Marie sous traitement.
    \item Quel conseil donner pour son projet de grossesse ? Quel risque pour le partenaire ? Pour l'enfant ?
\end{enumerate}

\tcblower

\textbf{Solution détaillée :}
\begin{enumerate}
    \item \textbf{Interprétation :} Le traitement est \textbf{efficace}. La charge virale est passée de $5 \times 10^4$ à $< 20$ copies/mL (indétectable), prouvant l'inhibition de la réplication virale (blocage transcriptase inverse + intégrase). Le taux de LT4 est remonté de 350 à 720/$\mu$L (reconstitution immunitaire), sortant de la zone de risque d'infections opportunistes ($> 500$). Marie est en \textbf{succès virologique et immunologique}.
    \item \textbf{Conseils :}
    \begin{itemize}
        \item \textbf{Partenaire :} Risque de transmission sexuelle \textbf{nul} (I=I : Indétectable = Intransmissible), condition : charge virale indétectable maintenue $> 6$ mois, observance parfaite, pas d'autres IST. Le préservatif reste recommandé pour prévenir autres IST/grossesse non désirée, mais non indispensable pour le VIH seul dans ce contexte.
        \item \textbf{Grossesse :} Projet \textbf{faisable et sécurisé}. Maintenir le traitement actuel (DTG compatible grossesse, surveillance folates). Suivi mensuel charge virale (cible indétectable à l'accouchement). Accouchement par voies basses si CV indétectable. Prophylaxie NVP/AZT au nouveau-né 4-6 semaines. Allaitement maternel exclusif 6 mois recommandé (sous couvert ARV maternel efficace). Dépistage enfant à la naissance (PCR ARN), 6 semaines, 9 mois, fin allaitement.
    \end{itemize}
\end{enumerate}
\end{exoresolu}

\begin{exercice}[Cas pratique : Fenêtre sérologique et prise de risque]
Paul a eu un rapport sexuel non protégé il y a 10 jours avec une partenaire dont il ignore le statut sérologique. Inquiet, il se rend au CEC du district de santé demain pour faire un test rapide.
\begin{enumerate}
    \item Le test rapide (recherche d'anticorps) peut-il être fiable aujourd'hui ? Justifier.
    \item Quelle conduite tenir immédiatement (aujourd'hui) et dans 3 mois ?
    \item Si le test rapide d'aujourd'hui est négatif, Paul peut-il être rassuré définitivement ?
\end{enumerate}
\end{exercice}

\begin{corrige}[Exercice : Fenêtre sérologique et prise de risque]
\begin{enumerate}
    \item \textbf{Non, le test n'est pas fiable aujourd'hui.} Le délai de 10 jours correspond à la \textbf{fenêtre sérologique}. Les anticorps anti-VIH ne sont pas encore produits en quantité détectable par les tests sérologiques standards (TDR/ELISA). Le test sera vraisemblablement \textbf{faux négatif} si Paul a été contaminé. La charge virale, elle, est déjà très élevée (primo-infection).
    \item \textbf{Conduite immédiate (Aujourd'hui, J+10) :} Paul est dans le délai pour bénéficier d'une \textbf{PEP (Prophylaxie Post-Exposition)} (délai max 48h, idéalement $< 4$h, mais on peut tenter jusqu'à 72h dans certains protocoles). Il doit se rendre \textbf{immédiatement} aux urgences ou au CEC pour évaluation et prescription d'ARV d'urgence (28 jours). \textbf{Conduite à 3 mois (J+90) :} Refaire un test de dépistage (anticorps) pour lever le doute définitif (fin de fenêtre sérologique).
    \item \textbf{Non.} Un test négatif à J+10 ne signifie pas « non infecté », il signifie « anticorps non détectés à ce jour ». Paul doit \textbf{impérativement} refaire le test à 3 mois (et idéalement à 6 semaines aussi) pour confirmer sa séro-négativité. En attendant, il doit se protéger (préservatif) pour ne pas transmettre le virus s'il est en primo-infection (très contagieux).
\end{enumerate}
\end{corrige}

\begin{savaistu}[I = I : Une révolution scientifique et sociale]
Depuis les études \textbf{HPTN 052 (2011)}, \textbf{PARTNER 1 \& 2 (2016, 2018)} et \textbf{Opposites Attract (2018)}, le consensus scientifique mondial (OMS, ONUSIDA, CDC, Société Française de Lutte contre le Sida, Programme National de Lutte contre le Sida du Cameroun) est formel : \textbf{Une personne vivant avec le VIH, sous traitement antirétroviral efficace depuis au moins 6 mois, avec une charge virale durablement indétectable ($< 50$ ou $< 200$ copies/mL selon les seuils), NE TRANSMET PAS le VIH par voie sexuelle.} C'est le principe \textbf{I = I (Indétectable = Intransmissible / U = U Undetectable = Untransmittable)}. Cette donnée change radicalement la vie des couples sérodifférents, la prévention (TasP : Treatment as Prevention) et la lutte contre la stigmatisation : le traitement \textbf{est} prévention.
\end{savaistu}

\newpage

\coursec{Activité d'intégration}

\coursec{Les mécanismes de l'immunité et ses dysfonctionnements}

\courssub{Objectifs de l'activité}
\begin{itemize}
    \item Mobiliser l'ensemble des savoirs et savoir-faire de la leçon pour analyser des documents expérimentaux et cliniques variés.
    \item Interpréter des résultats de greffe pour illustrer la reconnaissance du non-soi par le \cle{CMH}.
    \item Exploiter une courbe de cinétique anticorps pour caractériser la réponse immune primaire, secondaire et la \cle{mémoire immunitaire}.
    \item Analyser un cas clinique de diabète juvénile pour identifier un mécanisme \cle{auto-immun} (perte de tolérance au soi).
    \item Exploiter l'évolution des taux de LT4\texorpdfstring{+}{+} et de charge virale pour expliquer la physiopathologie du \cle{VIH/Sida}.
    \item Concevoir collectivement un outil de sensibilisation (affiche) intégrant prévention, modes de transmission et solidarité envers les PVVIH.
\end{itemize}

\courssub{Situation}
\textbf{Contexte :} Dans le cadre de la semaine « Santé et Citoyenneté » au Lycée Bilingue de Nkolbisson (Yaoundé), la classe de Terminale C est invitée par le Centre de Santé Intégré (CSI) du quartier à participer à une « Enquête Immunologique ». Le Dr \textsc{Mvondo}, médecin-chef, remet à chaque groupe de 4 élèves un dossier patient anonymisé composé de quatre documents. L'objectif est de rendre compte, lors d'une restitution orale, des mécanismes immunologiques en jeu et de proposer une affiche de prévention pour le hall du lycée.

\vspace{0.5em}
\noindent\textbf{Document 1 : Expérience de greffe de peau chez la souris (Modèle de Medawar, 1944, adapté).}
On prélève un fragment de peau (greffon) sur une souris donneuse de lignée pure \textbf{A} (homozygote pour son CMH) et on le greffe sur une souris receveuse de lignée pure \textbf{B}. On observe le sort du greffon au cours du temps. L'expérience est répétée dans quatre conditions :
\begin{center}
\begin{tabularx}{\linewidth}{|l|X|X|X|}\hline
\textbf{Condition} & \textbf{Donneur (Greffon)} & \textbf{Receveur (Hôte)} & \textbf{Résultat observé (Jour du rejet ou survie)} \\ \hline
1 & Souris A & Souris A (syngénique) & Greffon viable indéfiniment ($> \SI{100}{jours}$) \\ \hline
2 & Souris A & Souris B (allogénique) & Rejet complet vers le \textbf{12\textsuperscript{e} jour} \\ \hline
3 & Souris A & Souris B ayant déjà rejeté un greffon A (seconde greffe A) & Rejet accéléré vers le \textbf{6\textsuperscript{e} jour} \\ \hline
4 & Souris A & Souris B irradiée (dose létale : \SI{9}{Gy}) + injection de cellules de rate de souris B non irradiée (reconstitution hématopoïétique) & Greffon A rejeté vers le \textbf{12\textsuperscript{e} jour} \\ \hline
\end{tabularx}
\end{center}

\vspace{0.5em}
\noindent\textbf{Document 2 : Cinétique de la réponse humorale après vaccination (Vaccin tétanos, adjuvant aluminium).}
Un volontaire adulte, séronégatif pour la toxine tétanique, reçoit deux injections sous-cutanées du vaccin (J0 et J28). Le taux d'anticorps anti-toxine tétanique (IgG, en UI/mL) est dosé par ELISA à intervalles réguliers.
\begin{popfigure}
\begin{tikzpicture}
\begin{axis}[
    width=\linewidth, height=7cm,
    xlabel={Temps (jours)}, ylabel={Taux d'IgG anti-tétanos (UI/mL)},
    xmin=0, xmax=70, ymin=0, ymax=5.5,
    xtick={0,7,14,21,28,35,42,49,56,63,70},
    ytick={0,1,2,3,4,5},
    grid=major, grid style={popline},
    axis lines=middle,
    enlargelimits=false,
    clip=false,
    legend style={at={(0.98,0.98)}, anchor=north east, fill=white, draw=popdark},
]
% Primary response
\addplot[popblue, thick, smooth, domain=0:28, samples=100] 
    ({x}, {0.05*x*(x<7) + (0.35 + 1.5*exp(-(x-14)^2/30))*(x>=7)});
% Secondary response
\addplot[poporange, thick, smooth, domain=28:70, samples=100] 
    ({x}, {0.1 + 4.5*exp(-(x-38)^2/40)});
% Annotations
\draw[popblue, dashed] (14,0) -- (14,2.5) node[above, pos=0.5, right] {Pic 1\textsuperscript{er} pic};
\draw[poporange, dashed] (38,0) -- (38,4.8) node[above, pos=0.5, right] {Pic 2\textsuperscript{e} pic};
\draw[popmuted, dashed] (28,0) -- (28,1.2) node[above, pos=0.5, right] {Rappel (J28)};
\addlegendentry{Évolution du taux d'IgG}
\end{axis}
\end{tikzpicture}
\legende{Cinétique de la réponse humorale (IgG) après primo-vaccination (J0) et rappel (J28) contre le tétanos. Le pic primaire survient vers J14 (1,8 UI/mL), le pic secondaire vers J38 (4,6 UI/mL), soit 10 jours post-rappel.}
\end{popfigure}
\textit{Données chiffrées clés :} Pic primaire $\approx \SI{1,8}{UI/mL}$ au J14 ; Taux résiduel J28 $\approx \SI{0,4}{UI/mL}$ ; Pic secondaire $\approx \SI{4,6}{UI/mL}$ au J38 (soit 10 jours post-rappel).

\vspace{0.5em}
\noindent\textbf{Document 3 : Cas clinique — Diabète juvénile (DT1) à l'Hôpital Gynéco-Obstétrique et Pédiatrique de Yaoundé.}
\textbf{Patient :} K., 14 ans, scolarisé en 3\textsuperscript{e} au Lycée de Mfou.
\textbf{Antécédents :} Aucun. Vaccinations à jour (PEV). Pas de diabète familial au 1\textsuperscript{er} degré.
\textbf{Circonstances de découverte :} Polyuro-polydipsie (PUPD) majeure depuis 3 semaines, amaigrissement de \SI{5}{kg} en 1 mois, fatigue. Admis en urgence pour acidocétose diabétique (pH artériel \num{7,22} ; bicarbonates \SI{12}{mmol/L} ; glycémie \SI{28}{mmol/L} [\SI{504}{mg/dL}] ; cétonémie $+++$).
\textbf{Bilan immunologique (prélevé à J0, avant insulinothérapie) :}
\begin{itemize}
    \item Auto-anticorps anti-GAD65 (glutamate décarboxylase 65 kDa) : \textbf{\SI{125}{UI/mL}} (N $< \SI{5}{UI/mL}$).
    \item Auto-anticorps anti-IA2 (tyrosine phosphatase) : \textbf{\SI{87}{UI/mL}} (N $< \SI{7,5}{UI/mL}$).
    \item Auto-anticorps anti-insuline (IAA) : \textbf{\SI{42}{nmol/L}} (N $< \SI{0,4}{nmol/L}$).
    \item HLA-DR3/DR4 hétérozygote (phénotype à risque).
\end{itemize}
\textbf{Évolution :} Rééquilibrage métabolique sous insuline rapide (pompe) + insuline basale. Éducation thérapeutique initiée. Pronostic : insulinodépendance à vie.

\vspace{0.5em}
\noindent\textbf{Document 4 : Histoire naturelle de l'infection VIH-1 non traitée (Données cohortes ANRS / Cameroun).}
Un patient masculin, 32 ans, séroconverti estimé à $t=0$ (infection primaire symptomatique type mononucléose). Suivi sans ARV (refus initial, puis perte de vue). Mesures trimestrielles :
\begin{popfigure}
\begin{tikzpicture}
\begin{axis}[
    width=\linewidth, height=7cm,
    xlabel={Temps après infection (années)}, ylabel={Taux de LT4$^+$ (cellules/$\mu$L)},
    xmin=0, xmax=10, ymin=0, ymax=1300,
    xtick={0,1,2,3,4,5,6,7,8,9,10},
    ytick={0,200,400,600,800,1000,1200},
    grid=major, grid style={popline},
    axis y line=left, axis x line=bottom,
    legend style={at={(0.5,-0.15)}, anchor=north, fill=white, draw=popdark},
]
% CD4 curve
\addplot[popblue, very thick, smooth, mark=*, mark size=2] coordinates {
    (0, 850) (0.5, 550) (1, 600) (2, 550) (3, 500) (4, 450) (5, 380) (6, 300) (7, 220) (8, 150) (9, 80) (10, 40)
};
\addlegendentry{LT4$^+$ (échelle gauche)}
% Viral load - second axis
\end{axis}
\begin{axis}[
    width=\linewidth, height=7cm,
    xlabel={Temps après infection (années)}, ylabel={Charge virale (copies ARN/mL)},
    xmin=0, xmax=10, ymin=1, ymax=10000000,
    xtick={0,1,2,3,4,5,6,7,8,9,10},
    ytick={1,100,10000,1000000},
    ymode=log,
    axis y line=right, axis x line=none,
    legend style={at={(0.5,-0.15)}, anchor=north, fill=white, draw=popdark},
]
\addplot[poporange, very thick, smooth, mark=square*, mark size=2] coordinates {
    (0, 5000000) (0.5, 50000) (1, 30000) (2, 50000) (3, 80000) (4, 120000) (5, 200000) (6, 350000) (7, 600000) (8, 800000) (9, 1500000) (10, 3000000)
};
\addlegendentry{Charge virale (échelle droite, log)}
\end{axis}
\end{tikzpicture}
\legende{Évolution parallèle du taux de lymphocytes T CD4$^+$ (échelle linéaire, bleu) et de la charge virale VIH-1 (échelle logarithmique, orange) au cours de l'infection non traitée. Seuil Sida (LT4$^+$ < 200/µL) franchi vers l'année 7,5-8.}
\end{popfigure}
\textit{Repères cliniques :} Seuil Sida (stade C CDC) : LT4$^+$ $< \SI{200}{\per\micro\liter}$ ou infection opportuniste. Charge virale « set point » vers 1-2 ans $\approx 5 \times 10^4$ copies/mL.

\courssub{Consignes}
\textbf{Partie A : Analyse des documents (Travail de groupe, 45 min)}
\begin{enumerate}
    \item \textbf{Doc 1 :} Comparer les conditions 1 et 2. Quel rôle joue le CMH dans l'acceptation ou le rejet du greffon ? Expliquer le mécanisme cellulaire (lymphocytes T CD8$^+$ cytotoxiques).
    \item \textbf{Doc 1 :} Comparer les conditions 2 et 3. Quelle propriété fondamentale de l'immunité spécifique est mise en évidence ? Préciser le type de cellules responsables de cette accélération.
    \item \textbf{Doc 1 :} Interpréter la condition 4. Pourquoi la souris B irradiée reconstituée avec sa propre moelle rejette-t-elle le greffon A ? Que prouve cela sur l'origine des cellules effectrices ?
    \item \textbf{Doc 2 :} Décrire la courbe : phases de latence, croissance, pic, décroissance pour la primo-injection. Faire de même pour le rappel. Comparer les paramètres cinétiques (délai, amplitude, affinité supposée) et isotypes probables.
    \item \textbf{Doc 2 :} Relier la forme de la courbe secondaire à la mémoire immunitaire. Citer les deux populations cellulaires mémoires impliquées (LBm, LT4m) et leur rôle respectif.
    \item \textbf{Doc 3 :} À partir des auto-anticorps détectés (anti-GAD65, IA2, insuline), identifier la cible de l'attaque auto-immune. Quel type de diabète s'agit-il ? Expliquer le mécanisme : rupture de la tolérance centrale/périphérique $\rightarrow$ activation de LT4$^+$ autoréactifs $\rightarrow$ activation LT8$^+$ cytotoxiques et LB $\rightarrow$ destruction des cellules $\beta$ pancréatiques.
    \item \textbf{Doc 3 :} Quel est le rôle du génotype HLA-DR3/DR4 ? (Prédisposition génétique : présentation efficace de peptides auto aux LT4$^+$).
    \item \textbf{Doc 4 :} Décrire l'évolution parallèle des deux courbes (LT4$^+$ et charge virale) en 3 phases : infection primaire, phase de latence clinique (asymptomatique), phase Sida.
    \item \textbf{Doc 4 :} Expliquer la chute brutale des LT4$^+$ lors de l'infection primaire (lyse virale directe, apoptose bystander) puis la lente érosion pendant la latence (activation chronique, épuisement, destruction des niches lymphoïdes). Pourquoi la charge virale remonte-t-elle en phase terminale ?
    \item \textbf{Doc 4 :} À quel moment (année approximative) le patient bascule-t-il au stade Sida (critère immunologique LT4$^+$ $< \SI{200}{\per\micro\liter}$) ? Quelles infections opportunistes sont à craindre au Cameroun à ce stade (ex. : tuberculose, toxoplasmose cérébrale, cryptococcose, diarrhées à \emph{Cryptosporidium}) ?
\end{enumerate}

\textbf{Partie B : Production collective — Affiche de sensibilisation (30 min + restitution)}
\begin{enumerate}
    \setcounter{enumi}{10}
    \item En classe entière, concevoir une affiche A3 destinée au tableau d'affichage du lycée. Elle doit comporter :
    \begin{itemize}
        \item Un titre accrocheur.
        \item Les 3 modes de transmission (sang, sexuel, mère-enfant) avec pictogrammes.
        \item Les gestes \textbf{SANS RISQUE} (baiser, main, repas, moustique, piscine, toilettes) pour lutter contre la stigmatisation.
        \item Le message : « Vivre avec le VIH, c'est d'abord vivre. Soutenons les PVVIH. »
        \item Les numéros utiles au Cameroun : 117 (SOS Sida), 1510 (Allô Santé), Centre de Prise en Charge (PEC) le plus proche.
    \end{itemize}
\end{enumerate}

\courssub{Ressources à mobiliser}
\begin{itemize}
    \item \textbf{Reconnaissance soi/non-soi :} CMH de classe I (sur toutes cellules nucléées) et classe II (sur CPA : macrophages, cellules dendritiques, LB). Récepteur TCR (LT) / BCR (LB). Restriction par le CMH.
    \item \textbf{Réponse spécifique :} Médiation cellulaire (LT8$^+$ cytotoxiques $\rightarrow$ perforine/granzyme/FasL ; LT4$^+$ Th1 $\rightarrow$ IFN-$\gamma$ activation macrophages). Médiation humorale (LB $\rightarrow$ Plasmocytes $\rightarrow$ Ac ; LBm).
    \item \textbf{Mémoire :} LBm (réponse Ac rapide, haute affinité, IgG) + LT4m (aide rapide). Justifie la vaccination (rappel).
    \item \textbf{Auto-immunité :} Perte de tolérance (centrale thymique : sélection négative défaillante ; périphérique : Tregs défaillants, mimétisme moléculaire, libération d'antigènes cachés). Atteinte d'organe cible (cellules $\beta$ pancréatiques $\rightarrow$ DT1).
    \item \textbf{VIH/Sida :} VIH = rétrovirus à tropisme CD4 (LT4$^+$, macrophages, cellules dendritiques). Intégration provirus. Destruction LT4$^+$ $\rightarrow$ immunodépression. Traitement = ARV (trithérapie) $\rightarrow$ charge virale indétectable = intransmissible (I=I). Prévention : PrEP, PEP, préservatifs, dépistage, TASP.
\end{itemize}

\begin{popfigure}
\begin{tikzpicture}[node distance=1.2cm, >=Stealth, font=\small, every node/.style={align=center}]
\node (hsc) [draw, rounded rectangle, fill=popblueL, minimum width=3cm, minimum height=1cm, align=center]{Cellule souche hématopoïétique (CSH)\\(Moelle osseuse rouge)};
\node (myeloid) [draw, rectangle, fill=poporangeL, below left=of hsc, xshift=-1.5cm, minimum width=2.8cm, minimum height=1cm, align=center]{Lignée myéloïde\\(Globules rouges, Granulocytes,\\Monocytes $\rightarrow$ Macrophages/\textsc{cpa})};
\node (lymphoid) [draw, rectangle, fill=popgreenL, below right=of hsc, xshift=1.5cm, minimum width=2.8cm, minimum height=1cm, align=center]{Lignée lymphoïde\\(Pro-lymphocytes)};
\node (tcell) [draw, ellipse, fill=poppurpleL, below=of lymphoid, yshift=-0.8cm, minimum width=3cm, minimum height=1.2cm, align=center]{Précurseurs T $\rightarrow$ \textbf{Thymus}\\(Maturation : sélection +/−)\\\textbf{LT4$^+$ (CD4), LT8$^+$ (CD8)}};
\node (bcell) [draw, ellipse, fill=poppinkL, below=of tcell, yshift=-0.8cm, minimum width=3cm, minimum height=1cm, align=center]{Maturation B (Moelle osseuse)\\\textbf{LB naïfs (IgM/IgD de surface)}};
\draw[->, thick, popblue] (hsc) -- (myeloid);
\draw[->, thick, popblue] (hsc) -- (lymphoid);
\draw[->, thick, popblue] (lymphoid) -- (tcell);
\draw[->, thick, popblue] (lymphoid) -- (bcell);
\draw[->, thick, poporange, dashed] (tcell) -- ++(-2.5,0) node[left, text width=2.5cm] {Migration vers organes lymphoïdes secondaires (Ganglions, Rate, MALT)};
\draw[->, thick, poporange, dashed] (bcell) -- ++(2.5,0) node[right, text width=2.5cm] {Migration vers organes lymphoïdes secondaires (Ganglions, Rate, MALT)};
\end{tikzpicture}
\legende{Schéma récapitulatif : origine et maturation des cellules immunitaires. Toutes les cellules sanguines dérivent de la CSH (moelle osseuse). Les lymphocytes T mûrissent dans le thymus (sélection positive/négative sur CMH), les lymphocytes B dans la moelle osseuse. Les lymphocytes naïfs colonisent les organes lymphoïdes secondaires où ils rencontrent l'antigène.}
\end{popfigure}

\courssub{Démarche et corrigé commenté}
\begin{exoresolu}[Enquête immunologique : Corrigé guidé]
\textbf{Document 1 — Reconnaissance du non-soi et mémoire de greffe}
\begin{enumerate}
    \item \textbf{Cond 1 vs 2 :} Syngénique (CMH identique) $\rightarrow$ tolérance (pas de reconnaissance « non-soi »). Allogénique (CMH différent) $\rightarrow$ rejet aigu (J12). Les LT8$^+$ cytotoxiques du receveur reconnaissent les peptides allogéniques présentés par le CMH I du greffon (voie directe : LT8$^+$ hôte voient CMH I donneur + peptide ; voie indirecte : CPA hôte présentent peptides donneur sur CMH II hôte aux LT4$^+$). Destruction des cellules du greffon par perforine/granzyme B et voie Fas/FasL.
    \item \textbf{Cond 2 vs 3 :} Rejet en J12 (primo-réponse) vs J6 (secondaire). \textbf{Mémoire immunitaire}. Présence de LTm (CD8$^+$m et CD4$^+$m) spécifiques du CMH A, activés plus vite, sans phase de latence longue, proliférant massivement.
    \item \textbf{Cond 4 :} Irradiation \SI{9}{Gy} détruit la moelle osseuse (hématopoïèse) $\rightarrow$ pas de lymphocytes matures. Reconstitution par rate B (cellules souches hématopoïétiques + lymphocytes matures B). Le système immunitaire reconstruit est de phénotype B (CMH B). Il reconnaît le greffon A comme non-soi $\rightarrow$ rejet en J12 (kinétique primo-réponse). \textbf{Preuve :} Les cellules effectrices du rejet (LT) dérivent de la moelle osseuse (hématopoïèse), pas des tissus résidents radio-résistants.
\end{enumerate}

\textbf{Document 2 — Mémoire immunitaire et vaccination}
\begin{enumerate}
    \setcounter{enumi}{3}
    \item \textbf{Primo-injection (J0-J28) :} Phase de latence $\approx \SI{7}{jours}$ (activation LB naïfs, présentation par CPA, aide LT4$^+$, différenciation en plasmocytes). Pic IgM (non dosé ici) puis IgG vers J14 (\SI{1,8}{UI/mL}). Décroissance (demi-vie IgG $\approx \SI{21}{jours}$, mort des plasmocytes effecteurs). Taux résiduel J28 (\SI{0,4}{UI/mL}) = Ac résiduels + plasmocytes vieillissants.
    \item \textbf{Rappel (J28) :} Réactivation immédiate des LBm et LT4m. Pas de phase de latence (pic J38, soit 10 j post-rappel). Amplitude $\times 2,5$ (\SI{4,6}{UI/mL} vs \SI{1,8}{UI/mL}). Affinité supérieure (mutation somatique, sélection). Isotype : quasi exclusivement IgG (commutation de classe déjà faite).
    \item \textbf{Mécanisme :} LBm $\xrightarrow{\text{Ag + aide LT4m}}$ Plasmocytes (forte sécrétion Ac haute affinité) + nouveaux LBm. LT4m $\xrightarrow{\text{Ag sur CMH II}}$ sécrétion rapide IL-2, IL-4, IL-21 $\rightarrow$ prolifération/différenciation LBm. C'est la base de la vaccination : mimiquer l'infection sans pathologie pour générer cette mémoire.
\end{enumerate}

\textbf{Document 3 — Diabète de type 1 : Maladie auto-immune à médiation cellulaire et humorale}
\begin{enumerate}
    \setcounter{enumi}{6}
    \item \textbf{Auto-anticorps :} Anti-GAD65, IA2, Insuline = marqueurs spécifiques de la destruction auto-immune des \textbf{cellules $\beta$ des îlots de Langerhans}. Diagnostic : \textbf{Diabète de type 1 (DT1) / Diabète juvénile}.
    \item \textbf{Mécanisme :} Prédisposition génétique (HLA-DR3/DR4 : groove de liaison peptidique présentant efficacement peptides GAD/IA2/Insuline aux LT4$^+$). Déclencheur environnemental (virus ? Coxsackie B ?). Rupture tolérance périphérique (Tregs CD4$^+$CD25$^+$FoxP3$^+$ défaillants). Activation LT4$^+$ Th1 autoréactifs $\rightarrow$ IFN-$\gamma$, TNF-$\alpha$ $\rightarrow$ activation macrophages, expression CMH I/II sur cellules $\beta$ (habituellement CMH I faible, pas de CMH II). Activation LT8$^+$ cytotoxiques autoréactifs (reconnaissance peptides auto sur CMH I) $\rightarrow$ lyse cellules $\beta$. LB activés $\rightarrow$ auto-Ac (marqueurs diagnostiques, pas effecteurs principaux de la destruction). $\rightarrow$ Déficit absolu en insuline.
    \item \textbf{HLA-DR3/DR4 :} Allèles de classe II du CMH conférant un risque relatif $\times 15$ à $\times 20$. Hétérozygote DR3/DR4 = risque maximal (hétérozygotie avantageuse pour présentation de peptides variés, mais ici désavantageuse pour tolérance).
\end{enumerate}

\textbf{Document 4 — Dynamique VIH / LT4$^+$ : Vers le Sida}
\begin{enumerate}
    \setcounter{enumi}{9}
    \item \textbf{Phase 1 : Infection primaire (0-6 mois).} Pic charge virale ($> 10^6$ copies/mL) $\rightarrow$ dissémination. Chute brutale LT4$^+$ (\SI{850}{\per\micro\liter} $\rightarrow$ \SI{550}{\per\micro\liter}) : infection directe + lyse + apoptose bystander (GP120 soluble, activation immunitaire). Séroconversion.
    \item \textbf{Phase 2 : Latence clinique (6 mois - $\approx \SI{7-8}{ans}$).} Équilibre précaire : réplication virale continue (ganglions, rate, GALT) vs production LT4$^+$ (thymus, prolifération périphérique). « Set point » charge virale $\approx 5 \times 10^4$. Érosion lente LT4$^+$ ($\approx \SI{50-80}{\per\micro\liter/an}$) : activation chronique $\rightarrow$ épuisement (PD-1$^+$), fibrose du réseau réticulaire ganglionnaire (destruction niches stromales), mort par apoptose (Fas/FasL, TNF).
    \item \textbf{Phase 3 : Phase Sida (après $\approx \SI{8}{ans}$ ici).} LT4$^+$ franchissent seuil \SI{200}{\per\micro\liter} vers \textbf{l'année 7,5-8}. Charge virale repart à la hausse (perte contrôle immune, émergence variants X4 tropisme CXCR4 plus cytopathiques). Infections opportunistes (IO) : \emph{Mycobacterium tuberculosis} (TB pulmonaire/extrapulmonaire, 1\textsuperscript{re} cause décès VIH au Cameroun), \emph{Toxoplasma gondii} (encéphalite), \emph{Cryptococcus neoformans} (méningite), \emph{Pneumocystis jirovecii} (pneumonie), cytomégalovirus, diarrhées chroniques (\emph{Cryptosporidium}, \emph{Isospora}).
    \item \textbf{Basculement Sida :} Vers l'année \textbf{7,5 - 8} (LT4$^+$ $\approx \SI{200}{\per\micro\liter}$). Sans ARV, survie médiane $\approx \SI{1-2}{ans}$ après 1\textsuperscript{re} IO.
\end{enumerate}

\textbf{Partie B — Affiche de sensibilisation (Exemple de structure)}
\begin{center}
\begin{tabularx}{\linewidth}{|X|}\hline
\rowcolor{popblueL} \textbf{\Large \textsc{VIH/SIDA : CONNAÎTRE, PRÉVENIR, SOUTENIR}} \\ \hline
\begin{minipage}{\linewidth}
\textbf{COMMENT S'ATTRAPE-T-IL ? (3 voies)}
\begin{itemize}
    \item \textcolor{popred}{\textbf{Sang :}} Seringues partagées (drogue, tatouage, scarifications), accidents d'exposition (soignants), transfusion non testée (rare aujourd'hui).
    \item \textcolor{popred}{\textbf{Sexuel :}} Rapports non protégés (vaginal, anal, oral) avec partenaire séropositif non traité ou statut inconnu.
    \item \textcolor{popred}{\textbf{Mère-Enfant :}} Grossesse, accouchement, allaitement (évitable par ARV mère + enfant).
\end{itemize}
\textbf{GESTES SANS AUCUN RISQUE (Zéro transmission) :}
\begin{itemize}
    \item S'embrasser, se serrer la main, se prendre dans les bras.
    \item Manger ensemble, partager verre/couverts, toilettes, piscine, moustiques.
    \item Travailler, étudier, jouer ensemble.
\end{itemize}
\textbf{MES MOYENS DE PROTECTION :}
\begin{itemize}
    \item \textbf{Préservatif} (masculin/féminin) à chaque rapport.
    \item \textbf{Dépistage} régulier (gratuit, confidentiel au CSI/CEPEC).
    \item \textbf{PrEP} (Prophylaxie Pré-Exposition) si risque élevé.
    \item \textbf{PEP} (Post-Exposition) $< \SI{48}{h}$ après accident.
    \item \textbf{TASP} : Personne sous ARV, charge virale indétectable $=$ \textbf{NON TRANSMISSIBLE} (I = I).
\end{itemize}
\textbf{\textcolor{popgreen}{VIVRE AVEC LE VIH, C'EST D'ABORD VIVRE.}}
\textbf{Stop à la stigmatisation ! Une personne séropositive traitée est une personne en bonne santé.}
\vspace{0.5em}
\textbf{NUMÉROS UTILES CAMEROUN :}
\begin{itemize}
    \item \textbf{117} (SOS Sida / Écoute/Info)
    \item \textbf{1510} (Allô Santé / Ministère Santé Publique)
    \item \textbf{CEPEC le plus proche :} Hôpital Central / Hôpital Général / Hôpital de District de votre arrondissement.
\end{itemize}
\end{minipage} \\ \hline
\end{tabularx}
\end{center}
\end{exoresolu}

\begin{attention}[Pièges fréquents et points de vigilance pour le Bac]
\begin{itemize}
    \item \textbf{Greffe :} Ne pas confondre \emph{voie directe} (LT hôte reconnaissent CMH donneur intact sur cellules donneuses) et \emph{voie indirecte} (CPA hôte présentent peptides donneurs sur CMH hôte). Les deux coexistent. Le rejet aigu est médié principalement par les LT8$^+$ (voie directe) et LT4$^+$ Th1 (voie indirecte).
    \item \textbf{Mémoire :} La réponse secondaire est plus \emph{rapide} (pas de latence), plus \emph{intense} (amplitude), de \emph{meilleure affinité} (mutation somatique) et de \emph{classe IgG} (commutation déjà réalisée). Les LBm ne sécrètent pas d'Ac au repos ; ils se différencient en plasmocytes au contact de l'Ag.
    \item \textbf{Auto-immunité (DT1) :} Les auto-anticorps (anti-GAD, IA2, IAA) sont des \textbf{marqueurs diagnostiques} (dépistage précoce, prédiction), mais les \textbf{effecteurs de la destruction} sont les LT8$^+$ cytotoxiques (et macrophages activés par Th1). Ne pas écrire « les anticorps détruisent les cellules $\beta$ ».
    \item \textbf{VIH/Sida :} La charge virale « set point » détermine la vitesse d'évolution. L'épuisement des LT4$^+$ est multifactoriel (infection directe, apoptose bystander, activation chronique, fibrose lymphoïde, insuffisance thymique). \textbf{I = I} (Indétectable = Intransmissible) est un message majeur de santé publique : une PVVIH sous ARV efficace (charge virale $< \SI{50}{copies/mL}$) ne transmet pas le virus sexuellement.
    \item \textbf{Unités :} Toujours écrire les unités (\SI{}{\per\micro\liter} pour LT4, copies/mL pour charge virale, UI/mL pour Ac, Gy pour irradiation). Respecter les puissances de 10 (notation scientifique).
\end{itemize}
\end{attention}

\courssub{Bilan de l'activité}
Cette activité d'intégration a permis de :
\begin{itemize}
    \item \textbf{Consolider la notion de « Soi » et « Non-soi »} via le modèle expérimental de la greffe (Doc 1), fondement de la tolérance immunologique et de la surveillance anticancéreuse.
    \item \textbf{Quantifier la mémoire immunitaire} (Doc 2) : cinétique, amplitude, affinité, isotype — base rationnelle du calendrier vaccinal (rappels) et de la sérologie post-vaccinale.
    \item \textbf{Comprendre la rupture de tolérance} (Doc 3) : le DT1 comme modèle d'auto-immunité à médiation cellulaire (LT8$^+$ cytotoxiques) avec marqueurs humoraux (auto-Ac) accessibles au diagnostic précoce. Lien génétique (HLA) + environnement.
    \item \textbf{Modéliser la dynamique VIH/LT4$^+$} (Doc 4) : la course poursuite virus/hôte, l'épuisement du réservoir thymique et ganglionnaire, le seuil critique des \SI{200}{LT4^+/\micro\liter} définissant le Sida, et l'urgence du dépistage précoce et de l'ARV à vie.
    \item \textbf{Développer des compétences psychosociales} (Partie B) : transformer le savoir scientifique en message citoyen, lutter contre la sérophobie (gestes sans risque), promouvoir le dépistage et l'observance (I=I), connaître les ressources locales (numéros verts, CEPEC).
\end{itemize}
L'élève termine cette leçon capable d'argumenter scientifiquement sur une pathologie immune, d'interpréter un bilan sérologique ou une courbe de suivi VIH, et d'agir en acteur de santé publique dans son établissement.

\newpage

\coursec{Méthodes et exercices résolus}

\coursec{Les mécanismes de l'immunité et ses dysfonctionnements}

\courssub{Origine et maturation des cellules immunitaires}

\begin{definition}[Hématopoïèse et lignées immunitaires]
Toutes les cellules du système immunitaire dérivent de \cle{cellules souches hématopoïétiques} (CSH) situées dans la \cle{moelle osseuse rouge} (os plats : sternum, côtes, iliaques ; épiphyses des os longs). Deux grandes lignées se séparent précocement :
\begin{itemize}
\item la \textbf{lignée myéloïde} : donne les polynucléaires neutrophiles, éosinophiles, basophiles, les monocytes (devenant \textbf{macrophages} dans les tissus) et les cellules dendritiques ;
\item la \textbf{lignée lymphoïde} : donne les lymphocytes B (LB), les lymphocytes T (LT) et les cellules NK (Natural Killers).
\end{itemize}
\end{definition}

\begin{popfigure}
\begin{center}
\begin{tikzpicture}[
  node distance=1.2cm and 1.5cm,
  every node/.style={align=center, font=\small},
  box/.style={rectangle, draw=popblue, thick, fill=popblueL, rounded corners, minimum width=2.2cm, minimum height=0.8cm},
  arrow/.style={-Stealth, thick, popdark}
]
\node[box, align=center] (csh){Cellule souche\\ hématopoïétique\\(Moelle osseuse rouge)};
\node[box, below left=of csh] (myelo) {Lignée myéloïde};
\node[box, below right=of csh] (lymph) {Lignée lymphoïde};

\node[box, below left=of myelo, align=center] (mono){Monocyte $\rightarrow$\\ Macrophage (tissus)};
\node[box, below=of myelo, align=center] (poly){Polynucléaires\\(Neutro, Éosino, Baso)};
\node[box, below right=of myelo, align=center] (dc){Cellule\\ dendritique};

\node[box, below left=of lymph, align=center] (lb){Lymphocyte B\\ \textbf{Maturation : Moelle osseuse}};
\node[box, below=of lymph] (nk) {Cellule NK};
\node[box, below right=of lymph, align=center] (lt){Lymphocyte T\\ \textbf{Maturation : Thymus}};

\draw[arrow] (csh) -- (myelo);
\draw[arrow] (csh) -- (lymph);
\draw[arrow] (myelo) -- (mono);
\draw[arrow] (myelo) -- (poly);
\draw[arrow] (myelo) -- (dc);
\draw[arrow] (lymph) -- (lb);
\draw[arrow] (lymph) -- (nk);
\draw[arrow] (lymph) -- (lt);

\node[draw=poporange, thick, fill=poporangeL, rounded corners, fit=(lb)(lt), inner sep=4pt, label={[poporange,yshift=-2pt]above:\textbf{Organes lymphoïdes primaires}}] {};
\node[draw=popgreen, thick, fill=popgreenL, rounded corners, fit=(mono)(poly)(dc)(nk), inner sep=4pt, label={[popgreen,yshift=-2pt]above:\textbf{Effecteurs non spécifiques}}] {};
\end{tikzpicture}
\end{center}
\legende{Schématisation de l'hématopoïèse : distinction des organes de formation (moelle osseuse pour tous) et de maturation (moelle pour LB, thymus pour LT).}
\end{popfigure}

\begin{propriete}[Lieux de formation et de maturation]
\begin{itemize}
\item \textbf{Formation (production)} : moelle osseuse rouge pour \emph{toutes} les cellules immunitaires.
\item \textbf{Maturation (acquisition des récepteurs spécifiques et tolérance du soi)} :
      \begin{itemize}
      \item Lymphocytes \textbf{B} : moelle osseuse (mnémotechnique : \emph{B} = \emph{Bone marrow}).
      \item Lymphocytes \textbf{T} : \textbf{thymus} (mnémotechnique : \emph{T} = \emph{Thymus}). Sélection positive (réognition CMH propre) et négative (élimination autoréactifs).
      \item Polynucléaires, monocytes, NK : moelle osseuse.
      \end{itemize}
\item \textbf{Migration} : cellules matures (naïves) $\rightarrow$ circulation sanguine $\rightarrow$ \textbf{organes lymphoïdes secondaires} (ganglions, rate, amygdales, plaques de Peyer, MALT) où elles rencontrent les antigènes.
\end{itemize}
\end{propriete}

\courssub{Méthode 1 : Déterminer les lieux de formation et de maturation des cellules immunitaires}

\begin{methode}[Retrouver l'origine et le devenir d'une cellule immunitaire]
\begin{enumerate}
\item \textbf{Identifier la cellule étudiée} dans l'énoncé (lymphocyte B, lymphocyte T, macrophage, polynucléaire...).
\item \textbf{Se rappeler le point de départ commun} : toutes les cellules immunitaires dérivent des \cle{cellules souches hématopoïétiques} de la \cle{moelle osseuse rouge}.
\item \textbf{Distinguer formation et maturation} :
      \begin{itemize}
      \item lymphocytes \cle{B} : maturation dans la \textbf{moelle osseuse} ;
      \item lymphocytes \cle{T} : maturation dans le \textbf{thymus} ;
      \item polynucléaires et monocytes : maturation dans la moelle osseuse ; les monocytes deviennent \textbf{macrophages} dans les tissus.
      \end{itemize}
\item \textbf{Exploiter les documents} : greffe de moelle, thymectomie, irradiation $\rightarrow$ raisonner sur les déficits observés.
\item \textbf{Conclure} : la cellule mature migre vers les organes lymphoïdes secondaires (ganglions, rate) pour la rencontre antigénique.
\end{enumerate}
\textbf{Réflexe Bac} : ne jamais écrire « les lymphocytes T sont produits dans le thymus » : ils y \emph{mûrissent} mais sont \emph{produits} dans la moelle osseuse.
\end{methode}

\begin{exoresolu}[Thymectomie et déficit immunitaire chez la souris]
On réalise chez deux lots de souris nouveau-nées : Lot A : ablation du thymus (thymectomie) ; Lot B : témoins. Trois semaines plus tard, greffe de peau allogénique (autre souche) et mesure de la production d'anticorps après injection de sérum étranger. Résultats : Lot A accepte la greffe (pas de rejet) mais produit très peu d'anticorps ; Lot B rejette la greffe en 12 jours et produit normalement des anticorps. Interpréter.
\tcblower
\textbf{Raisonnement.} La thymectomie supprime le thymus sans toucher à la moelle : on teste le rôle du thymus dans la maturation des LT.

\textbf{Étape 1 — Lot A.} Absence de rejet $\Rightarrow$ réponse à \cle{médiation cellulaire} abolie (assurée par LT4 auxiliaires et LT8 cytotoxiques). Production d'anticorps très réduite $\Rightarrow$ l'activation complète des LB nécessite l'aide des LT4 (interleukines).

\textbf{Étape 2 — Lot B.} Rejet + anticorps normaux $\Rightarrow$ les deux réponses spécifiques fonctionnent avec un thymus intact.

\textbf{Étape 3 — Interprétation.} Le thymus est le lieu de \emph{maturation} des LT (acquisition TCR, tolérance du soi). La moelle, intacte chez A, produit précurseurs et LB $\rightarrow$ anticorps résiduels.

\textbf{Conclusion.} Toutes les cellules immunitaires sont \emph{formées} dans la moelle osseuse rouge ; les LB y \emph{mûrissent}, les LT migrent vers le \textbf{thymus} pour mûrir. Le thymus est un organe lymphoïde primaire indispensable à l'immunité cellulaire.
\end{exoresolu}

\courssub{La reconnaissance du soi et du non-soi}

\begin{definition}[Marqueurs d'identité et récepteurs]
\begin{itemize}
\item \textbf{Soi} : molécules du \cle{CMH} (Complexe Majeur d'Histocompatibilité), appelées \textbf{HLA} chez l'Homme. 
      \begin{itemize}
      \item \textbf{CMH classe I} : toutes les cellules nucléées (présentent peptides intracellulaires aux LT8/CD8).
      \item \textbf{CMH classe II} : cellules présentatrices d'antigènes professionnelles (CPA : macrophages, LB, cellules dendritiques) ; présentent peptides extracellulaires aux LT4/CD4.
      \end{itemize}
\item \textbf{Non-soi} : \cle{antigènes} (molécules étrangères : protéines virales/bactériennes, CMH de greffon, toxines). La partie reconnue est l'\cle{épitope} (déterminant antigénique).
\item \textbf{Récepteurs de reconnaissance} :
      \begin{itemize}
      \item \cle{TCR} (Récepteur des Lymphocytes T) : reconnaît un complexe \emph{peptide + CMH} (restriction du CMH).
      \item \cle{BCR} (Récepteur des Lymphocytes B = anticorps membranaire IgM/IgD) : reconnaît l'antigène \emph{natif} (conformation 3D), libre ou à la surface d'un pathogène.
      \item \cle{Anticorps solubles} (immunoglobulines) : même spécificité que le BCR, effecteurs de l'immunité humorale.
      \end{itemize}
\end{itemize}
\end{definition}

\begin{propriete}[Spécificité et complémentarité]
La reconnaissance repose sur une \textbf{complémentarité stéréochimique} (modèle clé-serrure) entre un épitope unique et le site de liaison d'un récepteur (TCR ou BCR). Un clone lymphocytaire ne reconnaît qu'un seul épitope. La diversité du répertoire ($>10^7$ spécificités) est générée par \textbf{recombinaison V(D)J} lors de la maturation (brassage génétique).
\end{propriete}

\courssub{Méthode 2 : Identifier les structures de reconnaissance du soi et du non-soi}

\begin{methode}[Reconnaître soi et non-soi dans un exercice]
\begin{enumerate}
\item \textbf{Repérer les partenaires} : cellule immunitaire (LT, LB, macrophage) et cible.
\item \textbf{Marqueurs du soi} : protéines du \cle{CMH} (HLA). Classe I = toutes cellules nucléées ; Classe II = CPA.
\item \textbf{Marqueurs du non-soi} : \cle{antigènes} (protéines étrangères, CMH greffon) ; partie reconnue = \cle{épitope}.
\item \textbf{Structures de reconnaissance} : \cle{TCR} (LT), \cle{BCR} (LB), anticorps solubles.
\item \textbf{Raisonner en complémentarité} : épitope donné $\leftrightarrow$ clone porteur du récepteur complémentaire.
\item \textbf{Conséquence} : non-soi $\Rightarrow$ activation clonale $\Rightarrow$ réponse ; soi seul $\Rightarrow$ tolérance (absence de réaction).
\end{enumerate}
\end{methode}

\begin{exoresolu}[Greffes de peau entre souches de souris]
Exp. 1 : greffe même souche pure $\rightarrow$ acceptée. Exp. 2 : greffe souche X sur Y $\rightarrow$ rejet en 12 j. Exp. 3 : re-greffe X sur même Y $\rightarrow$ rejet en 6 j. Expliquer.
\tcblower
\textbf{Exp. 1.} CMH identiques $\rightarrow$ greffon reconnu \textbf{soi}. Tolérance, pas d'activation.

\textbf{Exp. 2.} CMH X $\neq$ CMH Y $\rightarrow$ CMH X = \textbf{antigène} (non-soi) pour Y. \cle{TCR} des LT de Y reconnaissent épitopes CMH X. LT8 cytotoxiques détruisent greffon : \textbf{réponse primaire} (lente : sélection + amplification clonale = 12 j).

\textbf{Exp. 3.} Rejet accéléré (6 j) = \textbf{réponse secondaire}. Première greffe a généré \cle{lymphocytes mémoires} anti-CMH X. Reconnaissance immédiate, intense. Confirme spécificité et mémoire.

\textbf{Conclusion.} Soi = CMH ; non-soi reconnu par TCR/BCR. Réponse spécifique, mémorisée = immunité adaptative.
\end{exoresolu}

\courssub{La réponse immunitaire non spécifique (rappels)}

\begin{aretenir}[Immunité innée : premières lignes de défense]
\begin{enumerate}
\item \textbf{Barrières physiques/chimiques} : peau (kératine, pH acide, flore), muqueuses (cilia, mucus, lysozyme, pH gastrique).
\item \textbf{Inflammation} : signe clinique (rougeur, chaleur, gonflement, douleur). Médiateurs : histamine (mastocytes), prostaglandines, cytokines (IL-1, TNF-$\alpha$). Vasodilatation, perméabilité vasculaire $\rightarrow$ extravasation des polynucléaires neutrophiles (chimiotaxie).
\item \textbf{Phagocytose} : neutrophiles (premiers arrivés, vie courte), macrophages (résidents, longue vie, présentation d'antigène), cellules dendritiques (pont vers adaptative). Étapes : chémotaxie, adhérence, ingestion (phagosome), fusion lysosome $\rightarrow$ phagolysosome, digestion enzymatique (lysozyme, ROS, NO).
\item \textbf{Cellules NK (Natural Killers)} : lysent cellules cancéreuses ou infectées \emph{sans} sensibilisation préalable (reconnaissance « missing self » : absence CMH I). Perforines + granzymes $\rightarrow$ apoptose.
\item \textbf{Système du complément} : cascade protéique (voie classique, alternative, lectine) $\rightarrow$ opsonisation (C3b), chémotaxie (C5a), lyse directe (Complexe d'Attaque Membranaire C5b-9).
\end{enumerate}
\end{aretenir}

\courssub{La réponse immunitaire spécifique : médiation cellulaire et humorale}

\begin{experience}[Mise en évidence des deux médiation (transfert adoptif)]
\textbf{Protocole classique (souris).} Souris immunisée (antigène X) $\rightarrow$ prélèvement \textbf{sérum} (anticorps) OU \textbf{lymphocytes}. Transfert vers souris naïves $\rightarrow$ challenge antigène X.
\begin{itemize}
\item Transfert \textbf{sérum} protecteur $\Rightarrow$ \cle{médiation humorale} (anticorps contre antigènes libres : bactéries, toxines, virus extracellulaires).
\item Transfert \textbf{lymphocytes} protecteur $\Rightarrow$ \cle{médiation cellulaire} (LT8 cytotoxiques contre cellules infectées/cancéreuses/greffons ; LT4 auxiliaires orchestrent).
\end{itemize}
\end{experience}

\begin{popfigure}
\begin{center}
\begin{tikzpicture}[
  node distance=1.0cm and 1.2cm,
  every node/.style={align=center, font=\small},
  apc/.style={ellipse, draw=popblue, thick, fill=popblueL, minimum width=2.5cm},
  lymph/.style={circle, draw=poppurple, thick, fill=poppurpleL, minimum width=2cm},
  eff/.style={rectangle, draw=popgreen, thick, fill=popgreenL, rounded corners, minimum width=2.5cm},
  arrow/.style={-Stealth, thick, popdark},
  cytokine/.style={-Stealth, dashed, poporange, thick}
]
% APC
\node[apc, align=center] (apc){Cellule Présentatrice\\ d'Antigène (CPA)\\(Macrophage, Cell. dendritique, LB)\\\textbf{CMH II + Peptide}};
% Th
\node[lymph, right=of apc, align=center] (th){LT4 naïf (CD4+)\\\textbf{TCR + CD4}};
% Activation
\draw[arrow] (apc) -- node[above, font=\tiny] {Reconnaissance TCR/(Peptide-CMH II) + CD4 co-récepteur} (th);
\draw[cytokine] (th) -- ++(0,1.2) -| node[above, font=\tiny, pos=0.5] {IL-2, IFN-$\gamma$, IL-4...} (apc); % autocrine/paracrine

% Branche Humorale
\node[lymph, below right=of th, align=center] (lb){LB naïf\\\textbf{BCR (IgM/D) + CMH II}};
\draw[arrow] (th) -- node[above left, font=\tiny] {Contact CD40L/CD40 + Cytokines (IL-4, IL-5)} (lb);
\node[eff, below=of lb, align=center] (plasma){Plasmocyte\\ \textbf{Usine à anticorps (IgG, IgA, IgE...)}};
\node[eff, right=of plasma] (memB) {LB Mémoire};
\draw[arrow] (lb) -- node[left, font=\tiny] {Différenciation} (plasma);
\draw[arrow] (lb) -- (memB);

% Branche Cellulaire
\node[lymph, above right=of th, align=center] (tc){LT8 naïf (CD8+)\\\textbf{TCR + CD8}};
\draw[arrow] (th) -- node[above right, font=\tiny] {IL-2} (tc);
\node[eff, above=of tc, align=center] (ctlt){LT8 Cytotoxique (CTL)\\\textbf{Perforines, Granzymes, Fas/FasL}};
\node[eff, right=of ctlt] (memT) {LT Mémoire (CD4/CD8)};
\draw[arrow] (tc) -- (ctlt);
\draw[arrow] (tc) -- (memT);

% Cible Cellulaire
\node[apc, right=of ctlt, fill=poporangeL, draw=poporange, align=center] (cible){Cellule Cible Infectée\\ ou Cancéreuse\\ \textbf{CMH I + Peptide viral/tumoral}};
\draw[arrow, poporange] (ctlt) -- node[above, font=\tiny] {Reconnaissance TCR/(Peptide-CMH I) + CD8 co-récepteur} (cible);
\draw[arrow, poporange, bend left=20] (cible) to node[above, font=\tiny] {Apoptose} (ctlt);

% Legendes boites
\node[draw=popblue, fit=(apc), inner sep=2pt, label={[popblue]above left:\textbf{Présentation}}] {};
\node[draw=poppurple, fit=(th)(lb)(tc), inner sep=2pt, label={[poppurple]above left:\textbf{Activation / Sélection clonale}}] {};
\node[draw=popgreen, fit=(plasma)(memB)(ctlt)(memT), inner sep=2pt, label={[popgreen]above left:\textbf{Phase Effectrice}}] {};
\node[draw=poporange, fit=(cible), inner sep=2pt, label={[poporange]above left:\textbf{Cible}}] {};
\end{tikzpicture}
\end{center}
\legende{Schéma unifié de la réponse spécifique : présentation par CPA (CMH II) $\rightarrow$ activation LT4 $\rightarrow$ orchestration double : humorale (LB $\rightarrow$ Plasmocytes/Ac) et cellulaire (LT8 $\rightarrow$ CTL). Les flèches pleines = contact/cytokinique ; flèches pointillées = cytokines.}
\end{popfigure}

\begin{propriete}[Mécanismes effecteurs]
\begin{itemize}
\item \textbf{Médiation humorale (anti-Ag libres)} : Anticorps (IgG, IgM, IgA, IgE) $\rightarrow$ Neutralisation (blocage entrée virus/toxines), Opsonisation (phagocytose facilitée), Activation complément (lyse bactéries), ADCC (NK tuent cellule opsonisée via Fc$\gamma$R).
\item \textbf{Médiation cellulaire (anti-cellules altérées)} : LT8 cytotoxiques (CTL) reconnaissent \textbf{CMH I + peptide} sur cellule cible $\rightarrow$ lysis par \textbf{perforines} (pores) + \textbf{granzymes} (apoptose) ou voie \textbf{Fas/FasL}. LT4 Th1 activent macrophages (IFN-$\gamma$) pour tuer bactéries intracellulaires (ex. \emph{Mycobacterium tuberculosis}).
\end{itemize}
\end{propriete}

\courssub{Méthode 3 : Interpréter des résultats d'expériences sur les types de réponses immunitaires}

\begin{methode}[Analyser un protocole expérimental d'immunologie]
\begin{enumerate}
\item \textbf{Problème \& Hypothèse} : ex. « Immunité transmise par cellules ou sérum ? ».
\item \textbf{Lots} : Témoin / Expérimental(s) ; identifier la \emph{seule} variable modifiée.
\item \textbf{Comparaison quantitative} : citer chiffres \textbf{avec unités} (survie \%, titres Ac, délais).
\item \textbf{Attribution mécanistique} :
      \begin{itemize}
      \item Transfert \textbf{sérum} protecteur $\Rightarrow$ \cle{médiation humorale} (LB $\rightarrow$ Plasmocytes $\rightarrow$ Ac).
      \item Transfert \textbf{cellules} (lymphocytes) protecteur $\Rightarrow$ \cle{médiation cellulaire} (LT4, LT8).
      \item Destruction cellules infectées/greffées $\Rightarrow$ LT8 cytotoxiques.
      \item Protection immédiate, brève, sans mémoire après sérum $\Rightarrow$ \textbf{immunisation passive}.
      \end{itemize}
\item \textbf{Conclusion} : valider/infirmer hypothèse sans sur-généraliser.
\end{enumerate}
\textbf{Réflexe} : l'expérience de transfert (sérum ou cellules d'immunisé vers naïf) est le schéma classique pour distinguer humorale et cellulaire.
\end{methode}

\begin{exoresolu}[Transfert de sérum et de lymphocytes]
Cobayes guéris (immunisés) $\rightarrow$ sérum + lymphocytes. 3 lots naïfs : Lot 1 (sérum), Lot 2 (lymphocytes), Lot 3 (sérum physiologique témoin). Infection bactérie. Survie à J15 : L1=90\%, L2=20\%, L3=10\%. Interpréter.
\tcblower
\textbf{Hypothèse} : Immunité anti-bactérienne = facteurs sériques (Ac) ou cellules ?

\textbf{Lot 3 (Témoin)} : 10\% survie $\rightarrow$ bactérie pathogène, naïfs sans défense.

\textbf{Lot 1 (Sérum)} : 90\% survie $\rightarrow$ protection quasi totale. Sérum = \cle{anticorps} (immunoglobulines) des plasmocytes. Mécanismes : neutralisation, agglutination, opsonisation, activation complément.

\textbf{Lot 2 (Lymphocytes)} : 20\% $\approx$ témoin. Lymphocytes seuls inefficaces contre bactérie \emph{extracellulaire} (de plus, risque rejet allogénique si transfert entre individus).

\textbf{Interprétation} : Protection par \emph{sérum} = \cle{réponse à médiation humorale} (LB $\rightarrow$ Plasmocytes $\rightarrow$ Ac). Immunisation \textbf{passive} : immédiate, temporaire, pas de mémoire chez receveur.

\textbf{Conclusion} : Immunité acquise contre cette bactérie = type humoral. Anticorps seuls assurent la protection.
\end{exoresolu}

\courssub{La mémoire immunitaire et la vaccination}

\begin{propriete}[Cinétique de la réponse humorale : Primaire vs Secondaire]
\begin{center}
\begin{tabularx}{\linewidth}{|l|X|X|}\hline
\textbf{Paramètre} & \textbf{Réponse Primaire (1\iere{} contact)} & \textbf{Réponse Secondaire (Rappel, même Ag)} \\ \hline
\textbf{Latence} & Longue (5--10 jours) & Courte (1--3 jours) \\ \hline
\textbf{Pic de concentration} & Faible & Élevé (10 à 100$\times$ primaire) \\ \hline
\textbf{Ig dominantes} & IgM (puis IgG) & \textbf{IgG} (affinité mature) \\ \hline
\textbf{Durée / Persistance} & Courte (semaines) & Longue (mois/années) \\ \hline
\textbf{Base cellulaire} & Sélection clonale $\rightarrow$ Plasmocytes + \textbf{Lymphocytes Mémoire} (LBm, LTm) & Activation immédiate \textbf{Lymphocytes Mémoire} (nombreux, affinité haute, seuil activation bas) \\ \hline
\end{tabularx}
\end{center}
\end{propriete}

\courssub{Méthode 4 : Exploiter une courbe de cinétique des anticorps (mémoire immunitaire)}

\begin{methode}[Lire et interpréter une réponse primaire et secondaire]
\begin{enumerate}
\item \textbf{Axes} : Abscisse = Temps (jours) + dates injections ; Ordonnée = Titre Ac (conc. ou UA).
\item \textbf{Primaire} : Latence, pic \emph{faible/lent} (J10--15), IgM dominantes.
\item \textbf{Secondaire} : Même Ag injecté $\rightarrow$ Latence \emph{courte} (J2--3), pic \emph{élevé/durable}, IgG dominantes.
\item \textbf{Explication cellulaire} : 1\iere{} rencontre $\rightarrow$ sélection clonale + plasmocytes + \cle{lymphocytes mémoires}. 2\ieme{} rencontre $\rightarrow$ activation immédiate cellules mémoires.
\item \textbf{Lien vaccination} : Vaccin (Ag inoffensif) crée mémoire sans maladie ; rappel entretient cellules mémoires.
\item \textbf{Spécificité} : Ag différent au rappel $\rightarrow$ réponse primaire (pas de mémoire croisée).
\end{enumerate}
\end{methode}

\begin{exoresolu}[Cinétique anticorps après deux injections]
Lapin : Anatoxine tétanique J0 et J30. Titre (UA) : J0=0 ; J7=5 ; J15=40 (pic) ; J30=10 ; J33=60 ; J38=400 (pic) ; J60=250. Tracer et interpréter.
\tcblower
\textbf{Tracé.} Temps (jours) / Titre (UA). Courbe : pic modéré J15 (40 UA) ; 2\ieme{} injection J30 $\rightarrow$ montée rapide, pic élevé J38 (400 UA), décroissance lente.

\begin{popfigure}
\begin{center}
\begin{tikzpicture}
\begin{axis}[width=0.9\linewidth,height=6cm,
xlabel={Temps (jours)},ylabel={Titre d'anticorps (UA)},
xmin=0,xmax=65,ymin=0,ymax=450,
xtick={0,7,15,30,33,38,60},
grid=both,grid style={popline!40},
legend pos=north west]
\addplot[popcurve,smooth,thick] coordinates {(0,0) (7,5) (15,40) (22,25) (30,10) (33,60) (38,400) (45,330) (60,250)};
\addplot[poporange,dashed,thick] coordinates {(30,0) (30,450)};
\legend{Réponse anticorps, 2\ieme{} injection};
\end{axis}
\end{tikzpicture}
\end{center}
\legende{Cinétique des anticorps antitétaniques : réponse primaire (pic faible/tardif) puis secondaire (pic rapide/élevé) après rappel.}
\end{popfigure}

\textbf{Primaire (J0--J30).} Latence (5 UA J7), pic faible 40 UA J15. Temps sélection clone LB, amplification, différenciation plasmocytes. Apparition \cle{lymphocytes mémoires}.

\textbf{Secondaire (J30--J60).} Latence courte (60 UA J33 = 3 j), pic 10$\times$ plus haut (400 UA J38), persistant (250 UA J60). Cellules mémoires (spécifiques, nombreuses, seuil bas) activées immédiatement $\rightarrow$ production massive IgG.

\textbf{Conclusion.} Mémoire immunitaire = réponse secondaire plus rapide, intense, durable. Fondement de la \textbf{vaccination} (anatoxine = Ag vaccinal non pathogène).
\end{exoresolu}

\begin{aretenir}[Principes et types de vaccins]
\begin{itemize}
\item \textbf{Principe} : Mimer l'infection (Ag) sans provoquer la maladie $\rightarrow$ induire mémoire (LBm, LTm).
\item \textbf{Types} : 
      \begin{itemize}
      \item \textbf{Vivants atténués} (Rougeole, Rubéole, Oreillons, BCG, Fièvre jaune) : réplication limitée, immunité forte/durable (humorale + cellulaire), contre-indiqués immunodéprimés/grossesse.
      \item \textbf{Inactivés (tués)} (Coqueluche, Polio Salk, Rage, Hépatite A, Grippe) : pas de réplication, plus sûrs, souvent adjuvés (Alum), rappels nécessaires.
      \item \textbf{Sous-unitaires / Recombinants / ARN messager} (Hépatite B, HPV, COVID-19 ARNm, Pneumocoque conjugué) : Ag purifié ou code génétique $\rightarrow$ sécurité maximale, réponse ciblée.
      \item \textbf{Toxoïdes (Anatoxines)} (Tétanos, Diphtérie) : toxine inactivée $\rightarrow$ anti-toxines neutralisantes.
      \end{itemize}
\item \textbf{Couverture vaccinale} : Objectif OMS $\ge$ 90\% (immunité collective protège non-vaccinés).
\end{itemize}
\end{aretenir}

\courssub{Les dysfonctionnements : maladies auto-immunes}

\begin{definition}[Maladie auto-immune]
Rupture de la \textbf{tolérance du soi} (défaillance sélection négative thymique/médullaire ou régulation périphérique Treg) $\rightarrow$ activation de clones \textbf{autoréactifs} : production d'\textbf{auto-anticorps} et/ou \textbf{LT autoréactifs} dirigés contre antigènes propres $\rightarrow$ lésion tissulaire chronique.
\end{definition}

\begin{propriete}[Trois exemples au programme]
\begin{center}
\begin{tabularx}{\linewidth}{|l|X|X|X|}\hline
\textbf{Maladie} & \textbf{Cible auto-immune} & \textbf{Mécanisme lésionnel} & \textbf{Symptômes majeurs} \\ \hline
\textbf{Diabète juvénile (Type 1)} & Cellules $\beta$ des îlots de Langerhans (pancréas endocrine) & LT8 cytotoxiques + Auto-Ac anti-GAD/Insuline/IA-2 $\rightarrow$ destruction cellulaire & \textbf{Insulino-privation} $\rightarrow$ Hyperglycémie chronique ($>1,26$ g/L à jeun), glycosurie, polyurie, polydipsie, amaigrissement, cétoacidose. \\ \hline
\textbf{Myasthénie} & Récepteurs à l'acétylcholine (RACh) plaque motrice (jonction neuromusculaire) & \textbf{Auto-Ac anti-RACh} (IgG) $\rightarrow$ blocage/complément/destruction RACh $\rightarrow$ transmission neuromusculaire défaillante & \textbf{Fatigabilité musculaire} : ptosis, diplopie, dysphagie, faiblesse proximale, aggravée par l'effort, améliorée par le repos. \\ \hline
\textbf{Vitiligo} & Mélanocytes (cellules pigmentaires peau/cheveux) & LT8 cytotoxiques autoréactifs (réaction cellulaire) + Auto-Ac anti-mélanocytes $\rightarrow$ destruction mélanocytes & \textbf{Taches dépigmentées} (blanches, symétriques souvent), zones photo-exposées, muqueuses, cheveux blancs (leucotrichie). Non contagieux, impact psychologique fort. \\ \hline
\end{tabularx}
\end{center}
\end{propriete}

\courssub{Méthode 5 : Expliquer le mécanisme d'une maladie auto-immune}

\begin{methode}[Construire l'explication d'une maladie auto-immune]
\begin{enumerate}
\item \textbf{Définir} : rupture tolérance soi $\rightarrow$ auto-anticorps et/ou LT autoréactifs.
\item \textbf{Identifier la cible} (cellules $\beta$, RACh, mélanocytes...).
\item \textbf{Décrire le mécanisme} : Auto-Ac bloquants/destructeurs (complément, ADCC) OU LT8 cytotoxiques (perforines/granzymes).
\item \textbf{Relier lésion $\rightarrow$ symptômes} : disparition/blocage cible = signes cliniques.
\item \textbf{Citer exemples programme} : Diabète juvénile, Myasthénie, Vitiligo.
\end{enumerate}
\end{methode}

\begin{exoresolu}[Le diabète juvénile, une maladie auto-immune]
Adolescent 13 ans (CHU Essos, Yaoundé) : soif intense, polyurie, amaigrissement. Glycémie jeun = \SI{2,8}{g/L} (norme \num{0,7}--\SI{1,1}{g/L}). Insulinémie quasi nulle. Auto-Ac anti-cellules $\beta$ détectés. Expliquer mécanisme et justifier traitement insuline.
\tcblower
\textbf{Diagnostic.} Hyperglycémie chronique (\SI{2,8}{g/L}) + insulinopénie = \textbf{Diabète Type 1 (juvénile insulinodépendant)}.

\textbf{Mécanisme auto-immun.} Auto-Ac anti-cellules $\beta$ = preuve rupture tolérance. Système immunitaire reconnaît cellules $\beta$ propres comme non-soi. \textbf{LT8 cytotoxiques autoréactifs} + Auto-Ac $\rightarrow$ destruction progressive cellules $\beta$ (insulite).

\textbf{Conséquence physiologique.} Cellules $\beta$ uniques productrices d'\cle{insuline} (hormone hypoglycémiante : entrée glucose cellules, stockage glycogène/graisses). Destruction $\Rightarrow$ \textbf{absence insuline} $\Rightarrow$ glucose s'accumule sang (hyperglycémie) $\rightarrow$ seuil rénal dépassé $\rightarrow$ glycosurie $\rightarrow$ perte eau/électrolytes (polyurie, polydipsie) + catabolisme graisses/protéines (amaigrissement, corps cétoniques).

\textbf{Justification traitement.} Pancréas ne produit plus insuline $\rightarrow$ apport exogène obligatoire. \textbf{Injections quotidiennes d'insuline} (protéine digérée si voie orale) + automesure glycémie + régime équilibré.

\textbf{Conclusion.} Diabète juvénile = auto-immun (destruction cellules $\beta$ par auto-Ac/LT autoréactifs) $\rightarrow$ insulinopénie $\rightarrow$ hyperglycémie chronique traitée par insulinothérapie substitutive à vie.
\end{exoresolu}

\courssub{Le VIH/Sida : mécanismes, prévention et solidarité}

\begin{experience}[Cycle de réplication du VIH-1 (Rétrovirus)]
\begin{enumerate}
\item \textbf{Fixation} : Glycoprotéine virale \textbf{gp120} se lie au récepteur \textbf{CD4} (LT4, macrophages, cellules dendritiques) + co-récepteur \textbf{CCR5} (souches R5, transmission) ou \textbf{CXCR4} (souches X4, stade avancé).
\item \textbf{Entrée} : Fusion membrane virale / membrane cellulaire (gp41) $\rightarrow$ libération capside dans cytoplasme.
\item \textbf{Transcription inverse} : ARN viral $(+)$ $\xrightarrow{\text{\cle{Transcriptase Inverse (TI)}}}$ ADN double brin (provinus). \textbf{Étape cible thérapeutique majeure (I-NRTI, I-NNRTI)}.
\item \textbf{Intégration} : ADN viral $\xrightarrow{\text{Intégrase}}$ intégration dans génome hôte (chromosome). \textbf{Cible : Inhibiteurs d'Intégrase (II)}.
\item \textbf{Transcription/Traduction} : Machinery cellulaire $\rightarrow$ ARNm viraux $\rightarrow$ protéines (Gag, Pol, Env) + ARN génomique.
\item \textbf{Assemblage / Bourgeonnement} : Particules immatures $\rightarrow$ maturation par \textbf{Protéase virale} (clivage polyprotéines) $\rightarrow$ Virions matures infectieux. \textbf{Cible : Inhibiteurs de Protéase (IP)}.
\item \textbf{Destruction cellulaire} : Lyse, apoptose, syncytia, mort par activation immunitaire (bystander effect).
\end{enumerate}
\end{experience}

\begin{popfigure}
\begin{center}
\begin{tikzpicture}[
  node distance=1.0cm and 1.0cm,
  every node/.style={align=center, font=\small},
  stepbox/.style={rectangle, draw=popblue, thick, fill=popblueL, rounded corners, minimum width=3cm, minimum height=1cm},
  arrow/.style={-Stealth, thick, popdark},
  target/.style={rectangle, draw=poporange, thick, fill=poporangeL, rounded corners, minimum width=3.5cm, minimum height=0.8cm, font=\tiny}
]
\node[stepbox] (s1) {1. Fixation : gp120 -- CD4 + CCR5/CXCR4};
\node[stepbox, below=of s1] (s2) {2. Fusion / Entrée};
\node[stepbox, below=of s2, align=center] (s3){3. Transcription Inverse\\ (ARN $\rightarrow$ ADN)};
\node[stepbox, below=of s3, align=center] (s4){4. Intégration\\ (ADN viral dans génome hôte)};
\node[stepbox, below=of s4, align=center] (s5){5. Transcription / Traduction\\ (Protéines virales + ARN)};
\node[stepbox, below=of s5] (s6) {6. Assemblage / Bourgeonnement};
\node[stepbox, below=of s6, align=center] (s7){7. Maturation (Protéase)\\ $\rightarrow$ Virions infectieux};

\draw[arrow] (s1) -- (s2);
\draw[arrow] (s2) -- (s3);
\draw[arrow] (s3) -- (s4);
\draw[arrow] (s4) -- (s5);
\draw[arrow] (s5) -- (s6);
\draw[arrow] (s6) -- (s7);

% Cibles thérapeutiques
\node[target, right=of s3, xshift=1.5cm, align=center] (t1){Cible : \textbf{I-NRTI / I-NNRTI}\\(ex. Ténofovir, Efavirenz)};
\node[target, right=of s4, xshift=1.5cm, align=center] (t2){Cible : \textbf{Inhibiteurs Intégrase}\\(ex. Dolutégravir, Raltégravir)};
\node[target, right=of s7, xshift=1.5cm, align=center] (t3){Cible : \textbf{Inhibiteurs Protéase}\\(ex. Lopinavir, Darunavir)};
\draw[poporange, dashed, thick] (s3.east) -- (t1.west);
\draw[poporange, dashed, thick] (s4.east) -- (t2.west);
\draw[poporange, dashed, thick] (s7.east) -- (t3.west);
\end{tikzpicture}
\end{center}
\legende{Cycle de réplication du VIH et cibles des antirétroviraux (Trithérapie = association 3 molécules ciblant au moins 2 étapes différentes).}
\end{popfigure}

\begin{propriete}[Physiopathologie du Sida]
\begin{itemize}
\item \textbf{Cible principale} : \textbf{LT4 (CD4+)} = « chefs d'orchestre » (sécrètent IL-2, IFN-$\gamma$, IL-4... activant LB, LT8, macrophages).
\item \textbf{Dynamique} : Infection $\rightarrow$ réplication massive $\rightarrow$ destruction LT4 continue. Phase asymptomatique (années) : équilibre précaire production/destruction. Chute LT4 $<$ \SI{200}{\per\micro\liter} (norme \SIrange{800}{1200}{\per\micro\liter}) $\rightarrow$ \textbf{Stade Sida} (Syndrome d'ImmunoDéficience Acquise).
\item \textbf{Conséquence} : Effondrement \textbf{des deux réponses spécifiques} (humorale + cellulaire).
\item \textbf{Infections opportunistes} : Agents habituellement contrôlés (TB : \emph{Mycobacterium tuberculosis}, Pneumocystose : \emph{Pneumocystis jirovecii}, Candidoses, Toxoplasmose, CMV, Herpès, Kaposi HHV-8).
\end{itemize}
\end{propriete}

\courssub{Méthode 6 : Analyser l'action du VIH et produire un outil de sensibilisation}

\begin{methode}[Expliquer l'immunodéficience VIH et concevoir un message de prévention]
\begin{enumerate}
\item \textbf{Cible} : VIH (rétrovirus ARN) infecte \textbf{LT4 (CD4+)} et macrophages (récepteur CD4 + co-récepteurs CCR5/CXCR4).
\item \textbf{Cycle} : Fixation CD4 $\rightarrow$ Entrée $\rightarrow$ \textbf{Transcriptase inverse} (ARN $\rightarrow$ ADN) $\rightarrow$ Intégration $\rightarrow$ Production virions $\rightarrow$ Destruction LT4.
\item \textbf{Déficit} : LT4 = chefs d'orchestre (activent LB, LT8). Chute LT4 (norme \SIrange{800}{1200}{\per\micro\liter} ; Sida $<$ \SI{200}{\per\micro\liter}) $\rightarrow$ effondrement immunité spécifique.
\item \textbf{Conséquences} : Infections opportunistes (TB, Pneumocystose, Candidoses) + Cancers (Kaposi). \textbf{Sida = stade de la maladie, pas maladie distincte}.
\item \textbf{Outil de sensibilisation} : Support adapté (affiche, sketch, radio). Message clair :
      \begin{itemize}
      \item \textbf{Transmission} : Rapports sexuels non protégés, Sang contaminé (aiguilles, transfusions), Mère-enfant (grossesse, accouchement, allaitement).
      \item \textbf{Prévention} : Préservatif (M/F), Dépistage (connaître statut), Matériel stérile/usage unique, Trithérapie mère séropositive (PMTCT), PrEP/PEP.
      \item \textbf{Solidarité} : VIH \textbf{NE se transmet PAS} par poignée de main, repas, moustiques, piscine, toux. Discrimination = injustice, frein au dépistage/soins. « Vivre avec le VIH, pas le rejeter ».
      \end{itemize}
\end{enumerate}
\end{methode}

\begin{exoresolu}[Suivi biologique VIH et messages de campagne]
Patiente (Centre santé Douala) : LT4 : An 0 (contamination) \SI{1000}{\per\micro\liter} ; An 1 \num{650} ; An 4 \num{400} ; An 7 \SI{180}{\per\micro\liter} + tuberculose + candidoses. Expliquer évolution et proposer 2 messages prévention + 1 solidarité pour campagne scolaire.
\tcblower
\textbf{Étape 1 — Courbe LT4.} Chute progressive \SI{1000}{\per\micro\liter} $\rightarrow$ \SI{180}{\per\micro\liter} en 7 ans. VIH se réplique dans LT4 et les détruit en continu. Phase « silencieuse » (asymptomatique) = épuisement progressif stock LT4.

\textbf{Étape 2 — Immunologie.} LT4 sécrètent interleukines activant LB (Ac) et LT8 (cytotoxiques). $<$ \SI{200}{\per\micro\liter} $\rightarrow$ paralysie deux bras immunité spécifique = \textbf{Stade Sida}. TB et candidoses = \textbf{infections opportunistes} (germes commensaux/faiblement pathogènes non contrôlés).

\textbf{Étape 3 — Prise en charge.} \textbf{Trithérapie} (ARV : 2 I-NRTI + 1 I-NNRTI ou II ou IP) bloque réplication, permet remontée LT4, charge virale indétectable ($<$ \SI{50}{copies/mL}) $\rightarrow$ infection chronique contrôlée. Dépistage précoce = vital.

\textbf{Étape 4 — Messages campagne scolaire.}
\begin{itemize}
\item \textbf{Prévention 1} : « Protège ton avenir : le préservatif à \emph{chaque} rapport sexuel à risque, c'est la barrière qui arrête le VIH. »
\item \textbf{Prévention 2} : « Fais le test : connaître ton statut, c'est te protéger et protéger les autres. Jamais d'aiguille, de lame, de matériel de scarification partagé. »
\item \textbf{Solidarité} : « Le VIH ne se transmet ni par la main, ni par l'assiette, ni par les moustiques. Accueillons, soutenons, aimons les personnes vivant avec le VIH : la stigmatisation tue plus vite que le virus. »
\end{itemize}
\textbf{Conclusion.} VIH détruit LT4 (coordonnateurs immunité) $\rightarrow$ $<$ \SI{200}{\per\micro\liter} = Sida (infections opportunistes). Trois piliers : \textbf{Prévention} (préservatif, dépistage, matériel stérile), \textbf{Trithérapie} (vie longue/qualité), \textbf{Solidarité} (lutte contre stigmatisation).
\end{exoresolu}

\begin{attention}[Les 5 erreurs qui coûtent des points au Bac]
\begin{enumerate}
\item \textbf{Confondre formation et maturation} : « Lymphocytes T produits dans le thymus » = \textbf{FAUX}. Produits dans moelle, mûrissent dans thymus. LB mûrissent dans moelle (pas rate).
\item \textbf{Confondre médiation humorale et cellulaire} : Humorale = Ac (plasmocytes/LB) contre Ag \emph{libres} (bactéries, toxines). Cellulaire = LT8 cytotoxiques contre cellules \emph{infectées/étrangères} (cancéreuses, greffons). \textbf{JAMAIS} : « Les anticorps détruisent la cellule greffée ».
\item \textbf{Inverser réponse primaire/secondaire} : Secondaire = plus \emph{rapide}, plus \emph{intense}, plus \emph{durable} (lymphocytes mémoires). Attribuer pic élevé à 1\iere{} injection = erreur classique lecture courbe.
\item \textbf{Oublier unités et valeurs de référence} : Glycémie en \si{g/L} (norme \num{0,7}--\SI{1,1}{g/L}) ; LT4 en \si{\per\micro\liter} (seuil Sida \SI{200}{\per\micro\liter}). Comparaison résultats sans chiffres du doc = insuffisant.
\item \textbf{Confondre VIH et Sida / Voies transmission} : On transmet le \textbf{VIRUS (VIH)} ; Sida = \textbf{STADE} de la maladie. Voies \textbf{UNIQUEMENT} : Sexuelle, Sanguine, Mère-enfant. Dire « transmission Sida » ou « VIH par moustiques » = \textbf{zéro point} en sensibilisation.
\end{enumerate}
\end{attention}

\begin{savaistu}[Le VIH au Cameroun : données et enjeux]
Selon les estimations ONUSIDA 2023, le Cameroun compte environ \num{500000} personnes vivant avec le VIH (PVVIH), avec une prévalence autour de \SI{2,7}{\%} chez les 15--49 ans. La prévention de la transmission mère-enfant (PTME) a permis de réduire la transmission verticale à moins de \SI{5}{\%} grâce à la mise sous ARV précoce des femmes enceintes séropositives. Le défi majeur reste le \textbf{dépistage} (seulement \SI{70}{\%} des PVVIH connaissent leur statut) et l'\textbf{observance} thérapeutique pour atteindre l'objectif « 95-95-95 » (95\% diagnostiqués, 95\% traités, 95\% charge virale supprimée). La lutte contre la \textbf{stigmatisation} dans les écoles et communautés est un levier essentiel pour améliorer le recours aux soins.
\end{savaistu}

\newpage

\coursec{Exercices}

\courssub{Je vérifie mes connaissances}

\begin{exercice}[QCM : Origine et reconnaissance immunitaire]
\textbf{1.} Les lymphocytes B acquièrent leur compétence immunitaire (maturation) principalement dans :
\begin{itemize}
    \item[A.] le thymus.
    \item[B.] la moelle osseuse rouge.
    \item[C.] la rate.
    \item[D.] les ganglions lymphatiques.
\end{itemize}
\textbf{2.} Le Complexe Majeur d'Histocompatibilité (CMH) de classe I est exprimé par :
\begin{itemize}
    \item[A.] uniquement les cellules présentatrices d'antigènes (CPA).
    \item[B.] toutes les cellules nucléées de l'organisme.
    \item[C.] uniquement les lymphocytes T4.
    \item[D.] les érythrocytes.
\end{itemize}
\textbf{3.} La réponse immunitaire non spécifique (innée) se caractérise par :
\begin{itemize}
    \item[A.] une spécificité antigénique élevée et une mémoire immunologique.
    \item[B.] une mise en œuvre lente (plusieurs jours) et l'absence de mémoire.
    \item[C.] une mise en œuvre rapide (minutes à heures) et l'absence de mémoire.
    \item[D.] la production d'anticorps de haute affinité.
\end{itemize}
\textbf{4.} Quel type de lymphocyte est la cible principale du VIH ?
\begin{itemize}
    \item[A.] Les lymphocytes B.
    \item[B.] Les lymphocytes T cytotoxiques (LT8 / CD8+).
    \item[C.] Les lymphocytes T auxiliaires (LT4 / CD4+).
    \item[D.] Les cellules Natural Killer (NK).
\end{itemize}
\end{exercice}

\begin{exercice}[Vrai ou Faux ? Justifiez chaque réponse]
\textbf{1.} « Les maladies auto-immunes résultent d'une absence totale de réaction immunitaire contre les antigènes du soi. »
\textbf{2.} « La mémoire immunitaire est assurée uniquement par les lymphocytes B mémoires qui produisent des anticorps en permanence. »
\textbf{3.} « Dans le diabète de type 1 (juvénile), le système immunitaire détruit les cellules $\beta$ des îlots de Langerhans du pancréas. »
\textbf{4.} « Le dépistage du VIH par test rapide (TROD) recherche directement la présence du virus dans le sang par PCR. »
\end{exercice}

\begin{exercice}[Définitions et notions clés]
Donnez une définition précise et rigoureuse des termes suivants :
\begin{enumerate}
    \item Antigène.
    \item Épitope (déterminant antigénique).
    \item CMH (Complexe Majeur d'Histocompatibilité) de classe II.
    \item Tolérance immunologique (centrale et périphérique).
    \item Infection opportuniste (dans le contexte du Sida).
\end{enumerate}
\end{exercice}

\begin{exercice}[Calcul et lecture de données : Dynamique des populations lymphocytaires]
Chez un adulte sain, le taux de lymphocytes T CD4+ dans le sang périphérique est d'environ \SI{800}{\per\micro\liter} (soit \SI{40}{\%} des lymphocytes T). Chez un patient infecté par le VIH non traité, ce taux chute à \SI{200}{\per\micro\liter} au bout de 5 ans.
\begin{enumerate}
    \item Calculez le pourcentage de diminution du taux absolu de LT4.
    \item Si le total des lymphocytes T reste constant, quel devient le pourcentage de LT4 parmi les lymphocytes T ?
    \item Rappellez le seuil de LT4 (en \SI{}{\per\micro\liter}) en dessous duquel on parle de stade Sida (immunodéficience sévère) selon la classification CDC/OMS.
\end{enumerate}
\end{exercice}

\courssub{Je m'entraîne}

\begin{exercice}[Lieux de formation et de maturation : Synthèse]
Complétez le tableau suivant en indiquant pour chaque type cellulaire : son site de formation (hématopoïèse), son site de maturation (acquisition de la compétence immunitaire) et son marqueur de surface principal (CD).
\begin{center}
\begin{tabularx}{\linewidth}{|l|X|X|X|}\hline
\textbf{Cellule} & \textbf{Site de formation} & \textbf{Site de maturation} & \textbf{Marqueur CD principal} \\ \hline
Lymphocyte B &  &  &  \\ \hline
Lymphocyte T CD4+ (LT4) &  &  &  \\ \hline
Lymphocyte T CD8+ (LT8) &  &  &  \\ \hline
Macrophage / Monocyte &  &  &  \\ \hline
Cellule dendritique &  &  &  \\ \hline
\end{tabularx}
\end{center}
\end{exercice}

\begin{exercice}[Mécanismes effecteurs : Médiation cellulaire vs Médiation humorale]
On infecte deux groupes de souris avec une bactérie intracellulaire facultative (\emph{Listeria monocytogenes}) et un virus grippal.
\begin{enumerate}
    \item Pour chaque agent pathogène, précisez quelle réponse spécifique (cellulaire ou humorale) est \textbf{prépondérante} pour l'élimination de l'agent. Justifiez par la localisation de l'agent (extra/intracellulaire) et les effecteurs mobilisés.
    \item Décrivez le mécanisme de destruction d'une cellule cible infectée par un virus par un LT8 cytotoxique (perforines, granzymes, apoptose).
    \item Expliquez pourquoi le transfert de sérum immun (anticorps) d'une souris guérie protège une souris naïve contre le virus grippal, mais pas contre \emph{Listeria}.
\end{enumerate}
\end{exercice}

\begin{exercice}[Exploitation de courbes : Réponse primaire et réponse secondaire]
On injecte un antigène X (protéine purifiée) à un lapin au jour 0 (primo-injection) et une seconde injection (rappel) au jour 30. On mesure le taux d'anticorps anti-X dans le sérum pendant 60 jours. Les résultats sont schématisés ci-dessous (courbes schématiques).
\begin{popfigure}
\begin{tikzpicture}[scale=0.7]
\begin{axis}[
    width=\linewidth, height=8cm,
    xlabel={Temps (jours)}, ylabel={Taux d'anticorps (UA/mL)},
    xmin=0, xmax=65, ymin=0, ymax=110,
    xtick={0,7,14,30,37,44,60}, ytick={0,20,40,60,80,100},
    legend pos=north west,
    grid=major, grid style={dashed, gray!30},
]
% Primary response
\addplot[popblue, thick, smooth, mark=*] coordinates {(0,0) (3,5) (7,20) (10,45) (14,60) (21,40) (30,15)};
\addlegendentry{Réponse primaire (IgM puis IgG)}
% Secondary response
\addplot[poporange, thick, smooth, mark=square*] coordinates {(30,15) (32,30) (35,80) (37,100) (44,90) (60,70)};
\addlegendentry{Réponse secondaire (IgG)}
% Vertical line for booster
\draw[dashed, popdark] (30,0) -- (30,105) node[pos=0.5, right] {Rappel (J30)};
\end{axis}
\end{tikzpicture}
\legende{Évolution du taux d'anticorps sériques après primo-injection (J0) et rappel (J30) de l'antigène X.}
\end{popfigure}
\begin{enumerate}
    \item Comparez la latence, l'amplitude du pic, la cinétique de montée et la nature des anticorps (isotypes) entre la réponse primaire et la réponse secondaire.
    \item Identifiez les cellules responsables de la rapidité et de l'intensité de la réponse secondaire.
    \item Expliquez le principe du rappel vaccinal en vous appuyant sur ce graphique.
\end{enumerate}
\end{exercice}

\begin{exercice}[Cas clinique : Myasthénie auto-immune]
Un patient de 35 ans consulte pour une fatigue musculaire anormale à l'effort, une ptôsis (chute de la paupière) et une diplopie (vision double) s'aggravant en fin de journée. Le test à la néostigmine (inhibiteur de l'acétylcholinestérase) est positif (amélioration transitoire). Le dosage sanguin révèle la présence d'anticorps anti-récepteurs à l'acétylcholine (anti-RACh).
\begin{enumerate}
    \item Identifiez la pathologie. Précisez s'il s'agit d'une maladie auto-immune à médiation humorale ou cellulaire. Justifiez.
    \item Expliquez le mécanisme physiopathologique reliant la présence de ces auto-anticorps aux symptômes musculaires (blocage de la transmission neuromusculaire).
    \item Le thymus est souvent hyperplasique ou thymomeux chez ces patients. Quel rôle joue le thymus dans la genèse de cette auto-immunité (sélection négative défaillante, mimétisme moléculaire) ?
    \item Proposez deux pistes thérapeutiques ciblant le système immunitaire (autre que l'anticholinestérasique).
\end{enumerate}
\end{exercice}

\courssub{Je me prépare au Bac}

\begin{exercice}[Type Bac : Preuve expérimentale de la médiation cellulaire dans le rejet de greffe (Expérience de Mitchison, 1954)]
On réalise des greffes de peau allogéniques (donneur : souche A, receveur : souche B) chez la souris. On suit la survie de la greffe dans trois lots de souris receveurs (souche B) :
\begin{itemize}
    \item \textbf{Lot 1 (Témoins) :} Souris normales recevant une greffe de peau souche A.
    \item \textbf{Lot 2 :} Souris irradiées (système immunitaire détruit) recevant une greffe de peau souche A \textbf{ET} une injection de lymphocytes T purifiés de souris souche B \textbf{préalablement immunisées} contre la souche A.
    \item \textbf{Lot 3 :} Souris irradiées recevant une greffe de peau souche A \textbf{ET} une injection de sérum (anticorps) de souris souche B \textbf{préalablement immunisées} contre la souche A.
\end{itemize}
Résultats observés :
\begin{itemize}
    \item Lot 1 : Rejet de la greffe en \SI{12}{\pm 2} jours.
    \item Lot 2 : Rejet de la greffe en \SI{14}{\pm 3} jours.
    \item Lot 3 : Survie de la greffe \textgreater \SI{60}{jours} (pas de rejet).
\end{itemize}
\begin{enumerate}
    \item À partir des résultats des lots 1 et 3, que pouvez-vous conclure sur le rôle des anticorps (médiation humorale) dans le rejet de greffe de peau ?
    \item À partir des résultats du lot 2, identifiez la population cellulaire responsable du rejet. Quel type de réponse immunitaire est mis en évidence ?
    \item Précisez le phénotype (marqueurs CD) des lymphocytes T effecteurs qui reconnaissent les antigènes du greffon (CMH de classe I ou II ?). Décrivez brièvement leur mécanisme d'action cytotoxique.
    \item Pourquoi a-t-on utilisé des souris irradiées pour les lots 2 et 3 ? Quel contrôle supplémentaire aurait été nécessaire pour valider l'expérience (lot 4) ?
\end{enumerate}
\end{exercice}

\begin{exercice}[Type Bac : Dynamique de l'infection VIH et survenue des infections opportunistes]
Le document ci-dessous présente l'évolution simultanée du taux de lymphocytes T CD4+ (LT4) et de la charge virale (ARN VIH/mL plasma) chez un patient infecté par le VIH-1, non traité, suivi pendant 10 ans.
\begin{popfigure}
\begin{tikzpicture}[scale=0.75]
\begin{axis}[
    width=\linewidth, height=9cm,
    xlabel={Temps (années)}, ylabel={Taux de LT4 (\SI{}{\per\micro\liter})},
    xmin=0, xmax=10, ymin=0, ymax=1200,
    xtick={0,1,2,3,4,5,6,7,8,9,10},
    ytick={0,200,400,500,600,800,1000},
    axis y line*=left,
    legend style={at={(0.5,1.15)}, anchor=north, legend columns=2},
    grid=major, grid style={dashed, gray!30},
]
% CD4 curve
\addplot[popblue, very thick, smooth, mark=*] coordinates {
(0, 800) (1, 720) (2, 650) (3, 580) (4, 500) (5, 420) (6, 350) (7, 280) (8, 180) (9, 100) (10, 50)
};
\addlegendentry{LT4 CD4+ (\SI{}{\per\micro\liter})}
\end{axis}
\begin{axis}[
    xmin=0, xmax=10, ymin=1, ymax=1000000,
    ymode=log,
    axis y line*=right,
    ylabel={Charge virale (copies/mL)},
    ytick={1,100,1000,10000,100000,1000000},
    yticklabels={$10^0$,$10^2$,$10^3$,$10^4$,$10^5$,$10^6$},
    legend style={at={(0.5,1.3)}, anchor=north, legend columns=2},
]
% Viral load curve
\addplot[poporange, very thick, smooth, mark=square*] coordinates {
(0, 10000) (0.5, 500000) (1, 50000) (2, 30000) (3, 20000) (4, 30000) (5, 50000) (6, 100000) (7, 300000) (8, 500000) (9, 800000) (10, 1000000)
};
\addlegendentry{Charge virale VIH (log)}
% Phase lines
\draw[dashed, popdark] (2.5,1) -- (2.5,1000000);
\draw[dashed, popdark] (7.5,1) -- (7.5,1000000);
\node[popdark, rotate=90] at (1.2, 500000) {Phase 1};
\node[popdark, rotate=90] at (5, 500000) {Phase 2};
\node[popdark, rotate=90] at (8.7, 500000) {Phase 3};
\end{axis}
\end{tikzpicture}
\legende{Évolution de la numération des LT4 (échelle linéaire, axe gauche) et de la charge virale (échelle logarithmique, axe droit) au cours de l'infection VIH non traitée.}
\end{popfigure}
\begin{enumerate}
    \item Repérez sur le graphique les trois phases cliniques de l'infection VIH (Primaire / Aiguë, Chronique / Asymptomatique, Sida / Terminale). Indiquez les bornes temporelles approximatives pour chaque phase.
    \item Décrivez l'évolution de la charge virale et du taux de LT4 au cours de la \textbf{phase 1 (infection primaire)}. Expliquez le pic de virémie et la chute transitoire des LT4 par la réplication virale et la réponse immune innée/adaptative naissante.
    \item Pendant la \textbf{phase 2 (phase de latence clinique)}, la charge virale reste relativement basse et stable (point de consigne viral) tandis que les LT4 déclinent lentement. Expliquez ce déclin progressif malgré une charge virale contrôlée (turn-over viral, activation immune chronique, destruction des LT4 infectés et \emph{bystander apoptosis}).
    \item Caractérisez la \textbf{phase 3 (stade Sida)} par les seuils critiques de LT4 et de charge virale. Expliquez pourquoi les infections opportunistes (ex: \emph{Pneumocystis jirovecii}, \emph{Toxoplasma gondii}, \emph{Candida albicans}) surviennent spécifiquement à ce stade, en lien avec le rôle des LT4 dans l'activation des macrophages (IFN-$\gamma$) et des LT8.
\end{enumerate}
\end{exercice}

\begin{exercice}[Type Bac : Synthèse - Rôle central du LT4, immunodéficience VIH et prévention au Cameroun]
\textbf{Contexte :} Au Cameroun, la prévalence du VIH chez les adultes de 15-49 ans est estimée à \SI{2,7}{\%} (Enquête CAMPHIA 2017-2018), avec des disparités régionales et selon le genre. La prise en charge repose sur les ARV gratuits, mais la prévention reste un pilier majeur.
\begin{enumerate}
    \item \textbf{Rôle du LT4 :} En vous appuyant sur vos connaissances des réponses immunitaires spécifique (humorale et cellulaire), décrivez \textbf{trois} fonctions majeures du lymphocyte T auxiliaire (LT4 / CD4+) qui en font le « chef d'orchestre » de la réponse immune adaptative. (Indices : activation des LB, activation des macrophages, aide aux LT8).
    \item \textbf{Mécanisme de l'immunodéficience :} Expliquez de manière rigoureuse comment l'infection et la destruction progressive des LT4 par le VIH entraînent un effondrement \textbf{global} des deux bras de la réponse immune spécifique (humorale et cellulaire) et de l'immunité innée (macrophages), conduisant au Sida.
    \item \textbf{Prévention contextuelle :} Proposez \textbf{trois} mesures de prévention primaire ou secondaire adaptées au contexte camerounais (zones urbaines, zones rurales, populations clés, jeunes), en justifiant chacune par son efficacité prouvée et sa faisabilité socioculturelle. Intégrez les notions de PrEP, Tasp (Treatment as Prevention), éducation sexuelle complète, et lutte contre la stigmatisation.
\end{enumerate}
\end{exercice}

\newpage

\coursec{Corrigés des exercices}

\begin{corrige}[Exercice 1 — QCM : Origine et reconnaissance immunitaire]
\textbf{1. Réponse B — la moelle osseuse rouge.} Les lymphocytes B naissent et acquièrent leur compétence immunitaire dans la \cle{moelle osseuse rouge} (organe lymphoïde central), où s'effectue la sélection négative des LB autoréactifs. Le thymus (A) est le site de maturation des lymphocytes T ; la rate (C) et les ganglions (D) sont des organes lymphoïdes \emph{périphériques} où les lymphocytes matures sont \emph{activés}, et non formés.

\textbf{2. Réponse B — toutes les cellules nucléées de l'organisme.} Le \cle{CMH de classe I} (HLA-A, -B, -C chez l'Homme) est exprimé à la surface de toutes les cellules possédant un noyau : il présente les peptides endogènes aux LT8. Le CMH de classe II, lui, est restreint aux CPA (A). Les érythrocytes (D), anucléés, n'expriment pas de CMH — c'est d'ailleurs ce qui rend les transfusions possibles sans rejet lié au CMH.

\textbf{3. Réponse C — mise en œuvre rapide, sans mémoire.} L'immunité \cle{innée} (barrières, phagocytose, inflammation, complément, cellules NK) intervient en quelques minutes à heures, de façon identique à chaque contact : elle n'a ni spécificité clonale fine ni mémoire. Les propositions A et D décrivent l'immunité adaptative humorale ; B est contradictoire (la lenteur caractérise la réponse primaire adaptative, qui possède une mémoire).

\textbf{4. Réponse C — les lymphocytes T auxiliaires LT4 (CD4+).} La gp120 du \cle{VIH} se fixe sur le récepteur CD4 (et les corécepteurs CCR5/CXCR4). La destruction progressive des LT4, chefs d'orchestre de l'immunité adaptative, explique l'effondrement immunitaire du Sida.
\end{corrige}

\begin{corrige}[Exercice 2 — Vrai ou Faux ?]
\textbf{1. FAUX.} Les maladies auto-immunes résultent au contraire d'une \cle{rupture de la tolérance} : le système immunitaire réagit \emph{anormalement} contre des antigènes du soi (auto-antigènes), avec production d'auto-anticorps ou de LT autoréactifs. Ce n'est pas une absence de réaction, mais une réaction \emph{inappropriée et excessive} dirigée contre le soi.

\textbf{2. FAUX.} La mémoire immunitaire repose sur des \cle{lymphocytes B mémoires} \emph{et} des \cle{lymphocytes T mémoires} (LT4 et LT8 mémoires), cellules quiescentes à longue durée de vie. De plus, les LB mémoires ne sécrètent pas d'anticorps en permanence : ils se différencient en plasmocytes sécréteurs uniquement lors d'un nouveau contact avec l'antigène. (Seuls quelques plasmocytes à longue vie, dans la moelle, maintiennent un taux basal d'anticorps.)

\textbf{3. VRAI.} Dans le \cle{diabète de type 1}, une réaction auto-immune (auto-anticorps anti-îlots et LT8 cytotoxiques) détruit sélectivement les cellules $\beta$ des îlots de Langerhans, d'où l'absence de sécrétion d'insuline et l'hyperglycémie chronique, traitée par insulinothérapie substitutive.

\textbf{4. FAUX.} Le TROD (test rapide à orientation diagnostique) recherche des \cle{anticorps anti-VIH} (sérologie indirecte), parfois combinés à l'antigène p24 — jamais le génome viral. La PCR (recherche d'ARN ou d'ADN proviral) est un examen de biologie moléculaire réservé à des situations particulières (diagnostic précoce du nourrisson, fenêtre silencieuse, charge virale).
\end{corrige}

\begin{corrige}[Exercice 3 — Définitions et notions clés]
\textbf{1. Antigène :} toute molécule (protéine, polysaccharide, lipide, acide nucléique) étrangère ou perçue comme étrangère, capable d'être reconnue spécifiquement par un récepteur de lymphocyte (BCR, TCR) ou par un anticorps, et de déclencher une réponse immunitaire.

\textbf{2. Épitope :} petite région structurale précise de l'antigène (quelques acides aminés ou oses) effectivement reconnue par le site de liaison d'un anticorps ou d'un TCR. Un même antigène porte généralement plusieurs épitopes distincts.

\textbf{3. CMH de classe II :} glycoprotéines membranaires (HLA-DP, -DQ, -DR) exprimées uniquement par les cellules présentatrices d'antigènes (cellules dendritiques, macrophages, LB). Elles présentent des peptides d'origine \emph{exogène} (phagocytés puis digérés) aux lymphocytes T4 CD4+, conditionnant leur activation.

\textbf{4. Tolérance immunologique :} ensemble des mécanismes qui empêchent la réaction contre le soi. \emph{Tolérance centrale} : élimination (sélection négative) des lymphocytes autoréactifs dans le thymus (LT) et la moelle osseuse (LB). \emph{Tolérance périphérique} : anergie, apoptose ou régulation (lymphocytes T régulateurs) des clones autoréactifs ayant échappé à la sélection centrale.

\textbf{5. Infection opportuniste :} infection causée par un micro-organisme normalement inoffensif ou contrôlé chez l'individu immunocompétent, qui devient pathogène lorsque les défenses immunitaires sont effondrées (LT4 $<$ \SI{200}{\per\microliter}) : \emph{Pneumocystis jirovecii}, toxoplasmose cérébrale, candidoses, tuberculose, méningite à \emph{Cryptococcus}.
\end{corrige}

\begin{corrige}[Exercice 4 — Calcul et lecture de données : dynamique des LT4]
\textbf{1. Pourcentage de diminution :}
\[ \text{Diminution} = \frac{800 - 200}{800}\times 100 = \frac{600}{800}\times 100 = 75\,\%. \]
Le taux absolu de LT4 a diminué de \textbf{75\,\%} en 5 ans.

\textbf{2. Nouveau pourcentage de LT4 parmi les lymphocytes T :} Si les LT4 représentent 40\,\% des lymphocytes T au départ, le total des lymphocytes T vaut :
\[ \text{Total LT} = \frac{800}{0{,}40} = \SI{2000}{\per\microliter}. \]
Ce total restant constant, le nouveau pourcentage est :
\[ \frac{200}{2000}\times 100 = 10\,\%. \]
Les LT4 ne représentent plus que \textbf{10\,\%} des lymphocytes T (contre 40\,\%).

\textbf{3. Seuil du stade Sida :} on parle d'immunodéficience sévère (stade Sida, catégorie C de la CDC) lorsque le taux de LT4 descend en dessous de \textbf{\SI{200}{\per\microliter}} (ou en présence d'une infection opportuniste ou d'un cancer définissant le Sida, quel que soit le taux de LT4).
\end{corrige}

\begin{attention}[Points de vigilance — Exercice 4]
Bien distinguer \emph{taux absolu} (cellules/\si{\microliter}) et \emph{pourcentage} relatif ; ne pas confondre le seuil de \SI{200}{\per\microliter} (Sida) avec le seuil de \SI{500}{\per\microliter} (immunodépression modérée). Vérifier la cohérence : une chute de 800 à 200 est bien une division par 4, soit $-75\,\%$.
\end{attention}

\begin{corrige}[Exercice 5 — Lieux de formation et de maturation : synthèse]
\begin{center}
\begin{tabularx}{\linewidth}{|l|X|X|X|}\hline
\rowcolor{popblueL}
\textbf{Cellule} & \textbf{Site de formation} & \textbf{Site de maturation} & \textbf{Marqueur CD principal} \\ \hline
Lymphocyte B & Moelle osseuse rouge & Moelle osseuse rouge & CD19, CD20 (BCR = Ig membranaire) \\ \hline
LT CD4+ (LT4) & Moelle osseuse rouge & Thymus & CD3+, CD4+ \\ \hline
LT CD8+ (LT8) & Moelle osseuse rouge & Thymus & CD3+, CD8+ \\ \hline
Macrophage / Monocyte & Moelle osseuse rouge (lignée myéloïde) & Sang (monocyte) puis tissus (différenciation en macrophage) & CD14 \\ \hline
Cellule dendritique & Moelle osseuse rouge & Tissus périphériques (maturation après capture d'antigène, migration ganglionnaire) & CD11c \\ \hline
\end{tabularx}
\end{center}

\textbf{Remarques :} toutes les cellules immunitaires dérivent de la cellule souche hématopoïétique de la moelle osseuse rouge (hématopoïèse). Seuls les lymphocytes T migrent vers le \cle{thymus} pour y subir les sélections positive et négative ; les LB mûrissent sur place. Les cellules myéloïdes (monocytes, cellules dendritiques) n'ont pas de « maturation de compétence » au sens lymphocytaire : elles acquièrent leurs fonctions effectrices dans les tissus.
\end{corrige}

\begin{corrige}[Exercice 6 — Médiation cellulaire vs médiation humorale]
\textbf{1. Réponse prépondérante selon l'agent :}
\begin{itemize}
    \item \emph{Listeria monocytogenes} (bactérie \textbf{intracellulaire}, vivant dans le cytosol des macrophages) : réponse à \textbf{médiation cellulaire} prépondérante. Les anticorps ne peuvent atteindre une bactérie cachée dans une cellule. Effecteurs : LT4 Th1 qui activent les macrophages via l'IFN-$\gamma$ (destruction des bactéries phagocytées) et LT8 cytotoxiques qui détruisent les cellules infectées.
    \item \textbf{Virus grippal} : la réponse \textbf{humorale} est prépondérante pour neutraliser les particules virales \emph{extracellulaires} (les anticorps anti-hémagglutinine bloquent l'entrée du virus dans les cellules et opsonisent les virions). La réponse cellulaire (LT8) complète en éliminant les cellules déjà infectées.
\end{itemize}

\textbf{2. Mécanisme cytotoxique du LT8 :} le TCR du LT8 reconnaît le complexe [peptide viral + CMH de classe I] à la surface de la cellule infectée. La synapse immunologique se forme ; le LT8 exocyte ses granules : les \cle{perforines} polymérisent dans la membrane cible et y forment des pores ; les \cle{granzymes} (protéases) pénètrent par ces pores et activent les caspases, déclenchant l'\cle{apoptose} (mort cellulaire programmée) de la cellule infectée, sans lyse inflammatoire. La voie Fas/FasL complète ce mécanisme.

\textbf{3. Transfert de sérum :} le sérum immun contient des anticorps neutralisants qui agissent dans les liquides extracellulaires : ils protègent contre le virus grippal (particules virales accessibles lors de la diffusion). En revanche, \emph{Listeria} vit et se propage de cellule à cellule, à l'abri des anticorps : seule l'immunité cellulaire (macrophages activés, LT8) est efficace, d'où l'échec du transfert de sérum. C'est une preuve expérimentale classique de la complémentarité et de la spécificité d'action des deux bras de l'immunité adaptative.
\end{corrige}

\begin{corrige}[Exercice 7 — Réponse primaire et réponse secondaire]
\textbf{1. Comparaison :}
\begin{center}
\begin{tabularx}{\linewidth}{|l|X|X|}\hline
\rowcolor{popblueL}
\textbf{Critère} & \textbf{Réponse primaire (J0)} & \textbf{Réponse secondaire (J30)} \\ \hline
Latence & Longue ($\approx$ 5--7 jours) & Courte ($\approx$ 1--2 jours) \\ \hline
Cinétique de montée & Lente, pic vers J14 & Rapide, pic vers J37 ($\approx$ 7 jours après le rappel) \\ \hline
Amplitude du pic & Faible ($\approx$ 60 UA/mL) & Élevée ($\approx$ 100 UA/mL) \\ \hline
Isotypes & IgM d'abord, puis IgG de faible affinité & IgG majoritaires, de forte affinité (maturation d'affinité et commutation isotypique accomplies) \\ \hline
Persistance & Décroissance rapide & Taux élevé durable (plasmocytes longue vie) \\ \hline
\end{tabularx}
\end{center}

\textbf{2. Cellules responsables :} les \cle{lymphocytes B mémoires} (et les LT4 mémoires spécifiques), générés lors de la réponse primaire. Présents en grand nombre et déjà sélectionnés pour une forte affinité, ils se différencient rapidement en plasmocytes sécréteurs d'IgG dès le second contact avec l'antigène X.

\textbf{3. Principe du rappel vaccinal :} la primo-vaccination crée la mémoire immunologique sans provoquer la maladie ; le rappel (2\textsuperscript{e} injection) stimule les lymphocytes mémoires et induit une réponse secondaire plus rapide, plus intense et plus durable, garantissant un taux d'anticorps protecteur à long terme. Le graphique illustre exactement ce bénéfice : après le rappel à J30, le taux dépasse 100 UA/mL en moins d'une semaine et reste élevé à J60.
\end{corrige}

\begin{corrige}[Exercice 8 — Cas clinique : myasthénie auto-immune]
\textbf{1. Identification :} il s'agit de la \cle{myasthénie} (\emph{myasthenia gravis}), maladie auto-immune à \textbf{médiation humorale} : ce sont des \emph{auto-anticorps} (IgG anti-récepteurs à l'acétylcholine, détectés dans le sang) qui sont directement pathogènes, et non des lymphocytes cytotoxiques. La positivité du test à la néostigmine confirme le déficit de transmission neuromusculaire.

\textbf{2. Physiopathologie :} normalement, l'influx nerveux libère l'acétylcholine (ACh) dans la fente synaptique de la jonction neuromusculaire ; l'ACh se fixe sur les récepteurs RACh de la plaque motrice, provoquant la dépolarisation musculaire et la contraction. Les auto-anticorps anti-RACh (i) bloquent le site de fixation de l'ACh, (ii) accélèrent l'internalisation et la dégradation des récepteurs, (iii) activent le complément qui détruit la membrane post-synaptique. Résultat : diminution du nombre de récepteurs fonctionnels $\Rightarrow$ transmission neuromusculaire défaillante $\Rightarrow$ faiblesse musculaire fluctuante, majorée à l'effort et en fin de journée (épuisement du stock d'ACh), d'où ptôsis et diplopie. La néostigmine, en inhibant l'acétylcholinestérase, augmente la durée de vie de l'ACh dans la fente et améliore transitoirement la force musculaire.

\textbf{3. Rôle du thymus :} le thymus est le lieu de la \cle{sélection négative} des LT autoréactifs. Chez les myasthéniques, l'hyperplasie thymique (ou le thymome) témoigne d'une défaillance de cette tolérance centrale : les cellules épithéliales thymiques expriment des antigènes de type RACh, et une sélection négative défaillante laisse échapper des clones de LT4 autoréactifs anti-RACh, qui activent ensuite les LB producteurs d'auto-anticorps. Un \cle{mimétisme moléculaire} (ressemblance entre un épitope microbien et le RACh) peut initier la rupture de tolérance.

\textbf{4. Pistes thérapeutiques ciblant l'immunité :}
\begin{itemize}
    \item \textbf{Corticothérapie et immunosuppresseurs} (azathioprine, mycophénolate) : réduire la production d'auto-anticorps.
    \item \textbf{Thymectomie} : ablation chirurgicale du thymus hyperplasique ou du thymome, supprimant le site d'auto-sensibilisation.
    \item (Autres options acceptées : échanges plasmatiques ou immunoglobulines intraveineuses en crise aiguë, pour éliminer/neutraliser les auto-anticorps circulants.)
\end{itemize}
\end{corrige}

\begin{corrige}[Exercice 9 — Type Bac : expérience de Mitchison (1954)]
\textbf{1. Rôle des anticorps (lots 1 et 3) :} le lot 3, qui reçoit le sérum (anticorps) de souris immunisées, ne rejette pas la greffe (survie $>$ \SI{60}{jours}), alors que le lot témoin la rejette en \SI{12}{\pm 2}{jours}. Conclusion : les \textbf{anticorps seuls sont incapables de provoquer le rejet} d'une greffe de peau ; la médiation humorale n'est pas l'effecteur du rejet aigu.

\textbf{2. Lot 2 :} le transfert de \cle{lymphocytes T} purifiés de souris immunisées restaure le rejet (\SI{14}{\pm 3}{jours}, cinétique comparable au lot témoin). Conclusion : les lymphocytes T sont les \textbf{effecteurs du rejet de greffe} ; il s'agit d'une réponse immunitaire à \textbf{médiation cellulaire}. (La réponse est rapide car les cellules transférées sont déjà immunisées : réponse secondaire.)

\textbf{3. Phénotype et mécanisme :} les effecteurs sont des \textbf{LT8 cytotoxiques CD3+CD8+}, qui reconnaissent les peptides du donneur présentés par le \cle{CMH de classe I} (exprimé par toutes les cellules nucléées du greffon) ; des LT4 (reconnaissant le CMH de classe II des CPA du greffon) assurent la coopération (IL-2). Mécanisme : formation de la synapse immunologique, exocytose de \cle{perforines} (pores membranaires) et de \cle{granzymes} (activation des caspases) $\Rightarrow$ apoptose des cellules du greffon, d'où sa destruction et son rejet.

\textbf{4. Intérêt des souris irradiées et contrôle manquant :} l'irradiation détruit le système immunitaire propre du receveur : tout rejet observé est donc attribuable \emph{uniquement} aux cellules ou au sérum transférés (condition nécessaire pour conclure). \textbf{Lot 4 (contrôle) :} souris irradiées + greffe de souche A \emph{sans} transfert (ou avec transfert de lymphocytes de souris \emph{non} immunisées) : la greffe doit survivre, ce qui validerait que le rejet du lot 2 est bien dû aux lymphocytes T spécifiquement immunisés et non à un artefact.
\end{corrige}

\begin{attention}[Barème indicatif et points de vigilance — Exercice 9]
Barème proposé : Q1 : 1 pt ; Q2 : 1,5 pt ; Q3 : 1,5 pt ; Q4 : 1 pt. Attendus : citer explicitement la \emph{comparaison des lots} (démarche comparative), la condition « souris irradiées = système immunitaire aboli », le couple peptide--CMH I / TCR--CD8, et le rôle des perforines/granzymes. Erreur fréquente : attribuer le rejet aux anticorps par simple association « immunité = anticorps ».
\end{attention}

\begin{corrige}[Exercice 10 — Type Bac : dynamique de l'infection VIH]
\textbf{1. Les trois phases :}
\begin{itemize}
    \item \textbf{Phase 1 — infection primaire (aiguë) :} environ de 0 à 2,5 ans sur le graphique (pic de virémie précoce vers 6 mois, puis contrôle partiel).
    \item \textbf{Phase 2 — phase chronique (latence clinique) :} environ de 2,5 à 7,5 ans (charge virale stabilisée au « point de consigne », déclin lent des LT4).
    \item \textbf{Phase 3 — stade Sida :} au-delà de $\approx$ 7,5 ans (LT4 $<$ \SI{200}{\per\microliter}, charge virale $> 10^5$ copies/mL).
\end{itemize}

\textbf{2. Phase 1 :} la charge virale explose (pic à $\approx \num{5e5}$ copies/mL vers 6 mois) car le virus se réplique massivement dans les LT4 en l'absence de réponse adaptative ; cette réplication lyse les cellules infectées, d'où une chute transitoire des LT4 (de 800 à $\approx$ 720/\si{\microliter}). L'installation de la réponse immune (LT8 cytotoxiques puis anticorps neutralisants, séroconversion) fait chuter la virémie vers $\num{2e4}$--$\num{5e4}$ copies/mL et permet une remontée partielle des LT4 — mais sans éradication (réservoirs viraux).

\textbf{3. Phase 2 :} la charge virale se stabilise (point de consigne, $\approx \num{2e4}$ à $\num{1e5}$ copies/mL) : équilibre dynamique entre production virale (turn-over de l'ordre de $\num{1e9}$--$\num{1e10}$ virions/jour) et clairance par l'immunité. Malgré ce contrôle relatif, les LT4 déclinent lentement ($\approx$ 50 à 80 cellules/\si{\microliter}/an) à cause de : (i) la destruction directe des LT4 infectés par effet cytopathogène du virus et par les LT8 ; (ii) l'\cle{activation immune chronique} qui épuise et fait apoptoser des LT4 non infectés (\emph{bystander apoptosis}) ; (iii) la destruction progressive de l'architecture des organes lymphoïdes, limitant la régénération.

\textbf{4. Phase 3 (Sida) :} seuils critiques : LT4 $<$ \SI{200}{\per\microliter} (ici atteint vers l'année 8) et charge virale en forte recrudescence ($> \num{1e5}$ copies/mL, jusqu'à $\num{1e6}$). Les LT4 étant indispensables à l'activation des macrophages (via l'\cle{IFN-$\gamma$} des LT4 Th1) et à l'aide aux LT8 (IL-2) et aux LB, leur effondrement paralyse simultanément l'immunité cellulaire et humorale. Des micro-organismes normalement contrôlés (\emph{Pneumocystis jirovecii}, \emph{Toxoplasma gondii}, \emph{Candida albicans}, \emph{Cryptococcus}, bacilles tuberculeux) prolifèrent alors : ce sont les \cle{infections opportunistes}, définissant le stade Sida.
\end{corrige}

\begin{attention}[Barème indicatif et points de vigilance — Exercice 10]
Barème proposé : Q1 : 1 pt ; Q2 : 1,5 pt ; Q3 : 1,5 pt ; Q4 : 2 pts. Attendus : lecture chiffrée du graphique (valeurs et dates citées), mention obligatoire du rôle de l'IFN-$\gamma$ et de l'activation des macrophages pour expliquer les infections opportunistes, et distinction nette entre charge virale (axe droit, échelle log) et LT4 (axe gauche, échelle linéaire) — erreur de lecture très fréquente.
\end{attention}

\begin{corrige}[Exercice 11 — Type Bac : synthèse VIH et prévention au Cameroun]
\textbf{1. Trois fonctions majeures du LT4 :}
\begin{itemize}
    \item \textbf{Activation des lymphocytes B :} le LT4 spécifique reconnaît le peptide présenté par le CMH II du LB et lui délivre des signaux de coopération (CD40L, IL-4, IL-21) indispensables à sa différenciation en plasmocyte, à la commutation isotypique (IgM $\to$ IgG) et à la maturation d'affinité.
    \item \textbf{Activation des macrophages :} les LT4 Th1 sécrètent l'\cle{IFN-$\gamma$}, qui potentialise les pouvoirs phagocytaire et bactéricide des macrophages (destruction des germes intracellulaires).
    \item \textbf{Aide aux LT8 cytotoxiques :} par la sécrétion d'\cle{IL-2}, les LT4 permettent la prolifération clonale et la différenciation des LT8 en cellules cytotoxiques effectrices et en cellules mémoires.
\end{itemize}
Le LT4 orchestre donc à la fois l'immunité humorale, l'immunité cellulaire et l'immunité innée : c'est le « chef d'orchestre » de la réponse adaptative.

\textbf{2. Mécanisme de l'immunodéficience :} le VIH infecte préférentiellement les LT4 (fixation de la gp120 sur CD4 + corécepteurs CCR5/CXCR4), s'y réplique et les détruit (effet cytopathogène, lyse par les LT8, apoptose des cellules voisines). La chute des LT4 entraîne : (i) l'absence d'aide aux LB $\Rightarrow$ production d'anticorps effondrée (bras \textbf{humoral} défaillant) ; (ii) l'absence d'IL-2 $\Rightarrow$ défaut de prolifération des LT8 (bras \textbf{cellulaire} défaillant) ; (iii) l'absence d'IFN-$\gamma$ $\Rightarrow$ macrophages non activés (immunité \textbf{innée} affaiblie). L'ensemble des défenses s'effondre : c'est l'immunodéficience acquise, qui ouvre la porte aux infections opportunistes et aux cancers (Sida).

\textbf{3. Trois mesures de prévention adaptées au contexte camerounais :}
\begin{itemize}
    \item \textbf{Traitement comme prévention (TasP) et ARV gratuits :} une personne sous ARV avec charge virale indétectable ne transmet pas le virus (I = I, « indétectable = intransmissible »). Efficacité prouvée ; faisable grâce à la gratuité des ARV et au réseau des centres de traitement agréés — à condition de lutter contre la stigmatisation qui freine le dépistage et l'observance.
    \item \textbf{Dépistage généralisé et PrEP pour les populations clés :} dépistage rapide (TROD) dans les centres de santé, les campagnes communautaires et l'autodépistage chez les jeunes ; prophylaxie pré-exposition (PrEP) pour les personnes très exposées. Efficacité démontrée ; faisabilité renforcée par l'implication des associations communautaires et des leaders religieux et traditionnels.
    \item \textbf{Éducation sexuelle complète et promotion du préservatif :} information adaptée à l'âge dans les écoles, universités et médias communautaires (radios locales), distribution de préservatifs, compétences psychosociales (négociation, estime de soi, refus des rapports forcés). Mesure peu coûteuse, culturellement négociable via le dialogue avec les familles et les autorités coutumières, et indispensable pour inverser l'épidémie chez les jeunes.
\end{itemize}
La lutte contre la \cle{stigmatisation} et la discrimination est transversale : vivre avec le VIH, sous traitement, est compatible avec une vie scolaire, professionnelle et familiale normale ; la solidarité envers les personnes affectées est un savoir-être essentiel.
\end{corrige}

\begin{attention}[Barème indicatif et points de vigilance — Exercice 11]
Barème proposé : Q1 : 2 pts (trois fonctions justifiées, avec cytokines citées) ; Q2 : 2 pts (chaîne causale complète reliant destruction des LT4 aux trois défaillances) ; Q3 : 2 pts (mesures justifiées par efficacité + faisabilité socioculturelle). Attendus : vocabulaire exact (CD4, gp120, IFN-$\gamma$, IL-2), exemples locaux crédibles, ton respectueux et non stigmatisant. Piège : ne pas écrire que le VIH « détruit tous les lymphocytes » — sa cible principale est le LT4.
\end{attention}

\newpage

\coursec{Fiche bilan}

\courssub{Carte mentale de la leçon}

\begin{popfigure}
\begin{tikzpicture}[mindmap, grow cyclic, every node/.style={concept}, text width=2.6cm, align=center,
root concept/.append style={concept color=popblue, font=\bfseries, text=white, minimum size=3.2cm},
level 1 concept/.append style={level distance=4.6cm, sibling angle=51, font=\small},
level 2 concept/.append style={level distance=3.1cm, font=\scriptsize, text width=2.2cm}]
\node[root concept]{Mécanismes de l'immunité et dysfonctionnements}
  child[concept color=popgreen] { node {Origine des cellules immunitaires}
    child { node {Moelle osseuse rouge : PSC}}
    child { node {Maturation : thymus (LT), moelle (LB)}}
  }
  child[concept color=poporange] { node {Reconnaissance soi / non-soi}
    child { node {CMH (marqueurs du soi)}}
    child { node {Antigènes, anticorps, TCR}}
  }
  child[concept color=poppurple] { node {Immunité non spécifique}
    child { node {Barrières, phagocytose, inflammation}}
  }
  child[concept color=popblue!60] { node {Immunité spécifique}
    child { node {Humorale : LB $\to$ plasmocytes $\to$ Ac}}
    child { node {Cellulaire : LT8 $\to$ LTc $\to$ lyse}}
  }
  child[concept color=popgold] { node {Mémoire immunitaire}
    child { node {Cellules mémoires}}
    child { node {Vaccination, rappels}}
  }
  child[concept color=poppink] { node {Maladies auto-immunes}
    child { node {Vitiligo, diabète juvénile, myasthénie}}
  }
  child[concept color=popdark] { node {VIH / Sida}
    child { node {Rétrovirus, LT4 détruits}}
    child { node {Prévention, ARV, solidarité}}
  };
\end{tikzpicture}
\legende{Carte mentale de synthèse : les sept branches de la leçon, de l'origine des cellules immunitaires aux dysfonctionnements du système immunitaire.}
\end{popfigure}

\courssub{L'essentiel à retenir}

\begin{aretenir}[Définitions clés]
\begin{itemize}
\item \cle{Antigène} : substance étrangère (non-soi) ou modifiée, reconnue par le système immunitaire et capable de déclencher une réponse immunitaire.
\item \cle{Anticorps} (immunoglobuline) : protéine en Y sécrétée par les plasmocytes, qui se fixe spécifiquement sur un antigène pour le neutraliser.
\item \cle{CMH} (complexe majeur d'histocompatibilité, HLA chez l'Homme) : ensemble de glycoprotéines membranaires, marqueurs du \cle{soi}, propres à chaque individu (sauf vrais jumeaux).
\item \cle{Immunocompétence} : capacité d'une cellule immunitaire à reconnaître un antigène déterminé et à réagir contre lui.
\item \cle{Maladie auto-immune} : maladie due à une réaction du système immunitaire contre les propres constituants de l'organisme (rupture de la tolérance du soi).
\item \cle{Sida} : stade ultime de l'infection par le VIH, caractérisé par une immunodépression profonde (chute des LT4) et l'apparition de maladies opportunistes.
\end{itemize}
\end{aretenir}

\begin{aretenir}[Origine et maturation des cellules immunitaires]
\begin{center}
\begin{tabularx}{\linewidth}{|l|X|X|}
\hline
\rowcolor{popblueL} \textbf{Cellule} & \textbf{Origine} & \textbf{Lieu de maturation} \\
\hline
Toutes les cellules immunitaires & Moelle osseuse rouge (cellules souches hématopoïétiques) & --- \\
\hline
Lymphocytes B (LB) & Moelle osseuse rouge & Moelle osseuse rouge \\
\hline
Lymphocytes T (LT4, LT8) & Moelle osseuse rouge & Thymus \\
\hline
Phagocytes (polynucléaires, monocytes/macrophages) & Moelle osseuse rouge & Moelle puis sang et tissus \\
\hline
\end{tabularx}
\end{center}
\end{aretenir}

\begin{propriete}[Les deux voies de la réponse immunitaire spécifique]
\begin{center}
\begin{tabularx}{\linewidth}{|l|X|X|}
\hline
\rowcolor{popblueL} \textbf{Critère} & \textbf{Médiation humorale} & \textbf{Médiation cellulaire} \\
\hline
Effecteur & LB $\to$ plasmocytes & LT8 $\to$ lymphocytes T cytotoxiques (LTc) \\
\hline
Cible & Antigènes libres (microbes, toxines en milieu extracellulaire) & Cellules infectées ou cancéreuses (antigènes intracellulaires présentés par CMH I) \\
\hline
Arme & Anticorps (neutralisation, agglutination, opsonisation) & Contact direct : perforines et enzymes $\to$ lyse de la cellule cible \\
\hline
Coopération & LT4 auxiliaires (interleukines) & LT4 auxiliaires (interleukines) \\
\hline
\end{tabularx}
\end{center}
\textbf{Sélection clonale} : l'antigène sélectionne le clone lymphocytaire qui lui est spécifique $\to$ multiplication (amplification clonale) $\to$ différenciation en cellules effectrices et en \cle{cellules mémoires}.
\end{propriete}

\begin{aretenir}[Mémoire immunitaire, dysfonctionnements et VIH/Sida]
\begin{itemize}
\item \textbf{Réponse primaire} : lente (latence d'environ 7 à 10 jours), faible. \textbf{Réponse secondaire} : rapide, intense, durable (cellules mémoires) $\Rightarrow$ principe de la \cle{vaccination}.
\item \textbf{Maladies auto-immunes} : vitiligo (destruction des mélanocytes), diabète juvénile de type 1 (destruction des cellules $\beta$ des îlots de Langerhans $\to$ absence d'insuline), myasthénie (auto-anticorps anti-récepteurs de l'acétylcholine $\to$ faiblesse musculaire).
\item \textbf{VIH} : rétrovirus (génome ARN + transcriptase inverse) qui infecte et détruit les \cle{LT4} (récepteur CD4) $\to$ effondrement des réponses immunitaires $\to$ maladies opportunistes (tuberculose, méningites à cryptocoques, pneumocystose).
\item \textbf{Transmission} : sang, rapports sexuels non protégés, mère-enfant. \textbf{Pas de transmission} par la salive, les larmes, les poignées de main, les moustiques.
\item \textbf{Prévention} : préservatif, dépistage, fidélité réciproque, matériel d'injection à usage unique ; \textbf{traitement} : antirétroviraux (ARV) qui bloquent la multiplication virale sans guérir.
\end{itemize}
\end{aretenir}

\begin{methode}[Méthodes d'analyse et d'interprétation]
\begin{enumerate}
\item \textbf{Expérience de greffe} : greffe rejetée entre individus non apparentés, acceptée entre vrais jumeaux $\Rightarrow$ le CMH est le marqueur du soi.
\item \textbf{Test de Winn} : mélange cellules tumorales + lymphocytes T d'un animal immunisé $\to$ pas de tumeur chez l'animal receveur $\Rightarrow$ preuve de la médiation cellulaire.
\item \textbf{Courbe d'anticorps} : comparer les pentes et les maxima des réponses primaire et secondaire pour conclure à la mémoire immunitaire.
\item \textbf{Analyse d'un western blot ou d'un test ELISA} : présence d'anticorps anti-VIH = séropositivité (infection), à confirmer.
\end{enumerate}
\end{methode}

\begin{aretenir}[Checklist avant le Bac]
\begin{itemize}
\item[$\square$] Je sais situer la moelle osseuse rouge (formation) et le thymus (maturation des LT).
\item[$\square$] Je définis antigène, anticorps, CMH et immunocompétence sans hésiter.
\item[$\square$] J'explique pourquoi une greffe entre vrais jumeaux est acceptée.
\item[$\square$] Je décris les étapes de la phagocytose (réponse non spécifique).
\item[$\square$] Je distingue médiation humorale (LB, anticorps) et médiation cellulaire (LTc, lyse par contact).
\item[$\square$] Je sais que les LT4 auxiliaires coopèrent dans les deux voies.
\item[$\square$] J'interprète une courbe d'anticorps (réponses primaire et secondaire) et je la relie à la vaccination.
\item[$\square$] Je cite trois maladies auto-immunes et leur mécanisme (vitiligo, diabète juvénile, myasthénie).
\item[$\square$] J'explique comment le VIH détruit les LT4 et provoque l'immunodépression.
\item[$\square$] Je connais les trois modes de transmission du VIH et je réfute les idées reçues.
\item[$\square$] Je sais proposer des messages de prévention et de solidarité envers les PVVIH.
\end{itemize}
\end{aretenir}

\findecours

\end{document}
